% Generated by blueprint/scripts/check.py --write from nodes.json and sections.json.
% Edit those files, not this one.

\section{Introduction}
\label{section-1}
This blueprint follows Arthur Forey and Yimu Yin, \emph{Bounded integral and motivic Milnor fiber}, in its final form \href{https://github.com/daviddhilbert/diamond-integral/blob/main/paper/bounded-integral-motivic-milnor-fiber.pdf}{in this repository}, and links its numbered results to the Lean 4 formalization. Its sections and numbers are those of the paper: Lemma 7.5 here is Lemma 7.5 there. The paper's introduction has no results of its own; the results it announces are proved in Sections 2--8.

Each node is one numbered item of the paper: every definition, lemma, proposition, theorem and corollary, and every notation, terminology, convention or remark that a later statement uses or that makes a mathematical claim. Six items are left out: Convention 2.9 (a syntax convention), Notation 3.6 (an explanation of notation), Examples 6.3 and 6.25, Terminology 6.18 (a procedure used inside proofs) and Remark 8.12 (commentary). Equation (8.6), which is formalized separately, is the node Formula 8.6. The statements are those of the paper, shortened where noted; the proof outlines say which earlier results each proof combines. Amber nodes are the explicit hypotheses that the Lean theorems take as arguments in place of a proof: a consequence of a theorem of Reichstein and Youssin (Section 3), and the \'etale chart and \'etale coordinates of a smooth variety at a point (Section 8). Purple nodes are the three results of Section 5, which are recalled from the Fubini paper and proved in the Fubini-theorems package.

The main results are the Hrushovski--Kazhdan integral (Theorem~\ref{main-prop-HK}), the relative annular ideal and its compatibility (Notation~\ref{relative-annular-ideal}, Lemma~\ref{relative-annular-compatibility}), the bounded integral as a ring isomorphism onto the quotient by that ideal (Theorem~\ref{main-prop}, with Proposition~\ref{kernel-L} and Corollary~\ref{L-sur-c}) and its squares (Corollary~\ref{db-to-novol}), and the motivic Milnor fiber (Theorem~\ref{direct-mil}, with Lemma~\ref{lat-zeta}, Formula~\ref{z-coeff-change} and Remarks~\ref{rem-comp-milnor-fib} and~\ref{rem-comp-milnor-fib-bis}).

A node is \emph{formalized} when a Lean declaration states the paper's statement, possibly with the recorded explicit hypotheses; it is then marked as done, and its proof is marked as done. A \emph{partial} node has a Lean statement of a recorded narrower scope; its statement is marked, its proof is not. A node without Lean declarations is \emph{not formalized}. Each partial or unformalized node says what differs and who decided it; the decisions are collected in the repository's \href{https://github.com/daviddhilbert/diamond-integral/blob/main/FINAL-CHECK.md}{final check}.

The Lean statements share conventions that the paper states once. The structure is a saturated model of $\ACVF$ of equicharacteristic zero in a one-sorted encoding of the $\VF$, $\RV$ and $\Gamma$ sorts; the parameter base is generated by $\VF$ points and a subgroup of $\Gamma$, with nontrivial $\Gamma$ throughout; saturation, smallness and a valuation compatible with the language are hypotheses of most statements. For the main theorems they are supplied once, in the Puiseux ambient of Section 5, so that Theorems~\ref{main-prop} and~\ref{direct-mil} and Remarks~\ref{rem-comp-milnor-fib} and~\ref{rem-comp-milnor-fib-bis} also hold with no hypothesis (\texttt{Closed.theorem\_8\_11} and its companions). The explicit hypotheses are the Reichstein--Youssin consequence, used for Remark 3.15 and for the injectivity of the coefficient maps of Section 3, and the \'etale chart and coordinates, forms of Hensel's lemma used for smooth varieties in Section 8 and proved for affine space.

Source links point to a fixed revision of this repository, or to the pinned revision of the provider package that holds the proof: most modules of this repository are compatibility imports of hilbert-mathlib, Fubini-theorems and the upstream paper packages. A declaration is linked only if a section receipt of this repository records its type and axioms (only \texttt{propext}, \texttt{Classical.choice} and \texttt{Quot.sound} occur), or, in Section 5, if the public contract of the pinned Fubini-theorems package does.

\section{Technicalities}
\label{section-2}
Conventions on the language and the parameter base, discs and their names, dimension, continuity and the special bijections and correspondences used later. Most of these results reuse the special-transformations and integration-in-ACVF packages through checked adapters.

\setcounter{theorem}{0}
\begin{definition}[Coordinate projections]
\label{indexing}
For $n \in \mathbb{N}$ let $[n] = \{1, \ldots, n\}$. For $E \sub [n]$, $\operatorname{pr}_E(A)$, or $A_E$, is the projection of $A$ to the coordinates in $E$; a singleton $\{i\}$ is written $i$, $\tilde E = [n] \mi E$, and similarly $\operatorname{pr}_{\leq i}$, $A_{[i,j]}$, $A_{\VF}$, $A_{\RV}$ (projections to the $\VF$- and $\RV$-sort coordinates). A point $a \in A$ is a finite tuple, with subtuples $a_i$, $a_{\tilde i}$, $a_{\leq i}$, $a_{\VF}$, etc.; a singleton $\{t\}$ may be written $t$, as in $t \times A$. For $a \in A_{\tilde E}$ the fiber $\{b : (b,a) \in A\} \sub A_E$ is denoted $A_a$ and identified with $A_a \times \{a\}$. Functions are identified with their graphs, so for $f : A \fun B$ and $b \in B$ the pullback $f^{-1}(b)$ is also written $A_b$.

\emph{Formalization.} Notational convention with no separate declaration: the Lean statements use Mathlib tuples, coordinate maps and fibers directly (section 2/COVERAGE.md, items 2.1--2.7).
\end{definition}

\setcounter{theorem}{1}
\begin{definition}[The element $\infty$ of $\RV$]
\label{s2-rv-infinity}
\lean{DiamondIntegral.Section2.RVInfinity, DiamondIntegral.Section2.RV, DiamondIntegral.Section2.rvZero, DiamondIntegral.Section2.rv_eq_zero_iff}
\leanok
The $\RV$-sort has a special element $\infty = \rv(0)$. One writes $\RV = \rv(\VF^\times)$ for the $\RV$-sort without $\infty$, and $\RV_\infty = \rv(\VF)$ for the whole sort.
\end{definition}

\setcounter{theorem}{2}
\begin{definition}[Sets and subsets]
\label{nota-sub}
\lean{DiamondIntegral.Section2.Mixed, DiamondIntegral.Section2.MixedDefinable, DiamondIntegral.Section2.ParameterBase}
\leanok
\uses{s2-rv-infinity}
Throughout, $\mathbb{U}$ is a sufficiently saturated model of $\ACVF(0,0)$ and $\mathbb{S}$ is a substructure of $\mathbb{U}$ with $\Gamma(\mathbb{S})$ nontrivial; definable means $\mathbb{S}$-definable. A definable set in $\VF$ is a definable subset of $\VF^n$ for some $n$; similarly for other sorts or structures such as $\RV_\infty$, $\K$, $\mathcal{M}$, or a substructure $\mathbb{M}$ of $\mathbb{U}$. A definable set without further qualification is a definable subset of $\VF^n \times \RV_\infty^m$ for some $n, m \in \mathbb{N}$.

\emph{Formalization.} Definable sets are sets of mixed tuples definable over a parameter base of real ($\VF$, $\RV$) parameters and $\Gamma$-names. Saturation, smallness of the base and nontriviality of $\Gamma(\mathbb{S})$ are explicit hypotheses of the individual Lean theorems (section 2/COVERAGE.md).
\end{definition}

\setcounter{theorem}{3}
\begin{definition}[Relations as functions]
\label{s2-relations-as-functions}
\lean{DiamondIntegral.Section2.vfCorrespondenceGraph}
\leanok
\uses{indexing}
A relation $C \sub A \times B$ may be regarded as a function from $A$ into the powerset $\mathcal{P}(B)$; for instance a definable set $A \sub \VF^n \times \RV^m$ is a definable function from $A_{\VF}$ into $\mathcal{P}(\RV^m)$. A definable function $f : A \fun \mathcal{P}(B)$ is identified with the definable set $\bigcup_{a \in A} a \times f(a)$. Such an $f$ is finitary if $f(a)$ is finite for every $a \in A$.

\emph{Formalization.} A set-valued function is a map into sets whose graph (\texttt{vfCorrespondenceGraph}) is required to be definable; finitary is the hypothesis that every value is finite, not a separate predicate.
\end{definition}

\setcounter{theorem}{4}
\begin{definition}[Polydiscs and $\RV$-hulls]
\label{defn-disc}
\lean{DiamondIntegral.Section2.IsOpenDisc, DiamondIntegral.Section2.IsClosedDisc, DiamondIntegral.Section2.rvHull, DiamondIntegral.Section2.IsRVPullback}
\leanok
\uses{s2-rv-infinity}
A set $\mathfrak{p} = \prod_{i \leq n} \mathfrak{b}_i \times t \sub \VF^n \times \RV_\infty^m$ is an (open, closed, $\RV$-) polydisc if each $\mathfrak{b}_i$ is an (open, closed, $\RV$-) disc; $\rad(\mathfrak{p}) = (\rad(\mathfrak{b}_1), \ldots, \rad(\mathfrak{b}_n))$. The open polydisc centered at $(a,t)$ of radius $\gamma \in \Gamma^n$ is $\mathfrak{o}((a,t),\gamma)$; the closed one of radius $\gamma \in \Gamma_\infty^n$ is $\mathfrak{c}((a,t),\gamma)$, with $\mathfrak{c}(a,\infty) = \{a\}$; open discs have finite radius. Write $\mathcal{M}_\gamma = \mathfrak{o}(0,\gamma)$, $\mathcal{O}_\gamma = \mathfrak{c}(0,\gamma)$, $\mathcal{M} = \mathcal{M}_0$, $\mathcal{O} = \mathcal{O}_0$, $\RV^{\circ\circ}_\infty = \rv(\mathcal{M})$, $\RV^{\circ\circ} = \RV^{\circ\circ}_\infty \mi \infty$, $\RV^{\circ\circ}_\gamma = \rv(\mathcal{M}_\gamma \mi 0)$. An $\RV$-disc is a fiber $\rv^{-1}(t)$, $t \in \RV_\infty$, including $\rv^{-1}(\infty) = \{0\}$; an $\RV$-polydisc keeps the $\RV_\infty$-coordinates fixed. The $\RV$-hull $\operatorname{RVH}(A)$ is the union of the $\RV$-polydiscs meeting $A$; $A$ is an $\RV$-pullback if $A = \operatorname{RVH}(A)$.
\end{definition}

\setcounter{theorem}{5}
\begin{definition}[$\VF$-fibers and fiberwise properties]
\label{rvfiber}
\lean{DiamondIntegral.Section2.GraphFibreContinuous, DiamondIntegral.Section2.labelledComp, DiamondIntegral.Section2.GraphFibreContinuous.labelledComp}
\leanok
\uses{indexing}
A $\VF$-fiber of a set $A$ is a set $A_t$ with $t \in A_{\RV}$; a $\VF$-fiber of a function $f : A \fun B$ is a set $f_t$ with $t \in f_{\RV}$ ($f$ standing for its graph), which is a function. $A$ is open if every $\VF$-fiber of $A$ is, $f$ is continuous if every $\VF$-fiber of $f$ is, and so on. Composition is fiberwise and retains the intermediate $\RV$-coordinates: for fixed input, intermediate and output $\RV$-tuples, the corresponding $\VF$-fiber maps compose continuously on their natural domain. Thus if $f$, $g$ are continuous with $\operatorname{ran}(f) \sub \operatorname{dom}(g)$, the composite with these labels retained is continuous.

\emph{Formalization.} Fiberwise continuity and continuity of the labelled composite are formalized for total maps $X \times I \fun Y \times J$; the composition sentence for maps on their natural domains is not stated. Formalization-session choice, no author decision on record (FINAL-CHECK.md, section 3).
\end{definition}

\setcounter{theorem}{6}
\begin{definition}[The value group as a sort]
\label{gamwhat}
\lean{DiamondIntegral.Section2.Gamma, DiamondIntegral.Section2.GammaInfinity, DiamondIntegral.Section2.rvRawFibre, DiamondIntegral.Section2.vfValue_eq_iff}
\leanok
The value group $\Gamma$ is treated semantically as a definable sort, although syntactically references to it can be eliminated; $\gamma \in \Gamma$ is definable if and only if $\vv^{-1}(\gamma)$ is. For $\gamma \in \Gamma$ put $\gamma^\sharp = \vrv^{-1}(\gamma) \sub \RV$; for $I \sub \Gamma$ put $I^\sharp = \bigcup\{\gamma^\sharp : \gamma \in I\}$; for $U \sub \RV$ put $U^\sharp = \bigcup\{\rv^{-1}(t) : t \in U\}$. In particular $\gamma^{\sharp\sharp} = \{x \in \VF : \vv(x) = \gamma\}$.

\emph{Formalization.} $\Gamma$ is the quotient of the $\RV$-sort by equality of values; the $\sharp$-sets are written inline as preimages (\texttt{rvRawFibre} for a single $\RV$-disc), with no separate declaration for $I^\sharp$.
\end{definition}

\setcounter{theorem}{7}
\begin{definition}[The $\Gamma$-Jacobian]
\label{jcbG-op}
\lean{DiamondIntegral.Section2.gammaJacobian, DiamondIntegral.Section2.gammaJacobian_eq_companion}
\leanok
\uses{gamwhat}
Let $U \sub \RV^n \times \Gamma^m$, $V \sub \RV^{n'} \times \Gamma^{m'}$ and $C \sub U \times V$. The $\Gamma$-Jacobian of $C$ at $((u,\alpha),(v,\beta)) \in C$ is
\[ \jcb_{\Gamma}((u,\alpha),(v,\beta)) = -\Sigma(\vrv(u),\alpha) + \Sigma(\vrv(v),\beta), \]
where $\Sigma(\gamma_1, \ldots, \gamma_k) = \gamma_1 + \cdots + \gamma_k$. If $C$ is the graph of a function, one writes $\jcb_{\Gamma} C(u,\alpha)$ for $\jcb_{\Gamma}((u,\alpha), C(u,\alpha))$.
\end{definition}

\setcounter{theorem}{9}
\begin{definition}[The definable sort $\mathrm{DC}$ of discs]
\label{disc-exp}
\lean{DiamondIntegral.Section2.DC, DiamondIntegral.Section2.openDiscCode, DiamondIntegral.Section2.closedDiscCode, DiamondIntegral.Section2.openDiscCode_ne_closedDiscCode}
\leanok
\uses{defn-disc, gamwhat}
$\mathbb{U}^\bullet$ is the expansion of $\mathbb{U}$ by one sort $\mathrm{DC}$ whose elements are all open and closed discs; a point of $\VF$ is a closed disc of radius $\infty$. It has two uniform maps $\VF \times \Gamma \fun \mathrm{DC}$, $(a,\gamma) \efun \mathfrak{o}(a,\gamma)$ and $(a,\gamma) \efun \mathfrak{c}(a,\gamma)$, the second extended to $\VF \times \Gamma_\infty$ by $(a,\infty) \efun \{a\}$; the radius is an argument. The element of $\mathrm{DC}$ corresponding to a disc $\mathfrak{a}$ (its name) is $\ulcorner \mathfrak{a} \urcorner$, and $D \sub \mathrm{DC}$ is identified with $\{\mathfrak{a} : \ulcorner \mathfrak{a} \urcorner \in D\}$. Parameters in $\mathrm{DC}$ are used only when indicated, by names $\ulcorner \mathfrak{a} \urcorner$ or the phrase ``in $\mathbb{U}^\bullet$''. A disc centered at $0$ is identified with its radius and its open/closed designation; at a finite radius the open and closed discs are distinct.
\end{definition}

\setcounter{theorem}{10}
\begin{lemma}[Definable closure and henselianity]
\label{cut-to-hensel-substru}
\lean{DiamondIntegral.Section2.lemma_2_11}
\leanok
Suppose that $\mathbb{S}$ is $\VF$-generated (generated by elements of $\VF$) and $\Gamma(\mathbb{S})$ is nontrivial. Then $\operatorname{dcl}\mathbb{S} = \mathbb{S}$ if and only if the underlying valued field $(\VF(\mathbb{S}), \mathcal{O}(\mathbb{S}))$ of $\mathbb{S}$ is henselian.
\end{lemma}

\begin{proof}
\leanok
$\VF(\operatorname{acl}\mathbb{S})$ is the algebraic closure of $\VF(\mathbb{S})$, and the automorphisms of $\operatorname{acl}\mathbb{S}$ over $\mathbb{S}$ correspond to the valued-field automorphisms of this closure over $(\VF(\mathbb{S}), \mathcal{O}(\mathbb{S}))$, whose fixed field is $\VF(\mathbb{S})$ exactly when it is henselian. Hence $\VF(\operatorname{dcl}\mathbb{S}) = \VF(\mathbb{S})$ if and only if the field is henselian, and quantifier elimination gives $\RV(\operatorname{dcl}\mathbb{S}) = \RV(\mathbb{S})$.
\end{proof}

\setcounter{theorem}{11}
\begin{lemma}[Internalizing finite sets]
\label{finite-VF-project-RV}
\lean{DiamondIntegral.Section2.lemma_2_12, DiamondIntegral.Section2.GammaStable.lemma_2_12_gammaNames}
\leanok
\uses{nota-sub}
For any finite definable set $A$ in $\VF$ there is a definable bijection $f : A \fun U$ for some set $U$ in $\RV$.
\end{lemma}

\setcounter{theorem}{12}
\begin{lemma}[Algebraicity over $\RV$-parameters]
\label{dcl-to-ac}
\lean{DiamondIntegral.Section2.lemma_2_13}
\leanok
For $a, b \in \VF$ and $t \in \RV$, if $b$ is algebraic over $(a,t)$ then it is algebraic over $a$.
\end{lemma}

\setcounter{theorem}{13}
\begin{lemma}[Exchange principle]
\label{exchange}
\lean{DiamondIntegral.Section2.ExchangeGeneral.lemma_2_14_vf_general, DiamondIntegral.Section2.lemma_2_14_rv, DiamondIntegral.Section2.lemma_2_14_vf}
\leanok
The exchange principle holds for single elements in both sorts, with all algebraic closures taken over the same fixed parameter base:
\begin{itemize}
\item for $a, b \in \VF$, if $a \in \operatorname{acl}(b) \mi \operatorname{acl}(\0)$ then $b \in \operatorname{acl}(a)$;
\item for $t, s \in \RV$, if $t \in \operatorname{acl}(s) \mi \operatorname{acl}(\0)$ then $s \in \operatorname{acl}(t)$.
\end{itemize}
\end{lemma}

\setcounter{theorem}{14}
\begin{lemma}[Infinite fibers of finitary functions]
\label{finite-fib-def}
\lean{DiamondIntegral.Section2.lemma_2_15_vf_general, DiamondIntegral.Section2.lemma_2_15_rv}
\leanok
\uses{s2-relations-as-functions, nota-sub}
Let $f : A \fun \mathcal{P}(B)$ be a definable finitary function, where $A$, $B$ are subsets of $\VF$. Then there is a finite definable subset $B' \sub B$ such that, for every $b \in B$, the fiber $A_b = \{a \in A : b \in f(a)\}$ is infinite if and only if $b \in B'$. The same holds if $A$, $B$ are subsets of $\RV$.
\end{lemma}

\begin{proof}
\uses{exchange}
\leanok
If $A_b$ is infinite, compactness gives $a \in A_b \mi \operatorname{acl}(b)$; since $b \in \operatorname{acl}(a)$, Lemma 2.14 forces $b \in \operatorname{acl}(\0)$. A second compactness argument places all such $b$ in an algebraic subset of $B$, which is then trimmed to exactly the $b$ with $A_b$ infinite; the $\RV$ case is the same.
\end{proof}

\setcounter{theorem}{15}
\begin{definition}[$\VF$-dimension]
\label{s2-vf-dimension}
\lean{DiamondIntegral.Section2.vfDimension, DiamondIntegral.Section2.HasVFInteriorDim, DiamondIntegral.Section2.vfDimension_empty}
\leanok
\uses{indexing, rvfiber}
Set $\dim_{\VF}(\varnothing) = -\infty$. For a nonempty definable set $A$, the $\VF$-dimension $\dim_{\VF}(A)$ is the largest natural number $k$ such that, possibly after re-indexing the $\VF$-coordinates, $\operatorname{pr}_{\leq k}(A_t)$ has nonempty interior for some $t \in A_{\RV}$.
\end{definition}

\setcounter{theorem}{16}
\begin{lemma}[Full dimension and open polydiscs]
\label{full-dim-open-poly}
\lean{DiamondIntegral.Section2.lemma_2_17}
\leanok
\uses{s2-vf-dimension, defn-disc}
Let $A$ be a definable set with $A_{\VF} \sub \VF^n$. Then $\dim_{\VF}(A) = n$ if and only if there is a $t \in A_{\RV}$ such that $A_t$ contains an open polydisc.
\end{lemma}

\setcounter{theorem}{17}
\begin{lemma}[Algebraicity over the name of a polydisc]
\label{ball-to-ac}
\lean{DiamondIntegral.Section2.lemma_2_18_finite_names}
\leanok
\uses{s2-vf-dimension, defn-disc, disc-exp}
Let $b \in \VF$ and let $\mathfrak{a} \sub \VF^n$ be a polydisc with $\dim_{\VF}(\mathfrak{a}) = n$. If $b$ is algebraic over $\ulcorner \mathfrak{a} \urcorner$ then $b$ is algebraic.
\end{lemma}

\begin{proof}
\uses{dcl-to-ac, exchange}
\leanok
Put $\alpha = \rad(\mathfrak{a})$. For every $a \in \mathfrak{a}$, $b$ is algebraic over $(a,\alpha)$, hence over $a$ by Lemma 2.13. Choosing $a \in \mathfrak{a}$ with $a_i \notin \operatorname{acl}(a_{\tilde i}, b)$ for all $i \in [n]$, Lemma 2.14 shows that $b$ is algebraic.
\end{proof}

\setcounter{theorem}{18}
\begin{lemma}[Points in algebraic closed discs]
\label{effectiveness}
\lean{DiamondIntegral.Section2.lemma_2_19_algebraic, DiamondIntegral.Section2.lemma_2_19_definable}
\leanok
\uses{defn-disc}
If $\mathbb{S}$ is $(\VF, \RV, \Gamma)$-generated then every algebraic (respectively, definable) closed disc contains an algebraic (respectively, definable) point.
\end{lemma}

\setcounter{theorem}{19}
\begin{lemma}[Points in algebraic discs]
\label{algebraic-balls-definable-centers}
\lean{DiamondIntegral.Section2.lemma_2_20_algebraic, DiamondIntegral.Section2.lemma_2_20_definable, DiamondIntegral.Section3.lemma_2_20_algebraic_small, DiamondIntegral.Section3.lemma_2_20_definable_small}
\leanok
\uses{defn-disc}
If $\mathbb{S}$ is $(\VF, \Gamma)$-generated then every algebraic (respectively, definable) disc contains an algebraic (respectively, definable) point.

\emph{Formalization.} The base is a $\VF$-generated field base with $\Gamma$-names. The Section 2 declarations assume a size bound on the generated closure base; the Section 3 versions \texttt{lemma\_2\_20\_algebraic\_small} and \texttt{lemma\_2\_20\_definable\_small} assume only that the base and the $\Gamma$-names are small (section 2/GAP-REPAIRS.md, November 18).
\end{lemma}

\setcounter{theorem}{20}
\begin{corollary}[Definable points in split discs]
\label{open-disc-def-point}
\lean{DiamondIntegral.Section2.lemma_2_21}
\leanok
\uses{defn-disc}
Suppose that $\mathbb{S}$ is $(\VF, \RV, \Gamma)$-generated. Let $\mathfrak{a} \sub \VF$ be a definable disc and $A$ a definable subset of $\VF$. If $\mathfrak{a} \cap A$ is a nonempty proper subset of $\mathfrak{a}$ then $\mathfrak{a}$ contains a definable point.
\end{corollary}

\begin{proof}
\uses{effectiveness}
\leanok
By $C$-minimality, $\mathfrak{a}$ then contains a definable closed disc, which contains a definable point by Lemma 2.19.
\end{proof}

\setcounter{theorem}{21}
\begin{definition}[Complete types and atomic sets]
\label{s2-atomic}
\lean{DiamondIntegral.Section2.GeneratesCompleteType, DiamondIntegral.Section2.IsCompleteTypeDefinable, DiamondIntegral.Section2.IsAtomic, DiamondIntegral.Section2.isAtomic_iff_definable_completeTypeDefinable}
\leanok
Let $D$ be a set of parameters. A (not necessarily definable) nonempty set $A$ generates a (complete) $D$-type if, for every $D$-definable set $B$, either $A \sub B$ or $A \cap B = \0$. In that case $A$ is $D$-type-definable if no set properly containing $A$ also generates a $D$-type. If $A$ is $D$-definable and generates a $D$-type, or equivalently is both $D$-definable and $D$-type-definable, then $A$ is $D$-atomic, or atomic over $D$. Atomic means $\0$-atomic.
\end{definition}

\setcounter{theorem}{22}
\begin{lemma}[Atomicity after naming $\Gamma$-values]
\label{atomic-base-change}
\lean{DiamondIntegral.Section2.lemma_2_23, DiamondIntegral.Section3.integration_lemma_2_23, DiamondIntegral.Section2.lemma_2_23_disc_names}
\leanok
\uses{s2-atomic}
Suppose that $\Gamma(\mathbb{S})$ is nontrivial. Let $A$ be an atomic set. Then $A$ is $\gamma$-atomic for all $\gamma \in \Gamma$.

\emph{Formalization.} \texttt{lemma\_2\_23} takes the cited Lemma 3.3 of Yin's Integration in ACVF as an explicit argument; \texttt{Section3.integration\_lemma\_2\_23} supplies it from the formalized upstream package (section 2/verification/CITED-INPUTS.md and UPSTREAM-LINKS.md).
\end{lemma}

\setcounter{theorem}{23}
\begin{lemma}[Atomic open polydiscs over $\RV$-parameters]
\label{at-ball-t}
\lean{DiamondIntegral.Section2.lemma_2_24, DiamondIntegral.Section2.lemma_2_24_disc_names, DiamondIntegral.Section3.lemma_2_24_disc_names_discharged}
\leanok
\uses{s2-atomic, defn-disc}
Let $\mathfrak{a} \sub \VF^n$ be an atomic open polydisc. Then $\mathfrak{a}$ is $t$-atomic for all $t \in \RV$.

\emph{Formalization.} The case without disc parameters is proved directly. The case with finitely many named discs (\texttt{lemma\_2\_24\_disc\_names}) takes Lemma 3.9 of Integration in ACVF as an argument, discharged in \texttt{Section3.lemma\_2\_24\_disc\_names\_discharged} (section 2/COVERAGE.md).
\end{lemma}

\begin{proof}
\uses{open-disc-def-point, dcl-to-ac}
\leanok
If disc parameters have been named, this follows from Lemma 3.9 of Yin's Integration in ACVF. Otherwise, for $n = 1$, failure of $t$-atomicity would give, by Corollary 2.21 and Lemma 2.13, a point of $\mathfrak{a}$ algebraic over $\mathbb{S}$, and atomicity would put the infinite disc in a finite definable set. For $n > 1$, induction on $n$ shows that $\mathfrak{a}_1$ is $t$-atomic and $\mathfrak{a}_{>1}$ is $(a,t)$-atomic for $a \in \mathfrak{a}_1$.
\end{proof}

\setcounter{theorem}{24}
\begin{lemma}[Polydiscs in a type are atomic over their names]
\label{atom-self}
\lean{DiamondIntegral.Section2.lemma_2_25, DiamondIntegral.Section2.lemma_2_25_closed, DiamondIntegral.Section2.lemma_2_25_open, DiamondIntegral.Section3.lemma_2_25_discharged}
\leanok
\uses{s2-atomic, defn-disc, disc-exp}
Suppose that $\Gamma(\mathbb{S})$ is nontrivial. Let $A \sub \VF^n$ be a set that generates a type and $\mathfrak{b} \sub A$ an open (or closed) polydisc. Then $\mathfrak{b}$ is $\ulcorner \mathfrak{b} \urcorner$-atomic.

\emph{Formalization.} The Section 2 declarations take the unary Lemmas 3.3 and 3.4 of Integration in ACVF as arguments; \texttt{Section3.lemma\_2\_25\_discharged} is input-free, with saturation stated as a hypothesis (section 2/GAP-REPAIRS.md).
\end{lemma}

\begin{proof}
\uses{atomic-base-change}
\leanok
Induction on $n$, the case $n = 1$ being Lemma 3.4 of Yin's Integration in ACVF. For $n > 1$ write $\mathfrak{b} = \mathfrak{b}_1 \times \mathfrak{b}'$ and $\gamma = \rad(\mathfrak{b}_1)$; for $a \in \mathfrak{b}_1$ the fiber $A_a$ generates a type, so by Lemma 2.23 and induction $\mathfrak{b}'$ is $(a,\gamma,\ulcorner \mathfrak{b}' \urcorner)$-atomic, hence $\ulcorner \mathfrak{b} \urcorner$-atomic. Similarly $\mathfrak{b}_1$ is $\ulcorner \mathfrak{b} \urcorner$-atomic, so $\mathfrak{b}$ is.
\end{proof}

\setcounter{theorem}{25}
\begin{lemma}[$\RV$-values on atomic open polydiscs]
\label{at-rv-split}
\lean{DiamondIntegral.Section2.lemma_2_26, DiamondIntegral.Section2.lemma_2_26_disc_names, DiamondIntegral.Section3.lemma_2_26_disc_names_discharged}
\leanok
\uses{s2-atomic, defn-disc, s2-relations-as-functions}
Let $\mathfrak{a} \sub \VF^n$ be an atomic open polydisc and $f : \mathfrak{a} \fun \mathcal{P}(\VF)$ a definable finitary function such that $\rv \rest f(a)$ is injective for every $a \in \mathfrak{a}$. Then $\rv \circ f$ is constant.

\emph{Formalization.} The version over finitely many named discs (\texttt{lemma\_2\_26\_disc\_names}) takes Lemma 3.9 of Integration in ACVF as an argument, discharged in \texttt{Section3.lemma\_2\_26\_disc\_names\_discharged}.
\end{lemma}

\begin{proof}
\uses{at-ball-t}
\leanok
By Lemma 2.24, $\mathfrak{a}$ is $t$-atomic for every $t \in \rv(\bigcup f(\mathfrak{a}))$, and $f$ induces a $t$-definable function $\mathfrak{a} \fun t^\sharp$.
\end{proof}

\setcounter{theorem}{26}
\begin{lemma}[Images of atomic discs and thin annuli]
\label{hae-component-11}
\lean{DiamondIntegral.Section2.lemma_2_27, DiamondIntegral.Section3.integration_lemma_2_27}
\leanok
\uses{s2-atomic, defn-disc}
A thin annulus is a set $\mathfrak{c}(a,\gamma) \mi \mathfrak{o}(a,\gamma) = \{x \in \VF : \vv(x - a) = \gamma\}$ with $a \in \VF$ and $\gamma \in \Gamma$. Let $A \sub \VF$ be an atomic open disc, an atomic closed disc or an atomic thin annulus, and let $f : A \fun \VF$ be a definable function. Then $f(A)$ is also atomic of one of these three forms.

\emph{Formalization.} \texttt{lemma\_2\_27} proves atomicity of the image and takes the shape classification of Lemma 3.8 of Yin's Integration in ACVF as an argument; \texttt{Section3.integration\_lemma\_2\_27} supplies it from the upstream package (section 2/verification/CITED-INPUTS.md).
\end{lemma}

\setcounter{theorem}{27}
\begin{lemma}[Images of atomic open discs]
\label{atomic-open-ball-open}
\lean{DiamondIntegral.Section2.lemma_2_28, DiamondIntegral.Section3.integration_lemma_2_28}
\leanok
\uses{s2-atomic, defn-disc}
Let $A \sub \VF$ be an atomic open disc and $f : A \fun \VF$ a definable function. Then either $f$ is constant or $f(A)$ is also an atomic open disc.

\emph{Formalization.} \texttt{lemma\_2\_28} proves atomicity of the image and takes the geometric conclusion of Lemma 3.13 of Yin's Integration in ACVF as an argument; \texttt{Section3.integration\_lemma\_2\_28} supplies it from the upstream package (section 2/verification/CITED-INPUTS.md).
\end{lemma}

\setcounter{theorem}{28}
\begin{lemma}[Branches on an atomic open polydisc]
\label{atomic-open-multidisc}
\lean{DiamondIntegral.Section2.lemma_2_29, DiamondIntegral.Section2.lemma_2_29_disc_names, DiamondIntegral.Section3.lemma_2_29_discharged}
\leanok
\uses{at-rv-split, s2-atomic, defn-disc, s2-relations-as-functions, s2-rv-infinity, gamwhat}
Let $\mathfrak{a} \sub \VF^n$ be an atomic open polydisc and $f : \mathfrak{a} \fun \mathcal{P}(\VF)$ a definable finitary function such that $\rv \rest f(a)$ is injective for every $a \in \mathfrak{a}$ (as in Lemma 2.26), and suppose that $f$ is not constant. Then there are $t_1, \ldots, t_k \in \RV_\infty$ and, for each $i$, a $t_i$-definable function $f_i : \mathfrak{a} \fun t_i^\sharp$ such that $f = \bigcup_i f_i$ and each $f_i(\mathfrak{a})$ is either a point or a $t_i$-atomic open disc.

\emph{Formalization.} \texttt{lemma\_2\_29} takes the unary Lemma 3.13 of Integration in ACVF as an argument; \texttt{Section3.lemma\_2\_29\_discharged} is input-free (section 2/GAP-REPAIRS.md).
\end{lemma}

\begin{proof}
\uses{at-rv-split, at-ball-t, atomic-open-ball-open}
\leanok
By Lemma 2.26, $f$ splits into finitely many $t$-definable branches $g : \mathfrak{a} \fun t^\sharp$, constant ones being of the required form. For a nonconstant branch, $\mathfrak{a}$ is $t$-atomic by Lemma 2.24 and each coordinate line has image a point or an open disc by Lemma 2.28, so $g(\mathfrak{a})$ is open. Coordinate chains and $C$-minimality show that $g(\mathfrak{a})$ is a single maximal open disc, which is $t$-atomic.
\end{proof}

\setcounter{theorem}{29}
\begin{corollary}[Partition into polydisc branches]
\label{part-rv-cons}
\lean{DiamondIntegral.Section2.corollary_2_30, DiamondIntegral.Section2.corollary_2_30_disc_names, DiamondIntegral.Section3.integration_corollary_2_30}
\leanok
\uses{at-rv-split, s2-relations-as-functions, defn-disc, disc-exp, s2-rv-infinity, gamwhat}
Let $A \sub \VF^n$ be a definable set and $f : A \fun \mathcal{P}(\VF)$ a definable finitary function such that $\rv \rest f(a)$ is injective for every $a \in A$ (as in Lemma 2.26). Then there is a definable finite partition $(A_i)_i$ of $A$ such that, for each $i$ and every open polydisc $\mathfrak{a} \sub A_i$, there are finitely many $(\ulcorner \mathfrak{a} \urcorner, t_{\mathfrak{a} j})$-definable functions $f_{\mathfrak{a} j} : \mathfrak{a} \fun t^\sharp_{\mathfrak{a} j}$, where $t_{\mathfrak{a} j} \in \RV_\infty$, with $f \rest \mathfrak{a} = \bigcup_j f_{\mathfrak{a} j}$ and each $f_{\mathfrak{a} j}(\mathfrak{a})$ either a point or an open disc.

\emph{Formalization.} \texttt{corollary\_2\_30} takes the unary Lemmas 3.4, 3.9 and 3.13 of Integration in ACVF as arguments; \texttt{Section3.integration\_corollary\_2\_30} supplies them. The partition is over a small real base; \texttt{corollary\_2\_30\_disc\_names} also retains a fixed finite tuple of disc names (section 2/COVERAGE.md).
\end{corollary}

\begin{proof}
\uses{atom-self, atomic-open-multidisc}
\leanok
The required property $P(d)$ of an open-polydisc code $d$ is first-order. For $a \in A$, Lemma 2.25 makes every open polydisc in the type-definable set of $a$ atomic over its name, so Lemma 2.29 gives $P(d)$ for it. Two compactness arguments then give finitely many definable sets covering $A$ on which $P$ holds for all polydiscs; their partition works, the branches being the $\rv$-tagged pieces of $f$.
\end{proof}

\setcounter{theorem}{30}
\begin{definition}[Concentricity and the disc-to-disc property]
\label{defn-dtdp}
\lean{DiamondIntegral.Section2.FibreConcentric, DiamondIntegral.Section2.HasDiscToDiscProperty}
\leanok
\uses{rvfiber, defn-disc}
Let $f : A \fun B$ be a bijection between two sets $A$ and $B$, each with exactly one $\VF$-coordinate. Then $f$ is concentric if, for every $\VF$-fiber $f_t$ of $f$ and all open polydiscs $\mathfrak{a} \sub \operatorname{dom}(f_t)$, $f_t(\mathfrak{a})$ is also an open polydisc. If both $f$ and $f^{-1}$ are concentric then $f$ has the disc-to-disc property.
\end{definition}

\setcounter{theorem}{31}
\begin{lemma}[Disc-to-disc partition]
\label{open-pro}
\lean{DiamondIntegral.Section2.lemma_2_32, DiamondIntegral.Section3.integration_lemma_2_32}
\leanok
\uses{defn-dtdp}
Let $f : A \fun B$ be a definable bijection between two sets $A$ and $B$, each with exactly one $\VF$-coordinate. Then there is a definable finite partition $(A_i)_i$ of $A$ such that every restriction $f \rest A_i$ has the disc-to-disc property.

\emph{Formalization.} \texttt{lemma\_2\_32} takes Propositions 3.19 and 3.21 of Yin's Integration in ACVF as an argument; \texttt{Section3.integration\_lemma\_2\_32} supplies it under saturation (section 2/verification/CITED-INPUTS.md).
\end{lemma}

\setcounter{theorem}{32}
\begin{definition}[$\RV$-dimension]
\label{s2-rv-dimension}
\lean{DiamondIntegral.Section2.rvDimension, DiamondIntegral.Section2.RVHasFiniteMap, DiamondIntegral.Section2.rvDimension_empty}
\leanok
\uses{s2-rv-infinity}
Set $\dim_{\RV}(\varnothing) = -\infty$. For a nonempty definable subset $U \sub \RV_\infty^m$, the $\RV$-dimension $\dim_{\RV}(U)$ is the smallest number $k$ such that there is a definable finite-to-one function $f : U \fun \RV^k$, where $\RV^0$ is the singleton $\{\infty\}$.
\end{definition}

\setcounter{theorem}{33}
\begin{definition}[$\RV$-boundary and $\RV$-interior]
\label{s2-rv-boundary}
\lean{DiamondIntegral.Section2.rvBoundary, DiamondIntegral.Section2.rawRVInterior, DiamondIntegral.Section2.isRVPullback_iff_boundary_empty}
\leanok
\uses{gamwhat, defn-disc}
Let $A$ be a definable subset of $\VF^n$. The $\RV$-boundary $\partial_{\RV} A$ is the definable subset of $\rv(A)$ such that $t \in \partial_{\RV} A$ if and only if $t^\sharp \cap A$ is a proper nonempty subset of $t^\sharp$. The definable set $\operatorname{int}_{\RV}(A) = \rv(A) \mi \partial_{\RV} A$ is the $\RV$-interior of $A$. Thus $A$ is an $\RV$-pullback if and only if $\partial_{\RV} A$ is empty.
\end{definition}

\setcounter{theorem}{34}
\begin{lemma}[Dimension of the $\RV$-boundary]
\label{RV-bou-dim}
\lean{DiamondIntegral.Section2.lemma_2_35, DiamondIntegral.Section2.lemma_2_35_vf, DiamondIntegral.Section3.lemma_2_35_discharged}
\leanok
\uses{s2-rv-dimension, s2-rv-boundary}
Let $A$ be a definable subset of $\VF^n$. Then $\dim_{\RV}(\partial_{\RV} A) < n$.

\emph{Formalization.} The Section 2 declarations take the unary Lemmas 3.4, 3.9 and 3.13 of Integration in ACVF (through Corollary 2.30) as arguments; \texttt{Section3.lemma\_2\_35\_discharged} is input-free (section 2/GAP-REPAIRS.md).
\end{lemma}

\begin{proof}
\uses{open-disc-def-point, part-rv-cons}
\leanok
Induction on $n$; for $n = 1$ a finite Boolean disc decomposition leaves only finitely many partially met $\RV$-fibers. In the inductive step, Corollary 2.21 gives definable representatives of the boundary tags of the unary sections, and Corollary 2.30 makes their $\RV$-tags constant on the $\RV$-polydiscs of a partition of $\operatorname{pr}_{\tilde i}(A)$. The boundary is then covered by graphs finite over $n-1$ coordinates and by sets controlled by the induction hypothesis, both of $\RV$-dimension less than $n$.
\end{proof}

\setcounter{theorem}{35}
\begin{lemma}[Continuity almost everywhere]
\label{fun-almost-cont}
\lean{DiamondIntegral.Section2.lemma_2_36}
\leanok
\uses{s2-vf-dimension}
Let $f : \VF^n \fun \VF^m$ be a definable function. There is a definable closed set $A \sub \VF^n$ with $\dim_{\VF}(A) < n$ such that $f \rest (\VF^n \mi A)$ is continuous; that is, $f$ is continuous almost everywhere.
\end{lemma}

\setcounter{theorem}{36}
\begin{definition}[Derivatives and partial derivatives]
\label{defn-diff}
\lean{DiamondIntegral.Section2.HasDerivativeAt, DiamondIntegral.Section2.HasPartialDerivativeAt, DiamondIntegral.Section2.HasDerivativeAt.unique, DiamondIntegral.Section2.HasDerivativeAt.partial_matrix_entry}
\leanok
\uses{indexing}
Let $f : \VF^n \fun \VF^m$ be a definable function; for a nonzero tuple $x$ put $\vv(x) = \min_i \vv(x_i)$. $f$ is differentiable at $a \in \VF^n$ if there is a linear map $\lambda : \VF^n \fun \VF^m$ such that, for every $\epsilon \in \Gamma$, $\vv(f(a+b) - f(a) - \lambda(b)) - \vv(b) > \epsilon$ whenever $b \in \VF^n \mi \{0\}$ and $\vv(b)$ is sufficiently large; such $\lambda$ is unique and is the derivative of $f$ at $a$. For $f = (f_1, \ldots, f_m)$ and $a = (a_i, a_{\tilde i})$, the derivative of $f_j \rest (\VF \times a_{\tilde i})$ at $a_i$, if it exists, is the $ij$th partial derivative of $f$ at $a$; if $f$ is differentiable at $a$, it is the $(j,i)$-entry of the derivative matrix.
\end{definition}

\setcounter{theorem}{37}
\begin{lemma}[Jacobians exist almost everywhere]
\label{diff-almost-every}
\lean{DiamondIntegral.Section2.lemma_2_38}
\leanok
\uses{defn-diff, s2-vf-dimension, rvfiber}
Let $f : \VF^n \times \RV^m \fun \VF^{n'} \times \RV^{m'}$ be a definable function with $n = n'$. Outside a definable set of $\VF$-dimension less than $n$ the output $\RV$-tuple $s$ is constant near each input $(a,t)$, and the partial derivatives of $f$ at $(a,t)$ are those of the local $\VF$-map $f_{(t,s)}$; the Jacobian $\jcb_{\VF} f(a,t)$ is the determinant of the matrix of these partial derivatives. Then the Jacobian of $f$ exists and is a continuous function almost everywhere.
\end{lemma}

\begin{proof}
\leanok
Remove the locus where the output $\RV$-tuple is not locally constant. On each fiber with fixed input and output $\RV$-tuples, extend the $\VF$-map by zero and apply Corollary 9.9 of Yin's Special transformations to its partial derivatives; their nonexistence and discontinuity loci have $\VF$-dimension less than $n$, and so does the combined exceptional locus. In arity zero the Jacobian is the empty determinant $1$.
\end{proof}

\setcounter{theorem}{38}
\begin{definition}[Bounded and doubly bounded sets]
\label{bounded}
\lean{DiamondIntegral.Section2.valim, DiamondIntegral.Section2.BoundedBy, DiamondIntegral.Section2.DoublyBoundedBy, DiamondIntegral.Section2.IsDoublyBounded}
\leanok
\uses{s2-rv-infinity}
For $\gamma \in \Gamma$, a set $I \sub \Gamma_\infty^n$ is $\gamma$-bounded if $I \sub [\gamma, \infty]^n$, and doubly $\gamma$-bounded if $I \sub [-\gamma, \gamma]^n$. For $A \sub \VF^n \times \RV_\infty^m \times \Gamma_\infty^l$ let
\[ \operatorname{valim}(A) = \{(\vv(a), \vrv(t), \gamma) : (a, t, \gamma) \in A\} \sub \Gamma_\infty^{n+m+l}, \]
with $\vv(a)$ and $\vrv(t)$ taken coordinatewise; $A$ is (doubly) $\gamma$-bounded if $\operatorname{valim}(A)$ is. A set is bounded (doubly bounded) if it is (doubly) $\gamma$-bounded for some $\gamma \in \Gamma$.
\end{definition}

\setcounter{theorem}{39}
\begin{definition}[Contractions]
\label{defn-corr-cont}
\lean{DiamondIntegral.Section2.IsRVContractible, DiamondIntegral.Section2.rvContraction, DiamondIntegral.Section2.IsContractibleAlong, DiamondIntegral.Section2.contractionAlong_unique}
\leanok
A function $f : A \fun B$ is $\rv$-contractible if there is a (necessarily unique) function $f_{\downarrow} : \rv(A) \fun \rv(B)$, the $\rv$-contraction of $f$, such that $(\rv \rest B) \circ f = f_{\downarrow} \circ (\rv \rest A)$. Similarly, $f$ is $\operatorname{res}$-contractible (respectively, $\vv$-contractible) if the same holds with $\operatorname{res}$ (respectively, $\vv$ or $\vrv$, depending on the coordinates) in place of $\rv$.

\emph{Formalization.} The $\rv$-contraction has its own definition; the $\operatorname{res}$- and $\vv$-variants exist only through the generic coordinate-map definition \texttt{IsContractibleAlong}. Formalization-session choice, no author decision on record (FINAL-CHECK.md, section 3).
\end{definition}

\setcounter{theorem}{40}
\begin{lemma}[Definable images of doubly bounded sets]
\label{db-to-db}
\lean{DiamondIntegral.Section2.lemma_2_41, DiamondIntegral.Section2.lemma_2_41_gamma, DiamondIntegral.Section3.lemma_2_41_discharged, DiamondIntegral.Section3.lemma_2_41_gamma_discharged}
\leanok
\uses{bounded}
Let $U \sub \RV^k$ be a doubly bounded set and $f : U \fun \RV$ a definable function. Then $f(U)$ is also doubly bounded. The same holds if the codomain of $f$ is $\Gamma$.

\emph{Formalization.} The Section 2 declarations take piecewise rational-affinity of definable $\Gamma$-functions as an argument; it is proved from $\RV$ quantifier elimination (\texttt{Section5.GammaQE.gamma\_piecewise\_rational\_affine}), and the Section 3 \texttt{\_discharged} versions are input-free (section 2/GAP-REPAIRS.md).
\end{lemma}

\begin{proof}
\uses{defn-corr-cont}
\leanok
Pass to the $\Gamma$-valued $h = \vrv \circ f$ (or $h = f$). By induction on $k$, the fiberwise bounds $\gamma_t = \sup\{|h(t,x)| : x \in U_t\}$ exist by definable completeness of $\Gamma$ and form a definable function of $t \in \operatorname{pr}_1(U)$, bounded by the unary case. In the unary case, $C$-minimality lets $g$ be $\vrv$-contractible, and piecewise $\mathbb{Q}$-linearity of definable functions in $\Gamma$ makes the range of $g_{\downarrow}$ doubly bounded.
\end{proof}

\setcounter{theorem}{41}
\begin{definition}[Limit sets]
\label{s2-limit-set}
\lean{DiamondIntegral.Section2.HasLimitSet, DiamondIntegral.Section2.HasMinimalLimitSet, DiamondIntegral.Section2.HasMinimalLimitSet.isolated}
\leanok
\uses{defn-disc, s2-relations-as-functions}
Let $A \sub \VF^n$ and let $f : A \fun \mathcal{P}(\VF^m)$ be a definable function. For $a \in \VF^n$ and $L \sub \VF^m$, $L$ is a limit set of $f$ at $a$, written $\lim_{A \to a} f \sub L$, if for every $\epsilon \in \Gamma$ there is a $\delta \in \Gamma$ such that $f(c) \sub \bigcup_{b \in L} \mathfrak{o}(b, \epsilon)$ for every $c \in \mathfrak{o}(a,\delta) \cap (A \mi \{a\})$. A limit set $L$ of $f$ at $a$ is minimal if no proper subset of $L$ is a limit set of $f$ at $a$. In a minimal limit set every element is isolated.
\end{definition}

\setcounter{theorem}{42}
\begin{lemma}[Closures of minimal limit sets coincide]
\label{s2-minimal-limit-closure}
\lean{DiamondIntegral.Section2.lemma_2_43}
\leanok
\uses{s2-limit-set}
Let $A \sub \VF^n$, let $f : A \fun \mathcal{P}(\VF^m)$ be a definable function and $a \in \VF^n$. If $L_1, L_2 \sub \VF^m$ are two minimal limit sets of $f$ at $a$, then the closures of $L_1$ and $L_2$ coincide. One writes $\lim_{A \to a} f = L$ when $L$ is a closed, hence the unique, minimal limit set of $f$ at $a$.
\end{lemma}

\setcounter{theorem}{43}
\begin{lemma}[Finitary functions are almost everywhere continuous]
\label{fini-conti}
\lean{DiamondIntegral.Section2.lemma_2_44, DiamondIntegral.Section2.lemma_2_44_unary}
\leanok
\uses{s2-limit-set, s2-vf-dimension, s2-relations-as-functions}
Let $A \sub \VF^n$ be a definable set with $\dim_{\VF}(A) = n > 0$ and $f : A \fun \mathcal{P}(\VF^m)$ a definable finitary function. Call $f$ continuous at $a \in A$ if $\lim_{A \to a} f = f(a)$. Then $f$ is continuous almost everywhere, that is, outside a definable set of $\VF$-dimension less than $n$.
\end{lemma}

\begin{proof}
\uses{finite-VF-project-RV, full-dim-open-poly, fun-almost-cont}
\leanok
By Lemma 2.12 and compactness, $f$ is a union of definable branches $f_t$ indexed by a set in $\RV$. The discontinuity locus $E$ is definable; if $\dim_{\VF}(E) = n$, Lemma 2.17 and generic local constancy give an open polydisc in $E$ on which the set of branch labels is constant, and Lemma 2.36 applied to the finitely many branches gives a point where $f$ is continuous, a contradiction.
\end{proof}

\setcounter{theorem}{44}
\begin{lemma}[Existence of finite limits]
\label{lim-exists}
\lean{DiamondIntegral.Section2.lemma_2_45}
\leanok
\uses{s2-limit-set, s2-relations-as-functions, defn-disc}
Suppose that $f : \mathcal{M} \mi 0 \fun \mathcal{P}(A)$ is a definable finitary function, where $A$ is a closed and bounded set in $\VF$. Then there is a definable finite set $L \sub A$ such that $\lim_{\mathcal{M} \mi 0 \to 0} f = L$.
\end{lemma}

\setcounter{theorem}{45}
\begin{lemma}[Curve selection]
\label{curve-select}
\lean{DiamondIntegral.Section2.lemma_2_46, DiamondIntegral.Section2.lemma_2_46_affine}
\leanok
\uses{s2-limit-set, s2-relations-as-functions, defn-disc}
Suppose that $\mathbb{S}$ is $\VF$-generated. Let $A \sub \VF^n$ be definable and let $b$ be a point in its closure in the coordinatewise projective line $(\mathbb{P}^1(\VF))^n$, with the valuation topology. Then there is a $b$-definable continuous finitary function $g : \mathcal{M} \mi \{0\} \fun \mathcal{P}(A)$ whose values are nonempty and such that $\lim_{\mathcal{M} \mi \{0\} \to 0} g = \{b\}$.
\end{lemma}

\begin{proof}
\uses{fini-conti}
\leanok
A projective chart reduces to $b \in \VF^n$. Since $\operatorname{acl}(\mathbb{S}, b, c)$ is a model of $\ACVF(0,0)$, compactness gives a $b$-definable finitary $g_0$ with $\varnothing \ne g_0(c) \sub A \cap \mathfrak{o}(b, \vv(c))$, which converges to $\{b\}$. By Lemma 2.44, $g_0$ is continuous off a finite set; keep $g_0$ on a $b$-definable clopen punctured disc around $0$ where it is continuous and use a finite nonempty $b$-definable subset of $A$ elsewhere.
\end{proof}

\setcounter{theorem}{46}
\begin{remark}[Curve selection fails in $\Gamma_\infty$]
\label{s2-gamma-curve-selection-fails}
\lean{DiamondIntegral.Section2.remark_2_47, DiamondIntegral.Section2.remark_2_47_no_approach, DiamondIntegral.Section2.remark_2_47_diagonal, DiamondIntegral.Section3.remark_2_47_no_approach_discharged}
\leanok
\uses{gamwhat}
Although the $\Gamma$-sort is essentially an o-minimal group, curve selection fails because of the top element $\infty$: the set $\Gamma \times (0, \infty) \sub \Gamma_\infty^2$ contains curves approaching the boundary points $(0,0)$ and $(\infty,\infty)$, but no definable curve approaching the boundary point $(\infty, 0)$.

\emph{Formalization.} The negative claim (no definable curve approaches $(\infty,0)$) is formalized, with input-free versions in Section 3. For the positive claim, \texttt{remark\_2\_47\_diagonal} only records that the diagonal has points near $(0,0)$ and $(\infty,\infty)$; the diagonal case was not pursued, by the author-approved triage of 2026-11-18 (section 2/GAP-REPAIRS.md, FINAL-CHECK.md section 3).
\end{remark}

\begin{proof}
Every definable one-parameter curve in $\Gamma$ is piecewise rational-affine. Near a one-sided endpoint one affine piece remains, and a positive affine coordinate cannot tend to $0$ while the other coordinate tends to $\infty$.
\end{proof}

\setcounter{theorem}{47}
\begin{lemma}[Continuous images of closed bounded sets]
\label{cb-to-cb}
\lean{DiamondIntegral.Section2.lemma_2_48}
\leanok
Let $A \sub \VF^n$ be definable, closed and bounded, and let $f : A \fun \VF^m$ be definable and continuous. Then $f(A)$ is also closed and bounded.
\end{lemma}

\begin{proof}
\uses{lim-exists}
\leanok
Temporarily enlarge the base to a small $\VF$-generated base, over which algebraic closures are models, so compactness gives finite nonempty definable choices. For $b$ in the closure of $f(A)$, choices $F(c) \sub A$ with $f(F(c)) \sub \mathfrak{o}(b, \vv(c))$ have a nonempty finite minimal limit set by Lemma 2.45, and continuity forces $f(a) = b$ at its points. Unboundedness is excluded in the same way with choices satisfying $\vv(f_i(x)) < -\vv(c)$.
\end{proof}

\setcounter{theorem}{48}
\begin{corollary}[Injective continuous maps are homeomorphisms]
\label{conti-homeo}
\lean{DiamondIntegral.Section2.corollary_2_49}
\leanok
\uses{cb-to-cb}
Let $A \sub \VF^n$ be definable, closed and bounded, and let $f : A \fun \VF^m$ be definable, continuous and injective (as in Lemma 2.48). Then $f$ is a homeomorphism from $A$ onto $f(A)$.
\end{corollary}

\setcounter{theorem}{49}
\begin{lemma}[The closed graph criterion for continuity]
\label{closed-gra}
\lean{DiamondIntegral.Section2.lemma_2_50, DiamondIntegral.Section2.lemma_2_50_bounded}
\leanok
Let $A \sub \VF^n$ and $B \sub \VF^m$ be definable, with $B$ closed in affine space, and let $B^{\mathrm{p}}$ be the closure of $B$ in $(\mathbb{P}^1(\VF))^m$ with the valuation topology. A definable function $f : A \fun B$ is continuous if and only if its graph is closed in $A \times B^{\mathrm{p}}$. If $B$ is bounded then $B^{\mathrm{p}} = B$, so this is the usual relative closed-graph condition.
\end{lemma}

\begin{proof}
\uses{curve-select, lim-exists}
\leanok
Graphs of continuous maps into the Hausdorff space $B^{\mathrm{p}}$ are closed. Conversely, if $f$ is not continuous at $a$, choose a clopen projective polydisc $U$ around $f(a)$ with $a$ in the closure of $\{x : f(x) \notin U\}$; after enlarging to a small $\VF$-generated base, Lemma 2.46 gives a curve into this set converging to $a$, and the projective form of Lemma 2.45 gives a point $b' \notin U$ with $(a, b')$ in the closure of the graph, contradicting closedness.
\end{proof}

\setcounter{theorem}{50}
\begin{lemma}[Weak concentricity]
\label{lim-weak}
\lean{DiamondIntegral.Section2.lemma_2_51}
\leanok
\uses{defn-disc}
Let $f : A \fun B$ be a definable bijection between two open subsets of $\VF$, and suppose that $B$ is bounded. Then $f$ is weakly concentric: for every $a \in A$ there is a $b \in \VF$ such that, for every sufficiently large $\gamma \in \Gamma$, $f(\mathfrak{o}(a,\gamma) \mi a) = \mathfrak{o}(b,\delta) \mi b$ for some $\delta \in \Gamma$.
\end{lemma}

\begin{proof}
\uses{lim-exists}
\leanok
Lemma 2.45, applied with the closure of $B$, and Lemma 9.4 of Yin's Special transformations give a unique limit $b$ of $f$ at $a$. The valuation calculation of Lemma 8.3 there gives a positive rational-affine law $\vv(f(x) - b) = \rho(\vv(x - a))$ near $a$, and likewise for $f^{-1}$ near $b$; the two laws are inverse, so thin annuli around $a$ map onto thin annuli around $b$, and taking unions gives the claim.
\end{proof}

\setcounter{theorem}{51}
\begin{definition}[$\mathfrak{o}$-partitions]
\label{s2-o-partition}
\lean{DiamondIntegral.Section2.IsOpenDiscPartition, DiamondIntegral.Section2.IsMixedOpenDiscPartition}
\leanok
\uses{defn-disc}
Let $p : A \fun \Gamma$ be a definable function. Then $p$ is an $\mathfrak{o}$-partition of $A$ if, for every $a \in A$, $p$ is constant on $\mathfrak{o}(a, p(a)) \cap A$. When $A$ has no $\VF$-coordinates, every definable function $A \fun \Gamma$ is an $\mathfrak{o}$-partition of $A$.
\end{definition}

\setcounter{theorem}{52}
\begin{lemma}[Values of $\mathfrak{o}$-partitions are doubly bounded]
\label{vol-par-bounded}
\lean{DiamondIntegral.Section2.lemma_2_53, DiamondIntegral.Section3.lemma_2_53_discharged}
\leanok
\uses{s2-o-partition, bounded}
Let $p$ be an $\mathfrak{o}$-partition of $A$. Suppose that $A_{\RV}$ is doubly bounded and $A_t$ is closed and bounded for every $t \in A_{\RV}$. Then $p(A)$ is doubly bounded.

\emph{Formalization.} \texttt{lemma\_2\_53} takes piecewise rational-affinity of definable $\Gamma$-functions (used through Lemma 2.41) as an argument; \texttt{Section3.lemma\_2\_53\_discharged} is input-free (section 2/GAP-REPAIRS.md).
\end{lemma}

\begin{proof}
\uses{cb-to-cb, lim-exists, db-to-db}
\leanok
Without $\RV$-coordinates, $p(A)$ is bounded below since $A \sub \mathcal{O}_\eta^n$ for some $\eta$. If it were unbounded above, then over a small $\VF$-generated base (as in Lemma 2.48) finite choices $F(c) \sub \{x : p(x) > \vv(c)\}$ have a minimal limit point $e$ by Lemma 2.45, contradicting constancy of $p$ on $\mathfrak{o}(e, p(e))$. In general the fiberwise bounds form a definable $\Gamma$-valued function on $A_{\RV}$, whose image is doubly bounded by Lemma 2.41.
\end{proof}

\setcounter{theorem}{53}
\begin{lemma}[Fiberwise open and closed sets]
\label{lem-open-fiberwise}
\lean{DiamondIntegral.Section2.lemma_2_54, DiamondIntegral.Section2.lemma_2_54_open, DiamondIntegral.Section2.lemma_2_54_closed}
\leanok
\uses{indexing}
Let $A \sub \VF^n$ be a definable set such that, for each $a \in A_1$, the fiber $A_a \sub \VF^{n-1}$ is open (respectively, closed). Then there is a finite definable set $A' \sub A_1$ such that $A \cap ((A_1 \mi A') \times \VF^{n-1})$ is open (respectively, closed) in $(A_1 \mi A') \times \VF^{n-1}$.
\end{lemma}

\begin{proof}
\uses{cb-to-cb, fini-conti}
\leanok
Prove the open case over a unary base, the closed case following by relative complement. Let $S$ be the set of points of $A$ that are not interior, and $T_\gamma \sub S$ those whose fiber contains the $\gamma$-disc around them; if $S_1$ were infinite, some $(T_\gamma)_1$ would contain an open disc $D$ (Lemma 4.3 of Yin's Special transformations). Finite definable choices in $T_\gamma$ over $D$, over a small $\VF$-generated base as in Lemma 2.48, are continuous at some point by Lemma 2.44, which produces an open box inside $A$ meeting $S$, a contradiction.
\end{proof}

\setcounter{theorem}{54}
\begin{proposition}[Fusion of continuous fibers]
\label{prop-continuity-fiberwise}
\lean{DiamondIntegral.Section2.proposition_2_55}
\leanok
\uses{lem-open-fiberwise, indexing}
Let $A \sub \VF^n$ be a definable set (as in Lemma 2.54) and $f : A \fun \VF$ a definable function such that, for each $a \in A_1$, the induced function $f_a : A_a \fun \VF$ is continuous. Then there is a finite definable set $A' \sub A_1$ such that $f$ is continuous on $\bigcup_{a \in A_1 \mi A'} a \times A_a$.
\end{proposition}

\begin{proof}
\uses{lem-open-fiberwise, closed-gra}
\leanok
Let $G_a$ be the closure of the graph of $f_a$ in $\VF^{n-1} \times \mathbb{P}^1(\VF)$; these form a definable family with closed fibers. Lemma 2.54, applied in the two affine charts of $\mathbb{P}^1(\VF)$, gives a finite exceptional set off which the family is closed; fiberwise continuity makes the graph of $f$ relatively closed there, and the closed graph criterion (Lemma 2.50) gives continuity.
\end{proof}

\section{Hrushovski--Kazhdan style integral}
\label{section-3}
The measured and unmeasured categories of $\VF$, $\RV$ and $\RES$ objects, their Grothendieck semirings and rings, and the Hrushovski--Kazhdan integral as a ring isomorphism onto a quotient of the $\RV$ ring. The tensor comparisons and coefficient maps of Hrushovski and Kazhdan are proved in this repository rather than cited; the injectivity of the coefficient maps uses the Reichstein--Youssin hypothesis.

\setcounter{theorem}{0}
\begin{definition}[Grothendieck semigroup and Grothendieck ring]
\label{def-Kgroup}
\lean{DiamondIntegral.Section3.ScissorPresentation, DiamondIntegral.Section3.KPlus, DiamondIntegral.Section3.RingCompletion}
\leanok
Let $\mathcal{C}$ be a category endowed with a binary operation $\mi$ and a binary relation $\sub$ (subobject), both compatible with isomorphisms of objects (usually set subtraction and set inclusion). The Grothendieck semigroup $\gsk \mathcal{C}$ is the free semigroup generated by the isomorphism classes of $\mathcal{C}$, subject to the scissor relation $[A \mi B] + [B] = [A]$ whenever $B \sub A$, where $[A]$ denotes the isomorphism class of $A$. If $\mathcal{C}$ is also equipped with a binary operation (for example cartesian product) inducing a multiplication, then $\gsk \mathcal{C}$ is a commutative semiring. Its formal groupification, which is then a ring, is denoted by $\ggk \mathcal{C}$.
\end{definition}

\setcounter{theorem}{1}
\begin{definition}[$\RV$-fiber dimension]
\label{s3-rv-fiber-dimension}
\lean{DiamondIntegral.Section3.gammaGroundRVFibreDimension, DiamondIntegral.Section3.rvFibreDimension}
\leanok
\uses{s2-rv-dimension}
An $\RV$-fiber of a definable set $A$ is a set of the form $A_a$ with $a \in A_{\VF}$. The $\RV$-fiber dimension of $A$, denoted by $\dim^{\mathrm{fib}}_{\RV}(A)$, is the maximum of the $\RV$-dimensions of its $\RV$-fibers.
\end{definition}

\setcounter{theorem}{2}
\begin{lemma}[Yin, Special transformations, Lemma 4.13]
\label{RV-fiber-dim-same}
\lean{DiamondIntegral.Section3.gammaGroundLemma_3_3, DiamondIntegral.Section3.lemma_3_3}
\leanok
\uses{s3-rv-fiber-dimension}
Let $f : A \fun A'$ be a definable bijection. Then $\dim^{\mathrm{fib}}_{\RV}(A) = \dim^{\mathrm{fib}}_{\RV}(A')$.
\end{lemma}

\setcounter{theorem}{3}
\begin{definition}[$\VF$-categories]
\label{defn-VF-cat}
\lean{DiamondIntegral.Section3.GammaGroundVFObject}
\leanok
\uses{s3-rv-fiber-dimension, s2-vf-dimension}
The objects of the category $\VF[k]$ are the definable sets of $\VF$-dimension at most $k$ and $\RV$-fiber dimension $0$ (that is, all their $\RV$-fibers are finite). Any definable bijection between two such objects is a morphism of $\VF[k]$. Set $\VF_* = \bigcup_k \VF[k]$.
\end{definition}

\setcounter{theorem}{4}
\begin{definition}[$\VF$-categories with $\Gamma$-volume forms]
\label{defn-VF-mu-G}
\lean{DiamondIntegral.Section3.GammaGroundMeasuredVFData, DiamondIntegral.Section3.GammaGroundEssentialVFMap, DiamondIntegral.Section3.GammaGroundMeasuredVFHom}
\leanok
\uses{defn-VF-cat, defn-diff, diff-almost-every}
An object of $\mVF[k]$ is a definable pair $(A, \omega)$, where $A \in \VF[k]$, $A_{\VF} \sub \VF^k$, and $\omega : A \fun \Gamma$ is a function (a definable $\Gamma$-volume form on $A$). A morphism $(A, \omega) \fun (B, \sigma)$ is a definable essential bijection $F : A \fun B$, that is, a bijection defined outside definable subsets of $A$, $B$ of $\VF$-dimension less than $k$, such that for almost all $x \in A$, $\omega(x) = \sigma(F(x)) + \vv(\jcb_{\VF} F(x))$; such an $F$ is called $\Gamma$-measure-preserving. Set $\mVF[{\leq}k] = \coprod_{i \leq k} \mVF[i]$ and $\mVF[*] = \coprod_k \mVF[k]$.
\end{definition}

\setcounter{theorem}{6}
\begin{remark}[Morphisms up to agreement almost everywhere]
\label{mor-equi}
\lean{DiamondIntegral.Section3.gammaGroundEssentialAESetoid, DiamondIntegral.Section3.gammaGroundEssentialGroupoid, DiamondIntegral.Section3.gammaGroundEssentialQuotientEquiv}
\leanok
\uses{defn-VF-mu-G, def-Kgroup}
For two $\mVF[k]$-morphisms $F, G : (A, \omega) \fun (B, \sigma)$, put $F \sim G$ if $F(a) = G(a)$ for all $a \in A$ outside a definable subset of $\VF$-dimension less than $k$. This is a congruence on $\mVF[k]$; the morphisms of the quotient category $\mVF[k]/{\sim}$ are the classes $[F]$, every one of them is an isomorphism, and so $\mVF[k]/{\sim}$ is a groupoid. This quotient is what is used to form the Grothendieck semigroup of $\mVF[k]$ and to call two objects of $\mVF[k]$ isomorphic.
\end{remark}

\setcounter{theorem}{7}
\begin{definition}[$\RV$-categories]
\label{defn-c-RV-cat}
\lean{DiamondIntegral.Section3.GammaGroundRVObject, DiamondIntegral.Section3.RVObject}
\leanok
The objects of the category $\RV[k]$ are the pairs $(U, f)$ with $U$ a definable set in $\RV_\infty$ and $f : U \fun \RV^k$ a definable finite-to-one function. Given two such objects $(U, f)$, $(V, g)$, any definable bijection $F : U \fun V$ is a morphism of $\RV[k]$. As for the other categories, $\RV[{\leq}k] = \coprod_{i \leq k} \RV[i]$ and $\RV[*] = \coprod_k \RV[k]$.
\end{definition}

\setcounter{theorem}{8}
\begin{definition}[{Objects of $\RV[{\leq}k]$ and projection notation}]
\label{0coor}
\lean{DiamondIntegral.Section3.GammaGroundRVObject.projectionObject, DiamondIntegral.Section3.gammaGroundRVClass_projectionObject, DiamondIntegral.Section3.gammaGroundRVNonzeroStratum}
\leanok
\uses{defn-c-RV-cat, indexing}
If $(U, f) \in \RV[k]$ then $f(U) \sub \RV^k$, not $\RV^k_\infty$, while $\infty$ may occur in any coordinate of $U$; an object of $\RV[*]$ of the form $(U, \operatorname{id})$ is written $U$. More generally, a definable finite-to-one function $f : U \fun \RV^k_\infty$ gives the object $(U, f)$ of $\RV[{\leq}k]$; for example $\{\infty\} \fun \RV_\infty$ gives an object of $\RV[0]$ and $\{(1, \infty)\} \fun \RV^2_\infty$ an object of $\RV[1]$. Every object of $\RV[*]$ is isomorphic to one whose function is a coordinate projection; $(U, \operatorname{pr}_{\leq k})$ is written $U_{\leq k}$ and its class in $\gsk \RV[k]$ is written $[U]_{\leq k}$.
\end{definition}

\setcounter{theorem}{9}
\begin{definition}[$\RES$-categories]
\label{defn-RES-cat}
\lean{DiamondIntegral.Section3.GammaGroundRESObject}
\leanok
\uses{defn-c-RV-cat}
The category $\RES[k]$ is the full subcategory of $\RV[k]$ whose objects are the $(U, f) \in \RV[k]$ with $\vrv(U)$ finite.
\end{definition}

\setcounter{theorem}{10}
\begin{definition}[The correspondence $F^\dagger$ of an $\RV$-morphism]
\label{mor-ftf}
\lean{DiamondIntegral.Section3.GammaGroundRVObject.correspondence, DiamondIntegral.Section3.GammaGroundRVObject.correspondenceJacobian}
\leanok
\uses{defn-c-RV-cat, jcbG-op}
Every $\RV[k]$-morphism $F : (U, f) \fun (V, g)$ induces a definable finite-to-finite correspondence $F^\dagger \sub f(U) \times g(V)$, which may also be viewed as a finitary function $f(U) \fun \mathcal{P}(g(V))$. For $u \in U$, write $\jcb_{\Gamma} F^\dagger(u)$ for $\jcb_{\Gamma} F^\dagger(f(u), g(F(u)))$, the $\Gamma$-Jacobian of Notation 2.8.
\end{definition}

\setcounter{theorem}{11}
\begin{definition}[$\RV$- and $\RES$-categories with $\Gamma$-volume forms]
\label{defn-RV-mu}
\lean{DiamondIntegral.Section3.GammaGroundMeasuredRVObject, DiamondIntegral.Section3.GammaGroundMeasuredRVHom, DiamondIntegral.Section3.GammaGroundMeasuredRESObject}
\leanok
\uses{defn-c-RV-cat, defn-RES-cat, mor-ftf}
An object of $\mRV[k]$ is a definable triple $(U, f, \omega)$, where $(U, f)$ is an object of $\RV[k]$ and $\omega : U \fun \Gamma$ is a function (a definable $\Gamma$-volume form on $(U, f)$). A morphism $(U, f, \omega) \fun (V, g, \sigma)$ is an $\RV[k]$-morphism $F : (U, f) \fun (V, g)$ such that $\omega(u) = \sigma(F(u)) + \jcb_{\Gamma} F^\dagger(u)$ for every $u \in U$. The category $\mRES[k]$ is the full subcategory of $\mRV[k]$ of the objects with $(U, f) \in \RES[k]$.
\end{definition}

\setcounter{theorem}{12}
\begin{definition}[$\Gamma$-categories]
\label{def-Ga-cat}
\lean{DiamondIntegral.Section3.GammaGroundGammaObject, DiamondIntegral.Section3.GammaGroundGammaHom, DiamondIntegral.Section3.GammaGroundFiniteGammaObject}
\leanok
\uses{defn-corr-cont}
The objects of the category $\Gamma[k]$ are the finite disjoint unions of definable subsets of $\Gamma^k$. A definable bijection between two such objects is a morphism of $\Gamma[k]$ if and only if it is the $\vrv$-contraction of a definable bijection. The category $\Gamma^{\mathrm{fin}}[k]$ is the full subcategory of $\Gamma[k]$ of the finite objects.

\emph{Formalization.} Lean reads the witnessing definable bijection as a bijection between the full preimages $I^\sharp$ and $J^\sharp$, the reading confirmed by the author on 2026-11-12 (section 3/STATUS.md, ``Author clarification: intended Gamma morphisms'').
\end{definition}

\setcounter{theorem}{13}
\begin{remark}[Definable values and the divisible hull]
\label{s3-definable-value-hull}
\lean{DiamondIntegral.Section3.remark_3_14, DiamondIntegral.Section3.remark_3_14_converse, DiamondIntegral.Section3.remark_3_14_point}
\leanok
\uses{gamwhat}
If $\gamma \in \Gamma$ is definable then it lies in the divisible hull of $\Gamma(\mathbb{S})$. This does not mean that the definable set $\gamma^\sharp \sub \RV$ contains a definable point unless $\gamma \in \Gamma(\mathbb{S})$.

\emph{Formalization.} Lean proves the first sentence over a base generated, up to algebraic closure, by $\VF$-parameters (with the converse that roots of definable values are definable), and the second only for points of a $\VF$-generated base field. The versions over $\Gamma$- or $\RV$-named bases and for points definable over $\mathbb{S}$ were not pursued, by the author-approved triage of 2026-11-18 (section 3/GAP-REPAIRS.md, FINAL-CHECK.md).
\end{remark}

\begin{assumption}[Reichstein--Youssin classification of diagonal actions]
\label{hyp-reichstein-youssin}
\lean{DiamondIntegral.Section3.ReichsteinYoussinInput, DiamondIntegral.Section3.GammaGroundCoefficientClassification}
\leanok
Let $\K$ be an algebraically closed field of characteristic $0$, $A$ an abelian group and $a, b \in A^n$ whose entries together generate $A$. Suppose that the diagonal actions of the characters $\chi : A \fun \K^\times$ on $\mathbb{G}_m^n$ with weights $a$ and with weights $b$ are faithful, and that some $\K$-automorphism of $\K(x_1, \ldots, x_n)$ intertwines the scalings $x_i \efun \chi(a_i) x_i$ and $x_i \efun \chi(b_i) x_i$ for every $\chi$. Then $b = M a$ for some $M \in \operatorname{GL}_n(\Z)$. This consequence of Reichstein--Youssin, \emph{A birational invariant for algebraic group actions}, Theorem 7.1 and Corollary 2.5, is assumed for $\K$ a countable algebraically closed subfield of the residue field and $A = (\Delta + \sum_i \Z \alpha_i + \sum_i \Z \beta_i)/\Delta$ as in the proof after Notation 3.17.

\emph{Formalization.} Stated in Lean and taken as an explicit hypothesis, never proved: the author-approved cited input of 2026-11-12 (ACCEPTED-CITED-RESULTS.md, FINAL-CHECK.md). The paper cites it only in the unnumbered proof, after Notation 3.17, that the horizontal maps of the coefficient square are monomorphisms; in Lean it is also used for Remark 3.15, but not for Proposition 3.18 or Theorem 3.22.
\end{assumption}

\setcounter{theorem}{14}
\begin{remark}[Piecewise $\operatorname{GL}_k(\Z)$ characterization of $\Gamma$-morphisms]
\label{GLZ-char}
\lean{DiamondIntegral.Section3.gammaGroundContext_contraction_piecewiseIntegralAffine, DiamondIntegral.Section3.gammaGroundContext_GammaHom_iff_piecewiseGLZ, DiamondIntegral.Section3.gammaGround_code_contraction_piecewiseIntegralAffine, DiamondIntegral.Section3.gammaGroundGammaHom_iff_piecewiseGLZ}
\leanok
\uses{def-Ga-cat, defn-corr-cont}
If a definable function between two sets in $\Gamma$ is a $\vrv$-contraction, then it is piecewise $\Z$-affine, with constant terms of the form $\vrv(t)$ where $t \in \RV$ is definable, that is, $t \in \RV(\mathbb{S})$. Moreover, a definable bijection between two objects of $\Gamma[k]$ is a $\Gamma[k]$-morphism if and only if it is definably a piecewise $\operatorname{GL}_k(\Z)$-transformation.

\emph{Formalization.} Lean proves both statements locally, using neither the cited results of Yin nor Hrushovski--Kazhdan, but takes the Reichstein--Youssin consequence as an explicit hypothesis (author-approved input of 2026-11-12, ACCEPTED-CITED-RESULTS.md); the first sentence and the ``only if'' direction depend on it (FINAL-CHECK.md).
\end{remark}

\begin{proof}
\uses{hyp-reichstein-youssin}
\leanok
The first sentence is cited from Yin, \emph{Integration in algebraically closed valued fields with sections}, Remark 2.28. In the second, the ``if'' direction is clear; for the ``only if'' direction the paper cites Hrushovski--Kazhdan, \emph{Integration in valued fields}, Lemma 10.1, or Yin, \emph{Integration in algebraically closed valued fields with sections}, Lemma 2.29.
\end{proof}

\setcounter{theorem}{15}
\begin{definition}[$\Gamma$-categories with volume forms]
\label{def-Ga-cat-mu}
\lean{DiamondIntegral.Section3.GammaGroundMeasuredGammaObject, DiamondIntegral.Section3.GammaGroundMeasuredGammaHom, DiamondIntegral.Section3.GammaGroundMeasuredFiniteGammaObject}
\leanok
\uses{def-Ga-cat, jcbG-op}
An object of $\mG[k]$ is a definable pair $(I, \omega)$, where $I \in \Gamma[k]$ and $\omega : I \fun \Gamma$ is a function. A $\mG[k]$-morphism $(I, \omega) \fun (J, \sigma)$ is a $\Gamma[k]$-morphism $F : I \fun J$ such that $\omega(\alpha) = \sigma(F(\alpha)) + \jcb_{\Gamma} F(\alpha)$ for all $\alpha \in I$. The category $\mG^{\mathrm{fin}}[k]$ is the full subcategory of $\mG[k]$ of the objects with $I$ finite.
\end{definition}

\setcounter{theorem}{16}
\begin{definition}[Combined weights]
\label{comp-form}
\lean{DiamondIntegral.Section3.GammaGroundMeasuredRVObject.combinedWeight, DiamondIntegral.Section3.GammaGroundMeasuredGammaObject.combinedWeight}
\leanok
\uses{defn-RV-mu, def-Ga-cat-mu, jcbG-op}
For $(U, f, \omega) \in \mRV[k]$, write $\omega_f : U \fun \Gamma$ for the function $u \efun \Sigma(\vrv \circ f)(u) + \omega(u)$, where $\Sigma$ is the sum of the coordinates. Similarly, for $(I, \sigma) \in \mG[k]$, write $\sigma_I : I \fun \Gamma$ for the function $\gamma \efun \Sigma \gamma + \sigma(\gamma)$.
\end{definition}

\setcounter{theorem}{17}
\begin{proposition}[Positive tensor presentation]
\label{red-D-iso}
\lean{DiamondIntegral.Section3.proposition_3_18_unmeasured_kernel, DiamondIntegral.Section3.proposition_3_18_measured_kernel, DiamondIntegral.Section3.proposition_3_18_unmeasured, DiamondIntegral.Section3.proposition_3_18_measured}
\leanok
\uses{def-Kgroup, defn-RES-cat, def-Ga-cat, defn-RV-mu, def-Ga-cat-mu, 0coor, gamwhat}
Let $\Psi : T_+ \fun \gsk \RV[*]$, where $T_+ = \gsk \RES[*] \otimes_{\gsk \Gamma^{\mathrm{fin}}[*]} \gsk \Gamma[*]$, be induced by $[(U, f)] \otimes [I] \efun [(U \times I^\sharp, f \times \operatorname{id})]$, and $\mu\Psi : T_{\mu,+} \fun \gsk \mRV[*]$ its measured analogue on $T_{\mu,+} = \gsk \mRES[*] \otimes_{\gsk \mG^{\mathrm{fin}}[*]} \gsk \mG[*]$. For a $\RES$ object with underlying set $X$, let $\Lambda_X$ be the set of $\delta \in \Gamma \cap \operatorname{dcl}^{\mathrm{eq}}(\mathbb{S})$ admitting a definable map $X \fun \delta^\sharp$. Let $\mathcal{C}$ (respectively $\mathcal{C}_\mu$) be the semiring congruence generated by the pairs $[\boldsymbol{X}]_r \otimes [U]_m \sim [\boldsymbol{X}]_r \otimes [U + \delta e_i]_m$ (respectively $[\boldsymbol{X}]_r \otimes [(U, \sigma)]_m \sim [\boldsymbol{X}]_r \otimes [(U + \delta e_i, \sigma_\delta)]_m$ with $\sigma_\delta(\eta) = \sigma(\eta - \delta e_i) - \delta$), where $m \geq 1$, $U \in \Gamma[m]$ and $\delta \in \Lambda_X$. Then $\Psi$ and $\mu\Psi$ are surjective with kernel congruences $\mathcal{C}$ and $\mathcal{C}_\mu$, and induce graded semiring isomorphisms $T_+/\mathcal{C} \simeq \gsk \RV[*]$ and $T_{\mu,+}/\mathcal{C}_\mu \simeq \gsk \mRV[*]$.
\end{proposition}

\begin{proof}
\uses{mor-ftf, comp-form, def-Ga-cat}
\leanok
Surjectivity: the finite product normal form splits every $\RV[k]$ object into pieces $X \times U^\sharp$ with $X$ a $\RES$ object and $U \in \Gamma[m]$, after moving each constant measured weight onto the $\Gamma$ factor. The generating pairs lie in the kernels, since multiplying a free $\RV$ coordinate by a witness $c : X \fun \delta^\sharp$ is a bijection of $\Gamma$-Jacobian $\delta$. Conversely, by induction on the largest $\Gamma$ arity, an angular normal form $F(x, u) = (h(x), c(x) u^M)$ with $M \in \operatorname{GL}_m(\Z)$ on finitely many matched product cells is handled by $\Gamma$-morphisms and elementary pairs, and the residual pieces lie over integral hyperplanes, where one $\Gamma$ coordinate is moved into $\RES$; the measured case runs the same induction with the combined weights, without additive cancellation.
\end{proof}

\setcounter{theorem}{18}
\begin{definition}[Dropping the zero volume form]
\label{drop-0}
\lean{DiamondIntegral.Section3.gammaGroundOneForm}
\leanok
\uses{defn-c-RV-cat, defn-RV-mu}
The constant volume form $0$ is often dropped from the notation: if $\boldsymbol{U}$ is an object of $\RV[*]$, it may also denote the object $(\boldsymbol{U}, 0)$ of $\mRV[*]$ with the constant volume form $0$.

\emph{Formalization.} In Lean $\Gamma$ is written multiplicatively, so the constant form $0$ is the constant form $1$; this is notation only.
\end{definition}

\setcounter{theorem}{19}
\begin{definition}[The special congruences $\isp$ and $\mu\isp$]
\label{s3-special-congruence}
\lean{DiamondIntegral.Section3.gammaGroundRVSpecialCongruence, DiamondIntegral.Section3.gammaGroundMeasuredRVSpecialCongruence, DiamondIntegral.Section3.gammaGroundRVSpecialCompletionEquiv, DiamondIntegral.Section3.gammaGroundMeasuredRVSpecialCompletionEquiv}
\leanok
\uses{def-Kgroup, defn-c-RV-cat, defn-RES-cat, defn-RV-mu, 0coor, drop-0, defn-disc}
Let $[1] \in \gsk \RES[1]$ be the class of the singleton $\{1\}$; the class of $\{1\}$ in $\gsk \RES[0]$ is the identity $1$ of $\gsk \RES[*]$. Then $[\RV^{\circ\circ}_\infty] = [\RV^{\circ\circ}] + 1$ in $\gsk \RV[{\leq}1]$. Let $\isp$ be the (nonhomogeneous) semiring congruence on $\gsk \RV[*]$ generated by the pair $([1], [\RV^{\circ\circ}_\infty])$, and let $\boldsymbol{P} = [1] - [\RV^{\circ\circ}] \in \ggk \RV[1]$; the corresponding principal ideal of $\ggk \RV[*]$ is generated by $\boldsymbol{P} - 1$. Similarly, $\mu\isp$ is the homogeneous semiring congruence on $\gsk \mRV[*]$ generated by $([1], [\RV^{\circ\circ}])$, and the corresponding principal ideal of $\ggk \mRV[*]$ is generated by $\boldsymbol{P}$.
\end{definition}

\setcounter{theorem}{20}
\begin{definition}[The lifting maps $\mathbb{L}$ and $\mL$]
\label{RVcat-tag}
\lean{DiamondIntegral.Section3.gammaGroundObjectLift, DiamondIntegral.Section3.gammaGroundObjectLiftFinite, DiamondIntegral.Section3.gammaGroundMeasuredObjectLift}
\leanok
\uses{defn-c-RV-cat, 0coor, defn-VF-cat, defn-RV-mu, defn-VF-mu-G, gamwhat}
For each $\boldsymbol{U} = (U, f) \in \RV[k]$, let $U_f = \bigcup \{ f(u)^\sharp \times u : u \in U \}$. Let $\mathbb{L}_{\leq k} : \RV[{\leq}k] \fun \VF[k]$ be the map $\boldsymbol{U} \efun U_f$, and set $\mathbb{L} = \bigcup_k \mathbb{L}_{\leq k}$. Let $\mL_k : \mRV[k] \fun \mVF[k]$ be the map $(\boldsymbol{U}, \omega) \efun (\mathbb{L} \boldsymbol{U}, \mathbb{L} \omega)$, where $\mathbb{L} \omega$ is the function on $\mathbb{L} \boldsymbol{U}$ induced by $\omega$, and set $\mL = \bigoplus_k \mL_k$.
\end{definition}

\setcounter{theorem}{21}
\begin{theorem}[The Hrushovski--Kazhdan integral isomorphisms]
\label{main-prop-HK}
\lean{DiamondIntegral.Section3.theorem_3_22_bounded, DiamondIntegral.Section3.gammaGroundUnmeasuredRingIntegral, DiamondIntegral.Section3.gammaGroundDegreeMeasuredIntegral, DiamondIntegral.Section3.theorem_3_22_measured_graded}
\leanok
\uses{def-Kgroup, defn-VF-cat, defn-VF-mu-G, mor-equi, defn-c-RV-cat, defn-RV-mu, s3-special-congruence, RVcat-tag}
Suppose that $\mathbb{S}$ is $(\VF, \Gamma)$-generated. For each $k \geq 0$ there is a canonical semigroup isomorphism $\int_{+} : \gsk \VF[k] \fun \gsk \RV[{\leq}k] / \isp$ with $\int_{+} [A] = [\boldsymbol{U}] / \isp$ if and only if $[A] = [\mathbb{L} \boldsymbol{U}]$; these are compatible with the inductive systems, and the colimit of their groupifications is a canonical ring isomorphism $\int : \ggk \VF_* \fun \ggk \RV[*] / (\boldsymbol{P} - 1)$. Similarly, for each $k \geq 0$ there is a canonical semigroup isomorphism $\int_{+}^{\mu} : \gsk \mVF[k] \fun \gsk \mRV[k] / \mu\isp$ with $\int_{+}^{\mu} [\boldsymbol{A}] = [\boldsymbol{U}] / \mu\isp$ if and only if $[\boldsymbol{A}] = [\mL \boldsymbol{U}]$, and the direct sum of the groupifications is a canonical graded ring isomorphism $\int^{\mu} : \ggk \mVF[*] \fun \ggk \mRV[*] / (\boldsymbol{P})$.

\emph{Formalization.} Lean derives the theorem from the formalized \emph{Integration in algebraically closed valued fields} rather than from Hrushovski--Kazhdan; the generic integral constructions take lift-and-kernel records as parameters, and all these records are produced (section 3/GAP-REPAIRS.md). The ring isomorphism $\int$ is built directly, and its description as a colimit of the $\int_{+}$ is not a Lean statement (author-approved triage of 2026-11-18, FINAL-CHECK.md).
\end{theorem}

\begin{proof}
\uses{RVcat-tag, s3-special-congruence, defn-disc}
\leanok
The paper presents it as a combination of Hrushovski--Kazhdan, \emph{Integration in valued fields}, Theorems 8.8 and 8.29. For $(\VF, \Gamma)$-generated $\mathbb{S}$, $\mathbb{L}$ and $\mL$ induce surjective semiring homomorphisms onto $\gsk \VF_*$ and $\gsk \mVF[*]$ (Hrushovski--Kazhdan, Sections 4 and 6, or Yin, \emph{Special transformations in algebraically closed valued fields}, Corollaries 7.7 and 10.6); they kill $\boldsymbol{P} - 1$ and $\boldsymbol{P}$ respectively, since $\mathbb{L}([1]) = [\mathcal{M}]$ and $\mathbb{L}([\RV^{\circ\circ}]) = [\mathcal{M} \mi 0]$, and these relations generate the kernels.
\end{proof}

\section{Uniform retraction to $\RES$}
\label{section-4}
Euler characteristics on $\Gamma$, the retractions $\mathbb{E}_g$ and $\mathbb{E}_b$ to the residue rings, and the relative annular ideal $J_{\mathrm{var}}$ with its compatibility lemma, which Section 7 uses for the kernel of the bounded integral.

\setcounter{theorem}{0}
\begin{definition}[Dimension-free $\RES$-categories]
\label{df-res}
\lean{DiamondIntegral.Section4.RESDimensionFreeCategory.Object, DiamondIntegral.Section4.RESDimensionFreeCategory.Iso, DiamondIntegral.Section4.RESDimensionFreeCategory.Ring, DiamondIntegral.Section4.RESDimensionFreeCategory.definition_4_1_identification}
\leanok
\uses{defn-RES-cat}
The objects of the category $\RES$ are obtained from the objects $(U, f)$ of $\RES[*]$ by forgetting the function $f$. Any definable bijection between two such objects is a morphism of $\RES$.
\end{definition}

\setcounter{theorem}{1}
\begin{remark}[Geometric and bounded Euler characteristics]
\label{s4-euler-characteristics}
\lean{DiamondIntegral.Section4.GammaContinuousDecomposition.euler_characteristic_discharged, DiamondIntegral.Section4.GammaBijectiveRelations.integral_bijective_relation, DiamondIntegral.Section4.GammaBoxStabilization.box_stabilization, DiamondIntegral.Section4.EulerFunction.totalEuler_ray}
\leanok
\uses{def-Ga-cat, bounded}
Since the $\Gamma$-sort is an o-minimal group, there is the geometric o-minimal Euler characteristic $\chi_g : \ggk \Gamma[*] \to \Z$: for a definable $X \sub \Gamma^k$ choose a finite cell decomposition $X = \coprod_j C_j$ (cell dimensions may vary) and set $\chi_g(X) = \sum_j (-1)^{\dim C_j}$. This does not depend on the decomposition and is invariant under definable bijections. Moreover $\chi_g(X \cap [-\gamma, \gamma]^k)$ stabilizes as $\gamma \to +\infty$, and the induced map is the bounded Euler characteristic $\chi_b : \ggk \Gamma[*] \to \Z$. The two agree on doubly bounded sets, whereas $\chi_g((0, +\infty)) = -1$ and $\chi_b((0, +\infty)) = 0$. Both are written $\chi$ when no distinction is needed.

\emph{Formalization.} Bijection invariance and stabilization are stated with the unary $\Gamma$ rational-affine normal form as an argument; that input is proved (\texttt{Section5.GammaQE.gamma\_piecewise\_rational\_affine}), and the cell decomposition with its Euler sum has an input-free form (section 4/COVERAGE.md).
\end{remark}

\setcounter{theorem}{2}
\begin{definition}[Affine line and reduced residue rings]
\label{res-note}
\lean{DiamondIntegral.Section4.RESOrdinaryReduction.ReducedRing, DiamondIntegral.Section4.RESDimensionFreeReduction.Ring, DiamondIntegral.Section4.RESMeasuredReduction.fibreIdeal, DiamondIntegral.Section4.RESMeasuredReduction.affineClass}
\leanok
\uses{defn-RES-cat, defn-RV-mu}
Let $[\A] \in \ggk \RES$ be the class of the affine line over the residue field and $[\G_m] = [\A] - 1$ that of the multiplicative torus. In the graded ring $\ggk \RES[*]$, $[\A]$ denotes the homogeneous degree-one class $[\G_m] + [1]$, where $[1]$ is the degree-one point class and the identity is the degree-zero point class $1$; graded localizations by $[\A]$ use this class. Let $!I$ be the ideal of $\ggk \RES$ generated by the $[\gamma^\sharp] - [\G_m]$ with $\gamma \in \Gamma$ definable, and $!\ggk \RES = \ggk \RES / !I$; in the graded ring, $!I[*]$ is generated by the same differences in degree one and $!\ggk \RES[*] = \ggk \RES[*] / !I[*]$. Let $!\mu I[*]$ be the ideal of $\ggk \mRES[*]$ generated by the $[(\G_m, \alpha + \beta)] - [(\alpha^\sharp, \beta)] \in \ggk \mRES[1]$ with $\alpha, \beta \in \Gamma$ definable, and $!\ggk \mRES[*] = \ggk \mRES[*] / !\mu I[*]$.

\emph{Formalization.} Lean forms the measured quotient together with $([\A])$ (\texttt{RESMeasuredReduction.ReducedRing}), which is the target ring of Proposition 4.7.
\end{definition}

\setcounter{theorem}{3}
\begin{remark}[Localized targets and distinct retractions]
\label{eb-no-eg}
\lean{DiamondIntegral.Section4.RESAffineNonsurjective.remark_4_4, DiamondIntegral.Section4.RESPointLocalization.dimensionFreeEquiv, DiamondIntegral.Section4.RESAffineNormalization.comparison_point, DiamondIntegral.Section4.Section4LinkedRemarks.normalized_retractions_distinct}
\leanok
\uses{res-note, s4-euler-characteristics, red-D-iso, s3-special-congruence}
Let $\mathbb{E}_g : \ggk \RV[*] / (\mathbf{P} - 1) \to !\ggk \RES[*][[\A]^{-1}]$ and $\mathbb{E}_b : \ggk \RV[*] / (\mathbf{P} - 1) \to !\ggk \RES[*][[1]^{-1}]$ be the retractions of Section 4.1, induced by $x \otimes y \mapsto \chi(y) x [\G_m]^l$ for $y \in \ggk \Gamma[l]$ with $\chi = \chi_g$, resp. $\chi_b$. The piece $(!\ggk \RES[*][[\A]^{-1}])_0$ is canonically the colimit of the $!\ggk \RES[n]$ under $x \mapsto x[\A]^{m-n}$, and the natural homomorphism $(!\ggk \RES[*][[\A]^{-1}])_0 \to !\ggk \RES[[\A]^{-1}]$, $x \mapsto x[\A]^{-n}$ for $x \in !\ggk \RES[n]$, is not surjective. The analogous map $(!\ggk \RES[*][[1]^{-1}])_0 \to !\ggk \RES[[1]^{-1}] \cong !\ggk \RES$ is an isomorphism. With codomains in $!\ggk \RES[[\A]^{-1}]$, $\mathbb{E}_b \neq \mathbb{E}_g$, since $\mathbb{E}_b([1]) = 1$ and $\mathbb{E}_g([1]) = -[\G_m][\A]^{-1} + 1 = [\A]^{-1}$. Forcing $[\A] = 1$ equalizes them and gives the homomorphism of Hrushovski and Loeser.

\emph{Formalization.} The colimit description and the equalization by $[\A] = 1$ (with the comparison to Hrushovski and Loeser) are not formalized, by the author-approved triage of 2026-11-18 (section 4/GAP-REPAIRS.md and COVERAGE.md). The non-surjectivity, the isomorphism at $[1]$ and the distinctness are formalized.
\end{remark}

\setcounter{theorem}{4}
\begin{definition}[Generators of $\operatorname{FN}(\Gamma, \Z)$]
\label{gam-equ-cons}
\lean{DiamondIntegral.Section4.GammaFunctionRing.functionCommRing, DiamondIntegral.Section4.EulerFunction.point_ray, DiamondIntegral.Section4.EulerFunction.ray_mul, DiamondIntegral.Section4.EulerFunction.line_ray}
\leanok
\uses{s4-euler-characteristics}
Let $\operatorname{FN}(\Gamma, \Z)$ be the commutative ring of functions $f : \Gamma \to \Z$ with finite range and definable level sets, with pointwise addition and the convolution product $(f * g)(\gamma) = \sum_m m\, \chi(h_\gamma^{-1}(m))$, where $h_\gamma(\alpha) = f(\alpha) g(\gamma - \alpha)$. Let $r$ be the constant function $1$ and, for definable $\gamma \in \Gamma$, let $p_\gamma$, $q_\gamma$ be the characteristic functions of $\{\gamma\}$ and $(\gamma, \infty)$; $p_0$ is the identity. Then $\operatorname{FN}(\Gamma, \Z)$ is generated as a $\Z$-module by the $r, p_\gamma, q_\gamma$, and \[ r p_\beta = r, \quad p_\alpha p_\beta = p_{\alpha + \beta}, \quad p_\alpha q_\beta = q_{\alpha + \beta}, \quad q_\alpha q_\beta = -q_{\alpha + \beta}. \] Moreover $r^2 = r q_\beta = -r$ if $\chi = \chi_g$, and $r^2 = r$, $r q_\beta = 0$ if $\chi = \chi_b$.
\end{definition}

\setcounter{theorem}{5}
\begin{remark}[Euler profiles of measured $\Gamma$-objects]
\label{int-Gam}
\lean{DiamondIntegral.Section4.GammaRVChartPresentation.interpretation, DiamondIntegral.Section4.GammaMeasuredProfileInvariance.profile_volume_bijOn, DiamondIntegral.Section4.GammaProfileProducts.profile_product, DiamondIntegral.Section4.GammaPolynomialGradings.objectPolynomial_mem}
\leanok
\uses{def-Ga-cat-mu, gam-equ-cons, s4-euler-characteristics}
For $\mathbf{I} = (I, \mu) \in \mG[k]$ let $\mu_I(\gamma) = \mu(\gamma) + \Sigma\gamma$ on $I$ and $\lambda_{\mathbf{I}}(\gamma) = \chi(\mu_I^{-1}(\gamma))$. Then $\lambda_{\mathbf{I}} \in \operatorname{FN}(\Gamma, \Z)$ depends only on the class $[\mathbf{I}]$, is additive in each degree, and $\lambda_{[\mathbf{I}]} * \lambda_{[\mathbf{J}]} = \lambda_{[\mathbf{I}][\mathbf{J}]}$. Put $\Gamma_{\mathrm{def}} = \Gamma \cap \operatorname{dcl}^{\mathrm{eq}}(\mathbb{S})$ and let $\Z[\Gamma_{\mathrm{def}}]$ be the subring of $\operatorname{FN}(\Gamma, \Z)$ spanned by the $p_\gamma$, $\gamma \in \Gamma_{\mathrm{def}}$. Then $\ggk \mG[0] \cong \Z[\Gamma_{\mathrm{def}}]$, and $[\mathbf{I}] \mapsto \lambda_{[\mathbf{I}]} X^k$ for $\mathbf{I} \in \mG[k]$ defines a graded ring homomorphism \[ \lambda : \ggk \mG[*] \to \Z[\Gamma_{\mathrm{def}}] \oplus X \operatorname{FN}(\Gamma, \Z)[X]. \] There are two such, $\lambda_g$ and $\lambda_b$, for $\chi = \chi_g$ and $\chi = \chi_b$.

\emph{Formalization.} The isomorphism $\ggk \mG[0] \cong \Z[\Gamma_{\mathrm{def}}]$ is not formalized, by the author-approved triage of 2026-11-18 (section 4/GAP-REPAIRS.md and COVERAGE.md); the homomorphism $\lambda$, its invariance, multiplicativity and grading are.
\end{remark}

\begin{proof}
\uses{s4-euler-characteristics}
O-minimality gives finite range and definable level sets. A $\mG[k]$-morphism restricts to bijections between the fibres of the corrected weights, so $\lambda$ is invariant; a direct computation identifies $\lambda_{[\mathbf{I}]} * \lambda_{[\mathbf{J}]}$ with the Euler characteristics of the fibres of $(\alpha, \beta) \mapsto \mu(\alpha) + \sigma(\beta)$ on $I \times J$.
\end{proof}

\setcounter{theorem}{6}
\begin{proposition}[Measured Euler retractions]
\label{prop-eu-retr-k}
\lean{DiamondIntegral.Section4.Section3LinkedRetractionStatements.proposition_4_7_discharged, DiamondIntegral.Section4.Section3LinkedRetractionStatements.proposition_4_7}
\leanok
\uses{res-note, defn-RV-mu, s3-special-congruence}
Let $\mathbf{P} = [1] - [\RV^{\circ\circ}] \in \ggk \mRV[1]$ and $\mathfrak{h} = [(\RV^{\circ\circ}, -\vrv)] \in \ggk \mRV[1]$. There are two graded ring homomorphisms \[ \mgE_g, \mgE_b : \ggk \mRV[*] \to !\ggk \mRES[*] / ([\A]) \] such that \begin{itemize} \item $\mathbf{P}$ vanishes under both of them, \item $\mgE_g(x) = \mgE_b(x) = x / ([\A])$ for all $x \in \ggk \mRES[*]$, \item they are distinguished by $\mgE_g(\mathfrak{h}) = +[1] / ([\A])$ and $\mgE_b(\mathfrak{h}) = 0$. \end{itemize}

\emph{Formalization.} \texttt{proposition\_4\_7} takes the unary $\Gamma$ rational-affine normal form as an argument; it is proved (\texttt{Section5.GammaQE.gamma\_piecewise\_rational\_affine}) and supplied in \texttt{proposition\_4\_7\_discharged} (section 4/GAP-REPAIRS.md). The Lean statement also records $\mgE_g \neq \mgE_b$.
\end{proposition}

\begin{proof}
\uses{red-D-iso, int-Gam, gam-equ-cons, res-note, s4-euler-characteristics}
\leanok
The product map $(\iota / ([\A])) \times (\psi \circ \lambda)$, with $\lambda$ from Remark 4.6 and $\psi$ the graded homomorphism given by $p_\gamma X^k \mapsto [\G_m]^{k-1}[(\G_m, \gamma)]$, $q_\gamma X^k \mapsto [\G_m]^{k-1}[(\{1\}, \gamma)]$, $r X^k \mapsto 0$, is bilinear over $\ggk \mG^{\mathrm{fin}}[*]$ because of $!\mu I[*]$, and identifies the measured coefficient-translation pairs because translation preserves corrected weights; so it descends through the kernel congruence of $\mu\Psi$ (Proposition 3.18). The values on $\mathbf{P}$ and $\mathfrak{h}$ are computed on explicit tensor preimages, using $\lambda_g([((0, \infty), -\mathrm{id})]) = -p_0 X$ and $\chi_b((0, \infty)) = 0$. Finally $[1] \neq 0$ in the target: map to pure residue-field definable sets modulo $\mathbb{L}$, specialize by the Lefschetz principle to the algebraic closure of a finite field, and count $\mathbb{F}_q$-points modulo $q$.
\end{proof}

\setcounter{theorem}{7}
\begin{remark}[Alternative assignment for $\chi_g$]
\label{s4-alternative-psi}
\lean{DiamondIntegral.Section4.Section4LinkedRemarks.alternativeEvaluation, DiamondIntegral.Section4.Section4LinkedRemarks.alternativeEvaluation_res, DiamondIntegral.Section4.RESAlternativePsi.gradedPsi, DiamondIntegral.Section4.RESAlternativeForget.forget}
\leanok
\uses{res-note, gam-equ-cons, int-Gam}
Let $\psi : \Z[\Gamma_{\mathrm{def}}] \oplus X \operatorname{FN}(\Gamma, \Z)[X] \to !\ggk \mRES[*] / ([\A])$ be the graded homomorphism given for $k \geq 1$ by $p_\gamma X^k \mapsto [\G_m]^{k-1}[(\G_m, \gamma)]$, $q_\gamma X^k \mapsto [\G_m]^{k-1}[(\{1\}, \gamma)]$, $r X^k \mapsto 0$. For $\chi = \chi_g$ one may instead send $r X^k \mapsto [\G_m]^{k-1}[1]$, the same element as $q_0 X^k$. Let $J$ be the enlargement of $!\mu I[*]$ by the elements $[1] - [(\{1\}, \gamma)] \in \ggk \mRES[1]$. Then $\psi$ is still a graded ring homomorphism into $\ggk \mRES[*] / (J + ([\A]))$, and one obtains a ring homomorphism \[ \ggk \mRV[*] \to \ggk \mRES[*] / (J + ([\A])) \xrightarrow{\ \phi\ } !\ggk \RES[*] / ([\A]), \] where $\phi$ is induced by the forgetful functor.
\end{remark}

\setcounter{theorem}{8}
\begin{definition}[Doubly bounded $\RV$-categories]
\label{db-RV}
\lean{DiamondIntegral.Section4.RVMeasuredBounds.DomainBounded, DiamondIntegral.Section4.RVBoundedRing.Ring, DiamondIntegral.Section4.GammaBoundedRVPresentation.ChartRing, DiamondIntegral.Section4.RVStrictBoundedComparison.equivalence}
\leanok
\uses{defn-c-RV-cat, defn-RV-mu, def-Ga-cat-mu, bounded}
An object $(U, f) \in \RV[k]$ is doubly bounded if $U$ is. The full subcategory of $\RV[k]$ of doubly bounded objects is denoted by $\RV^{\db}[k]$, and similarly for the categories $\mRV^{\db}[k]$ and $\mG^{\db}[k]$.

\emph{Formalization.} Lean's bound leaves the $\infty$-coordinates of a chart unconstrained, so its ring for $\mRV^{\db}$ has more generators; the ring on $\infty$-free charts (the paper's objects) is proved isomorphic to it, compatibly with the grading, $(\mathbf{P}_\Gamma)$ and $J_{\mathrm{var}}$ (author decision of 2026-11-18, section 4/GAP-REPAIRS.md). The unmeasured $\RV^{\db}[k]$ has no Lean category of its own; its ring is among the items of Notation 4.13 left out by the triage of 2026-11-18.
\end{definition}

\setcounter{theorem}{9}
\begin{remark}[Bounded tensor product map]
\label{tensor-bd}
\lean{DiamondIntegral.Section4.RVBoundedDecomposition.coordinate_image_bounded, DiamondIntegral.Section4.RVBoundedDecomposition.volume_image_bounded, DiamondIntegral.Section4.RVBoundedPositiveTensor.comparison, DiamondIntegral.Section4.RVBoundedPositiveTensor.comparison_charts}
\leanok
\uses{db-RV, red-D-iso}
If $(U, f, \omega) \in \mRV^{\db}[k]$ then $f(U)$ and $\omega(U)$ are doubly bounded. The product construction gives a homomorphism \[ \mu\Psi^{\db} : \gsk \mRES[*] \otimes \gsk \mG^{\db}[*] \to \gsk \mRV^{\db}[*], \] the tensor product being over $\gsk \mG^{\mathrm{fin}}[*]$. It is used to compute images of explicit tensor classes; the bounded retraction itself is obtained by restricting the graded map $\widetilde{\mathbb{E}}_b$ after forgetting volume forms.
\end{remark}

\begin{proof}
\uses{db-to-db, red-D-iso}
\leanok
The first sentence is Lemma 2.41 applied to the coordinates of $f$ and to $\omega$; the map sends $[\mathbf{X}] \otimes [\mathbf{I}]$ to the class of the product of $\mathbf{X}$ with the full $\RV$-preimage of $\mathbf{I}$, as for $\mu\Psi$ in Proposition 3.18.
\end{proof}

\setcounter{theorem}{10}
\begin{definition}[Generators of $\operatorname{FN}^{\db}(\Gamma, \Z)$]
\label{db-generator}
\lean{DiamondIntegral.Section4.GammaBoundedFunctions.boundedBaseSubring, DiamondIntegral.Section4.EulerFunction.boundedInterval, DiamondIntegral.Section4.GammaBoundedFunctions.base_generator_expansion, DiamondIntegral.Section4.EulerFunction.point_mul_boundedInterval}
\leanok
\uses{gam-equ-cons, bounded}
Let $\operatorname{FN}^{\db}(\Gamma, \Z)$ be the subgroup of $\operatorname{FN}(\Gamma, \Z)$ of functions whose nonzero locus is doubly bounded. For definable $\gamma$ put $o_\gamma = \mathbf{1}_{(0, \gamma)}$ if $\gamma > 0$, $o_\gamma = \mathbf{1}_{(\gamma, 0)}$ if $\gamma < 0$, and $o_0 = -p_0$ (a convention, not the characteristic function of an empty interval). Then $\operatorname{FN}^{\db}(\Gamma, \Z)$ is generated as a $\Z$-module by the $p_\gamma, o_\gamma$; $o_\gamma = q_0 - q_\gamma - p_\gamma$ for $\gamma > 0$ and $o_\gamma = q_\gamma - q_0 - p_0$ for $\gamma < 0$. For $\beta \neq 0$, with $a = \min(\alpha, \alpha + \beta)$ and $b = \max(\alpha, \alpha + \beta)$, $p_\alpha o_\beta = \mathbf{1}_{(a, b)}$ equals $o_b - o_a - p_a$, $o_a - o_b - p_b$ or $o_a + o_b + p_0$ according as $0 \leq a < b$, $a < b \leq 0$ or $a < 0 < b$; $p_\alpha o_0 = -p_\alpha$; and $o_\alpha o_\beta = -o_{\alpha + \beta}$ if $\alpha, \beta \geq 0$ or $\alpha, \beta \leq 0$, while $o_\alpha o_\beta = -o_\alpha - o_\beta - p_0$ if $\alpha < 0 < \beta$ or $\beta < 0 < \alpha$.
\end{definition}

\setcounter{theorem}{11}
\begin{definition}[Bounded profile homomorphism]
\label{db-gam-fn}
\lean{DiamondIntegral.Section4.GammaBoundedChartMap.interpretation, DiamondIntegral.Section4.GammaBoundedChartMap.interpretation_boundedClass, DiamondIntegral.Section4.GammaBoundedProfileRange.boundedProfile_changeKind, DiamondIntegral.Section4.GammaBoundedComparison.baseBoundedEuler_changeKind}
\leanok
\uses{int-Gam, db-RV, db-generator, s4-euler-characteristics}
The homomorphism \[ \lambda : \ggk \mG^{\db}[*] \to \Z[\Gamma_{\mathrm{def}}] \oplus X \operatorname{FN}^{\db}(\Gamma, \Z)[X], \qquad [\mathbf{I}] \mapsto \lambda_{[\mathbf{I}]} X^k \ (\mathbf{I} \in \mG^{\db}[k]), \] is constructed as in Remark 4.6. There is only one such homomorphism, since $\chi_g$ and $\chi_b$ agree on doubly bounded sets.
\end{definition}

\setcounter{theorem}{12}
\begin{definition}[Radius generators $\mathbf{P}_\gamma$]
\label{pgammabis}
\lean{DiamondIntegral.Section4.RVRadiusIdeal.pGamma, DiamondIntegral.Section4.RVRadiusIdeal.radiusIdeal, DiamondIntegral.Section4.RVOrdinaryRadius.radiusIdeal_le_ker, DiamondIntegral.Section4.RVMeasuredRadius.radiusIdeal_le_comap_pIdeal}
\leanok
\uses{defn-disc, db-RV, def-Ga-cat, s3-special-congruence, s3-definable-value-hull}
Let $\RV^{\circ\circ}_\gamma = \rv(\mathcal{M}_\gamma \mi 0)$, and let $D_{\RV}$ consist of the nonnegative definable $\gamma \in \Gamma$ such that $\gamma^\sharp$ contains an $\mathbb{S}$-definable point. For $\gamma \in D_{\RV}$ let \[ \mathbf{P}_\gamma = [\RV^{\circ\circ} \mi \RV^{\circ\circ}_\gamma] + [\{t_\gamma\}] - [1] \in \ggk \RV^{\db}[1], \] where $t_\gamma \in \gamma^\sharp$ is a definable point; it also denotes the corresponding element of $\ggk \mRV^{\db}[1]$ with constant volume form $0$. Although $[\{t_\gamma\}] = [1]$ in $\ggk \RV^{\db}[1]$, $[\{t_\gamma\}] \neq [1]$ in $\ggk \mRV^{\db}[1]$ unless $\gamma = 0$. $\mathbf{P}_\gamma$ does not depend on $t_\gamma$. Let $(\mathbf{P}_\Gamma)$ be the ideal of $\ggk \mRV^{\db}[*]$ generated by the $\mathbf{P}_\gamma$, $\gamma \in D_{\RV}$. Its images under the forgetful homomorphisms to $\ggk \RV[*]$ and $\ggk \mRV[*]$ lie in $(\mathbf{P} - 1)$ and $(\mathbf{P})$.

\emph{Formalization.} The unmeasured ring $\ggk \RV^{\db}[*]$ and the inequality $[\{t_\gamma\}] \neq [1]$ in $\ggk \mRV^{\db}[1]$ are not formalized, by the author-approved triage of 2026-11-18 (section 4/GAP-REPAIRS.md and COVERAGE.md); $\mathbf{P}_\gamma$, its independence of $t_\gamma$, the ideal and the two forgetful containments are formalized in the measured ring (volume forgetting even sends each $\mathbf{P}_\gamma$ to $0$).
\end{definition}

\setcounter{theorem}{13}
\begin{remark}[Bounded $\Gamma$ assignment to $\RES$]
\label{db-gam-res}
\lean{DiamondIntegral.Section4.GammaBoundedQuotient.boundedBaseQuotientEquiv, DiamondIntegral.Section4.GammaBoundedQuotient.baseQuotientMap_eq_intEuler, DiamondIntegral.Section4.RESBoundedPsi.psi, DiamondIntegral.Section4.RESBoundedPsi.psi_object}
\leanok
\uses{db-generator, db-gam-fn, res-note, s4-euler-characteristics}
Let $!P$ be the ideal of $\operatorname{FN}^{\db}(\Gamma, \Z)$ generated by the elements $p_\gamma - p_0$ and $o_\gamma + p_0$. Then $\operatorname{FN}^{\db}(\Gamma, \Z) / !P \cong \Z$, and the quotient map is $\chi$: if $\mathbf{I} = (I, \mu) \in \mG^{\db}[*]$ then $\lambda_{\mathbf{I}} / !P = \chi(I)$. The assignments $p_\gamma X^k \mapsto [\G_m]^k$, $o_\gamma X^k \mapsto -[\G_m]^k$ induce a graded ring homomorphism \[ \psi^{\db} : \Z[\Gamma_{\mathrm{def}}] \oplus X \operatorname{FN}^{\db}(\Gamma, \Z)[X] \to !\ggk \RES[*], \] and $\psi^{\db} \circ \lambda$ is the forgetful map $[(I, \mu)] \mapsto \chi(I) [\G_m]^k$ for $(I, \mu) \in \mG^{\db}[k]$.
\end{remark}

\begin{proof}
\uses{db-generator}
\leanok
The quotient is computed from the product table of Notation 4.11, and the same table shows that the assignments respect all relations among the generators.
\end{proof}

\setcounter{theorem}{14}
\begin{proposition}[Bounded retraction to $\RES$]
\label{prop-eu-retr-k-db}
\lean{DiamondIntegral.Section4.Section3LinkedRetractionStatements.proposition_4_15_discharged, DiamondIntegral.Section4.Section3LinkedRetractionStatements.proposition_4_15}
\leanok
\uses{db-RV, pgammabis, res-note}
Let $\phi : \ggk \mRES[*] \to !\ggk \RES[*]$ be the forgetful homomorphism. There is a graded ring homomorphism \[ \mgE^{\db} : \ggk \mRV^{\db}[*] \to !\ggk \RES[*] \] such that $(\mathbf{P}_\Gamma)$ vanishes under it and $\mgE^{\db}(x) = \phi(x)$ for all $x \in \ggk \mRES[*]$.

\emph{Formalization.} Stated on Lean's ring for $\mRV^{\db}$, whose charts may have $\infty$-coordinates; it is isomorphic to the paper's ring compatibly with the grading and $(\mathbf{P}_\Gamma)$ (Definition 4.9, author decision of 2026-11-18). \texttt{proposition\_4\_15} takes the proved unary $\Gamma$ rational-affine normal form as an argument, which \texttt{proposition\_4\_15\_discharged} supplies.
\end{proposition}

\begin{proof}
\uses{tensor-bd, db-gam-res, pgammabis, red-D-iso, s4-euler-characteristics}
\leanok
Set $\mgE^{\db} = \widetilde{\mathbb{E}}_b \circ F$, where $F : \ggk \mRV^{\db}[*] \to \ggk \RV[*]$ forgets volume forms and $\widetilde{\mathbb{E}}_b : \ggk \RV[*] \to !\ggk \RES[*]$ is the graded map of Section 4.1, obtained by descending $x \otimes y \mapsto \chi_b(y) x [\G_m]^l$ through the kernel congruence of $\Psi$ (Proposition 3.18). On bounded tensors it is $\phi \times (\psi^{\db} \circ \lambda)$ by Remark 4.14. For the radius generators, $\mathbf{P}_0 = 0$, and for $\gamma > 0$ the preimage $1 \otimes [(0, \gamma)] + 1 \otimes [\{\gamma\}] + [\{t_\gamma\}] \otimes 1 - [1] \otimes 1$ of $\mathbf{P}_\gamma$ under $\mu\Psi^{\db}$ maps to $-[\G_m] + [\G_m] + [1] - [1] = 0$.
\end{proof}

\setcounter{theorem}{15}
\begin{remark}[$\mathbb{E}_g$ does not commute with $\mathbb{E}^\diamond$]
\label{eb-not-eg}
\lean{DiamondIntegral.Section4.Section4LinkedRemarks.diagram_edb_eb_com, DiamondIntegral.Section4.Section4LinkedRemarks.normalized_retractions_distinct}
\leanok
\uses{prop-eu-retr-k-db, eb-no-eg, db-gam-res}
Let $\mathbb{E}^\diamond$ be $\mgE^{\db}$ followed by the forgetful map $!\ggk \RES[*] \to !\ggk \RES$. By Remark 4.14 it equals $\mathbb{E}_b \circ F$, where $F : \ggk \mRV^{\db}[*] \to \ggk \RV[*]$ forgets volume forms and $\mathbb{E}_b$ is followed by the identification $(!\ggk \RES[*][[1]^{-1}])_0 \cong !\ggk \RES$ of Remark 4.4 (the commutative triangle (4.8)). After changing $!\ggk \RES$ to $!\ggk \RES[[\A]^{-1}]$ in (4.8), $\mathbb{E}_g$ does not commute with $\mathbb{E}^\diamond$: $\mathbb{E}_g([1]) = [\A]^{-1}$, whereas $\mathbb{E}^\diamond([1]) = \mathbb{E}_b([1]) = 1$. One could redefine $\mathbb{E}^\diamond$ through $!\ggk \RES[*] \to !\ggk \RES[[\A]^{-1}]$ to make it compatible with $\mathbb{E}_g$, but $\mathbb{E}_b$ is the better choice.

\emph{Formalization.} The second paragraph (the redefinition of $\mathbb{E}^\diamond$, which the paper rejects) is not formalized, by the author-approved triage of 2026-11-18 (section 4/GAP-REPAIRS.md and COVERAGE.md). The triangle (4.8) and the distinctness are formalized.
\end{remark}

\setcounter{theorem}{16}
\begin{definition}[Relative annular ideal]
\label{relative-annular-ideal}
\lean{DiamondIntegral.Section4.RelativeAnnularGenerators.AnnularDatum, DiamondIntegral.Section4.RelativeAnnularGenerators.AnnularDatum.generator, DiamondIntegral.Section4.RelativeAnnularGenerators.manuscriptIdeal, DiamondIntegral.Section4.RelativeAnnularScalar.radiusIdeal_le_manuscriptIdeal}
\leanok
\uses{db-RV, pgammabis, defn-RV-mu}
Put $B = \ggk \mRV^{\db}[*]$ and $I = (\mathbf{P}_\Gamma)$. For $\mathbf{U} = (U, f, \omega) \in \mRV^{\db}[k]$ with $k > 0$ and a coordinate $j$ such that $\operatorname{pr}_{\widetilde{j}} \rest f(U)$ is finite-to-one, let $\tau : U \to \RV$ be definable with $0 \leq \gamma(u) := \vrv(\tau(u)) \leq \Delta$ for one finite named $\Delta$. The relative annular object is the separately tagged union \[ U^{\tau, j} = \{(u, s) : 0 < \vrv(s) \leq \gamma(u)\} \uplus \{(u, \tau(u)) : u \in U\}, \] with structural coordinates $s f_j(u)$ in position $j$ and $f_i(u)$ in every other position, and volume form $\omega(u)$. Let $J_{\mathrm{var}}$ be the ideal of $B$ generated by the $[\mathbf{U}^{\tau, j}] - [\mathbf{U}]$ for all these data; the full aperture graph is retained, even when its image is infinite. Constant apertures give $I \sub J_{\mathrm{var}}$; no equality is asserted.

\emph{Formalization.} Defined in Lean's ring for $\mRV^{\db}$, whose charts may have $\infty$-coordinates (Definition 4.9); under the isomorphism with the paper's ring, the ideal generated by the data with $\infty$-free source corresponds exactly to it (\texttt{RVStrictAnnularGenerators.map\_strictAnnularIdeal}).
\end{definition}

\setcounter{theorem}{17}
\begin{lemma}[Compatibility with the relative annular ideal]
\label{relative-annular-compatibility}
\lean{DiamondIntegral.Section4.RelativeAnnularGenerators.manuscriptIdeal_le_forget_kernel, DiamondIntegral.Section4.RelativeAnnularGenerators.manuscriptIdeal_homogeneous, DiamondIntegral.Section4.RelativeAnnularGenerators.manuscript_quotient_injective_on_degree_zero, DiamondIntegral.Section4.RelativeAnnularGenerators.manuscriptIdeal_maps_into_p}
\leanok
\uses{relative-annular-ideal, prop-eu-retr-k-db, s3-special-congruence}
Let $B = \ggk \mRV^{\db}[*]$ and $J_{\mathrm{var}}$ be as in Notation 4.17. The volume-forgetting map $F : B \to \ggk \RV[*]$ kills $J_{\mathrm{var}}$. Hence the bounded retractions of this section factor through $B / J_{\mathrm{var}}$, with the same restrictions to the residue rings. The ideal $J_{\mathrm{var}}$ is homogeneous and meets the entire degree-zero subring trivially. Its image in the full measured ring $\ggk \mRV[*]$ lies in $(\mathbf{P})$.

\emph{Formalization.} Stated for Lean's ring for $\mRV^{\db}$ ($\infty$ convention of Definition 4.9). The factorization clause is \texttt{Section4LinkedRemarks.manuscript\_section4\_compatibility\_discharged}.
\end{lemma}

\begin{proof}
\uses{relative-annular-ideal, prop-eu-retr-k-db}
\leanok
With $\mathbf{W} = (U, f_{\widetilde{j}}, \omega + \vrv(f_j))$ and $\mathbf{W}_\gamma$ its shift by $\gamma$, the source and the endpoint are $\mathbf{W}$ and $\mathbf{W}_\gamma$ times the structural unit point. After forgetting volume, $(u, s) \mapsto (u, \tau(u) s)$ identifies the positive-ray bundle with its deep subray, whose complement is the shell, so the shell vanishes by scissors and the endpoint is isomorphic to the source; thus $F$ kills every generator. Generators have positive homogeneous degree, which gives homogeneity and the degree-zero claim, and in the measured ring the same shear gives $\operatorname{image}([\mathbf{U}^{\tau, j}] - [\mathbf{U}]) = ([\mathbf{W}_\gamma] - [\mathbf{W}]) \mathbf{P}$. The factorization follows by composing $F$ with the Euler maps, as in Proposition 4.15.
\end{proof}

\section{Interlude: RV decomposition and the realization in varieties}
\label{section-5}
The three items of this section are recalled from Forey and Yin, \emph{Fubini theorems for motivic volume}, \href{https://github.com/daviddhilbert/Fubini-theorems}{Fubini-theorems repository}, file \texttt{paper/fubini-theorems.tex}, and are proved in the pinned Fubini-theorems package. The Lean names keep the numbers 5.12, 5.14 and 5.15 that these items had before the Fubini theorems moved to their own paper.

\setcounter{theorem}{0}
\begin{imported}[Monomial decomposition]
\label{coh-decom}
\lean{DiamondIntegral.Section5.monomial_decomposition_5_12}
\leanok
Let $U \sub (\RV^\times)^n$ be definable. There is a finite definable partition $(D_i)_i$ of $\vrv(U)$ such that, writing $U_i = U \cap \vrv^{-1}(D_i)$, there are $M_i \in \operatorname{GL}_n(\Z)$, $0 \le m_i \le n$, a nonempty definable set $V_i \sub (\RV^\times)^{m_i}$ and a nonempty definable set $I_i \sub \Gamma^{n-m_i}$ with $\vrv(V_i) = \{\beta_i\}$ and $h_{M_i}(U_i) = V_i \times I_i^\sharp$. Here $h_M(u)_a = \prod_b u_b^{m_{ab}}$, $I^\sharp$ is the full preimage of $I$ under $\vrv$, and all sets and maps are definable over the parameter base.

\emph{Formalization.} The Lean theorem keeps the old number 5.12 in its name. It is stated over a VF-generated real base with a nontrivial definable value group, in an $\aleph_0$-saturated model of $\ACVF$ of equicharacteristic zero, and gives the matrices in a standard block order.
\end{imported}

\begin{proof}
\leanok
Proved as Lemma 2.1 of the Fubini paper (Forey and Yin, \emph{Fubini theorems for motivic volume}, in the Fubini-theorems repository); the Lean proof is in that package.
\end{proof}

\setcounter{theorem}{1}
\begin{imported}[Equivariant residue realization]
\label{bipo-gdv}
\lean{DiamondIntegral.Section5.EquivariantVariety.ModifiedRing, DiamondIntegral.Section5.PuiseuxOriginal.theta_5_14_concrete}
\leanok
\uses{coh-decom}
Work in the complex Puiseux series field $\puC$ over the base $\C((t))$, with value group $\Q$. Let $\gmv$ be the Grothendieck ring of complex varieties with good $\hat\mu$-actions (the action factors through some $\mu_N$ and every orbit lies in an affine open set), modulo $[X \times (\A^r, \rho)] - [X \times \A^r]$ for every linear action $\rho$ of some $\mu_N$ on $\A^r$. There is a canonical ring isomorphism $\Theta : !\ggk\RES \to \gmv$, built from the Kummer actions of $\hat\mu$ on the points of $\RES$ objects rational over $\C((t^{1/N}))$. For a monomial decomposition $h_{M_i}(U_i) = V_i \times I_i^\sharp$ of an object $\bm U$ of $\RV[*]$ with $\ell_i$ free $\RV$ coordinates, $(\Theta \circ \mathbb{E}_b)([\bm U]) = \sum_i \chi_b(I_i)\, \Theta([V_i])\, [\G_m]^{\ell_i}$.

\emph{Formalization.} The Lean isomorphism keeps the old number 5.14 in its name. Section 8 does not use the global formula (7.1): it evaluates on product generators instead (section 8/LEAN-INTERFACE.md).
\end{imported}

\begin{proof}
\uses{coh-decom}
\leanok
Notation 7.3 and formula (7.1) of the Fubini paper. The displayed formula is an evaluation of the retraction $\mathbb{E}_b$ of Section 4, so it does not depend on the decomposition.
\end{proof}

\setcounter{theorem}{2}
\begin{imported}[Motivic volume]
\label{mov-vol}
\lean{DiamondIntegral.Section5.PuiseuxOriginal.puiseuxBoundedVFVolume}
\leanok
\uses{bipo-gdv}
The motivic volume operator is $\Vol = \Theta \circ \mathbb{E}_b \circ \int : \ggk\VF_* \to \gmv$, where $\mathbb{E}_b$ is applied after the quotient by $(\bm P - 1)$, on which it is defined in Section 4.

\emph{Formalization.} In Lean the volume is defined in a given Puiseux ambient, with the bounded retraction and the lifting input of the Fubini package.
\end{imported}

\section{Special covariant bijections}
\label{section-6}
Special covariant bijections with bounded variable apertures, proper invariant sets, and the measured category $\mVF^{\diamond}$ with its Grothendieck ring.

\setcounter{theorem}{0}
\begin{definition}[Invariance and covariance]
\label{defn-binv}
\lean{DiamondIntegral.Section6.AlphaBetaCovariant, DiamondIntegral.Section6.AlphaInvariant, DiamondIntegral.Section6.ballProjection}
\leanok
\uses{rvfiber}
For $\gamma\in\Gamma$ let $\pi_\gamma:\VF\fun\VF/\mathcal{M}_\gamma$ be the quotient map; for $\gamma=(\gamma_1,\ldots,\gamma_n)\in\Gamma^n$, $\pi_\gamma$ is the product of the maps $\pi_{\gamma_i}$. Let $\alpha\in\Gamma^n$ and $\beta\in\Gamma^m$. A function $f:A\fun B$ with $A\sub\VF^n$ and $B\sub\VF^m$ is \emph{$(\alpha,\beta)$-covariant} if it contracts to a function $f_\downarrow:\pi_\alpha(A)\fun\pi_\beta(B)$, that is, $\pi_\beta\circ f=f_\downarrow\circ\pi_\alpha$. A set in $\VF^n$ is \emph{$\alpha$-invariant} if its characteristic function is $(\alpha,0)$-covariant. For sets $A$, $B$ with $\RV$-coordinates, a function $f:A\fun B$ is \emph{covariant} if every one of its $\VF$-fibers is covariant, and a set is \emph{invariant} if every one of its $\VF$-fibers is invariant.

\emph{Formalization.} Lean encodes covariance of a map with $\RV$-coordinates through its graph fibers, inside the definition of proper covariance (\texttt{ProperCovariant}).
\end{definition}

\setcounter{theorem}{1}
\begin{definition}[Proper covariance and proper invariance]
\label{prop-pseu}
\lean{DiamondIntegral.Section6.ProperCovariant, DiamondIntegral.Section6.ProperInvariant}
\leanok
\uses{defn-binv, bounded, rvfiber}
A definable function $f:A\fun B$ is \emph{proper covariant} if the sets $A_{\VF}$, $f(A)_{\VF}$ are bounded, the sets $A_{\RV}$, $f(A)_{\RV}$ are doubly bounded, and for each $\VF$-fiber $f_t$ of $f$ there is a $t$-definable tuple $(\alpha_t,\beta_t)$ in $\Gamma$ such that $f_t$ is $(\alpha_t,\beta_t)$-covariant, $\operatorname{dom}(f_t)$ is $\alpha_t$-invariant and $\operatorname{ran}(f_t)$ is $\beta_t$-invariant. A definable set $A$ is \emph{proper invariant} if $A_{\VF}$ is bounded, $A_{\RV}$ is doubly bounded, and each $\VF$-fiber $A_t$ is $\alpha_t$-invariant for some $t$-definable tuple $\alpha_t$ in $\Gamma$. For a powerset-valued map $f$ one uses its membership relation $R_f=\{(a,s):s\in f(a)\}$: the same global bounds, and proper invariance of every membership section.
\end{definition}

\setcounter{theorem}{3}
\begin{remark}[Closed sources need no image bound]
\label{s6-closed-image-db}
\lean{DiamondIntegral.Section6.remark_6_4, DiamondIntegral.Section6.remark_6_4_properCovariant, DiamondIntegral.Section6.remark_6_4_discharged}
\leanok
\uses{prop-pseu}
In Definition 6.2, if $A$ is closed and $\infty$ does not occur in the $\RV$-coordinates of $f(A)$, then the requirement that $f(A)_{\RV}$ be doubly bounded is implied by the other conditions of proper covariance.

\emph{Formalization.} The hypothesis that $\infty$ does not occur in the $\RV$-coordinates of $f(A)$ was added to the paper on November 22, after a counterexample and the author's approval of the correction on November 17 (section 6/GAP-REPAIRS.md; FINAL-CHECK.md, section 3). The Lean statement has this hypothesis and reads closedness on the $\VF$-fibers of $A$.
\end{remark}

\begin{proof}
\uses{vol-par-bounded, invar-db}
\leanok
There is an $\mathfrak{o}$-partition $p:A\fun\Gamma$ such that $f$ maps each open polydisc $\mathfrak{o}(a,p(a))\sub A$ into the range of a single $\VF$-fiber of $f$; by Lemma 2.53, $p(A)$ is doubly bounded. The definable surjection $A\fun\vrv(f(A)_{\RV})$ induced by $f$ then satisfies the hypothesis of Lemma 6.22, so $f(A)_{\RV}$ is doubly bounded.
\end{proof}

\setcounter{theorem}{4}
\begin{lemma}[Criterion for proper invariance]
\label{inv-crit}
\lean{DiamondIntegral.Section6.lemma_6_5}
\leanok
\uses{prop-pseu, bounded}
Let $A$ be a definable set such that $A_{\VF}$ is bounded and $A_{\RV}$ is doubly bounded. Then $A$ is proper invariant if and only if $A$ is clopen.
\end{lemma}

\begin{proof}
\uses{vol-par-bounded}
\leanok
Invariant fibers and their complements are unions of open polydiscs, which gives one direction. Conversely, on each fiber the least radius $\gamma\geq0$ with $\mathfrak{o}(a,\gamma)\sub A_t$ is an $\mathfrak{o}$-partition; since the fiber is closed and bounded, Lemma 2.53 bounds these radii, and their supremum is a $t$-definable invariance radius.
\end{proof}

\setcounter{theorem}{5}
\begin{lemma}[Boolean combinations of proper invariant sets]
\label{int-inv}
\lean{DiamondIntegral.Section6.lemma_6_6}
\leanok
\uses{prop-pseu, defn-disc}
If $A$, $B$ are proper invariant subsets of $\VF^n\times\RV^m$ then $A\cap B$, $\operatorname{RVH}(A)\mi A$ and $B\mi A$ are also proper invariant, where $\operatorname{RVH}(A)$ is the $\RV$-hull of $A$.
\end{lemma}

\begin{proof}
\uses{inv-crit}
\leanok
First, $\operatorname{RVH}(A)$ is proper invariant: a radius $\delta$ for which $A_t$ is $\delta$-invariant also works for $\operatorname{RVH}(A)_t$, and the hull preserves the bounds. Intersections and differences of clopen fibers are clopen and inherit the bounds, so Lemma 6.5 applies.
\end{proof}

\setcounter{theorem}{6}
\begin{lemma}[Continuous surjections are proper covariant]
\label{conti-covar}
\lean{DiamondIntegral.Section6.lemma_6_7}
\leanok
\uses{prop-pseu}
Let $f:A\fun B$ be a definable continuous surjection between two proper invariant sets such that every $\operatorname{dom}(f_t)$ is open. If every $\operatorname{ran}(f_t)$ is clopen then $f$ is proper covariant. If $f$ is bijective, the same holds under the weaker assumption that every $\operatorname{ran}(f_t)$ is open.
\end{lemma}

\begin{proof}
\uses{vol-par-bounded, inv-crit}
\leanok
The least radius of an open polydisc around a point of $\operatorname{ran}(f_{r,s})$ inside that range is an $\mathfrak{o}$-partition of a closed set, so its values are doubly bounded by Lemma 2.53. Openness of the domains and continuity then give an $\mathfrak{o}$-partition $p$ of $A$ whose polydiscs are mapped into polydiscs of one fixed radius, and Lemmas 6.5 and 2.53 bound $p(A)$. In the bijective case the fiber index is a locally constant definable function on $B$ and the same argument applies.
\end{proof}

\setcounter{theorem}{7}
\begin{lemma}[Proper covariant homeomorphisms]
\label{conv-unif}
\lean{DiamondIntegral.Section6.lemma_6_8, DiamondIntegral.Section6.lemma_6_8_discharged}
\leanok
\uses{prop-pseu}
If $f:A\fun B$ is a proper covariant surjection then both $A$ and $B$ are proper invariant. If, in addition, $f$ is continuous and bijective then $f$ is a proper covariant homeomorphism, that is, $f^{-1}$ is also proper covariant and continuous. In that case, if $A'$ is a proper invariant subset of $A$ then $f\rest A'$ is a proper covariant homeomorphism as well.
\end{lemma}

\begin{proof}
\uses{db-to-db, conti-homeo, conti-covar}
\leanok
Since $f_{\RV}$ is doubly bounded, Lemma 2.41 gives one sufficiently large $\epsilon$ such that every fiber $f_t$ is $(\epsilon,\beta_t)$-covariant and $\operatorname{dom}(f_t)$, $\operatorname{ran}(f_t)$ are $\epsilon$-invariant. By Corollary 2.49 each $f_t$ is a homeomorphism, and continuity of $f^{-1}$ gives its covariance as in Lemma 6.7. For $A'$ the fiber domains and ranges remain clopen, so the argument of Lemma 6.7 and the second claim apply.
\end{proof}

\setcounter{theorem}{8}
\begin{lemma}[Composites of proper covariant homeomorphisms]
\label{pr-co-home-com}
\lean{DiamondIntegral.Section6.lemma_6_9, DiamondIntegral.Section6.lemma_6_9_discharged}
\leanok
\uses{conv-unif}
If $f:A\fun B$ and $g:B\fun C$ are proper covariant homeomorphisms then $g\circ f$ is also a proper covariant homeomorphism.
\end{lemma}

\begin{proof}
\uses{conv-unif, conti-covar}
\leanok
By Lemmas 6.8 and 6.7 it suffices to show that the domain and range of every $\VF$-fiber of $g\circ f$ are open. With radii $\epsilon,\epsilon',\epsilon''$, each sufficiently larger than the next, the fibers of $f$ are $(\epsilon,\epsilon')$-covariant and those of $g$ are $(\epsilon',\epsilon'')$-covariant, so the fiber domains of $g\circ f$ are $\epsilon$-invariant and hence clopen; the same holds for $(g\circ f)^{-1}$.
\end{proof}

\setcounter{theorem}{9}
\begin{definition}[Regularization]
\label{conv-can}
\lean{DiamondIntegral.Section6.regularization, DiamondIntegral.Section6.regularizeMap, DiamondIntegral.Section6.regularization_properInvariant_iff}
\leanok
\uses{prop-pseu, bounded}
For each definable set $A$ let $\operatorname{can}(A)=\{(a,\rv(a),t):(a,t)\in A,\ a\in A_{\VF}\}$. The definable bijection $\operatorname{can}:A\fun\operatorname{can}(A)$ is the \emph{regularization} of $A$. For a proper invariant set $A$, $\operatorname{can}(A)$ is proper invariant if and only if $A_{\VF}$ is doubly bounded. When $A$ is doubly bounded and proper invariant, $\operatorname{can}(A)$ is (tacitly) substituted for $A$.
\end{definition}

\setcounter{theorem}{10}
\begin{definition}[Special covariant bijections]
\label{defn-special-bijection}
\lean{DiamondIntegral.Section6.ContractedApertureData, DiamondIntegral.Section6.RelativeApertureData, DiamondIntegral.Section6.RelativeApertureChain, DiamondIntegral.Section6.SpecialCovariantChain}
\leanok
\uses{conv-can, prop-pseu, defn-binv, defn-disc}
Let $A$ be a doubly bounded proper invariant set (or a finite disjoint union of such), regularized as in Convention 6.10. A \emph{centripetal transformation} in the first $\VF$-coordinate has as data an $\RV$-pullback $C\sub\operatorname{RVH}(A)$; a definable continuous focus $\lambda:\operatorname{pr}_{>1}(C\cap A)\fun\VF$ with graph in $C$, $(\beta_{\mathfrak p},\alpha_{\mathfrak p})$-covariant on each $\RV$-polydisc $\mathfrak p\sub C$, where $\mathfrak p\cap A$ is $(\alpha_{\mathfrak p},\beta_{\mathfrak p})$-invariant; and an \emph{aperture} $\tau:U\fun\RV$ on the retained labels with $0\leq\vrv(\tau)\leq\Delta$ for a definable $\Delta$. With $\gamma_{\mathfrak p}=\rad(\operatorname{pr}_1\mathfrak p)+\vrv(\tau(u_{\mathfrak p}))\geq\alpha_{\mathfrak p}$ and $t_{\mathfrak p}=\tau(u_{\mathfrak p})\rv(\operatorname{pr}_1\mathfrak p)$, it maps $(a,x)\in\mathfrak p\cap A$ to $(a-\lambda(x),x)$ if $\vv(a-\lambda(x))\leq\gamma_{\mathfrak p}$ and to $(t_{\mathfrak p},x)$ otherwise, with distinct tags, and is the identity off $C$. A \emph{bijective lift} adds a definable $b$ with values in $\VF^\times$, continuous on label fibers, with $\rv(b)=t_{\mathfrak p}$, and uses $(a-\lambda(x)+b(x),x)$ in the second clause. Alternating finite compositions of these transformations and regularizations are \emph{special covariant transformations}, of length the number of centripetal components; with lift data for every component they are \emph{special covariant bijections}, with image $A^\flat$.

\emph{Formalization.} \texttt{SpecialCovariantChain} is the constant-aperture subclass (with regularizations and the zero-centred preparations of Remark 6.13); \texttt{RelativeApertureChain} is the bounded variable-aperture class, and \texttt{ContractedApertureData} is the transformation before lifting (section 6/README.md, section 6/GAP-REPAIRS.md).
\end{definition}

\setcounter{theorem}{11}
\begin{remark}[Increasing the aperture]
\label{bigger-aper}
\lean{DiamondIntegral.Section6.ContractedApertureData.ApertureIncrease, DiamondIntegral.Section6.ContractedApertureData.ApertureIncrease.agree_outside, DiamondIntegral.Section6.ContractedApertureData.ApertureIncrease.old_core, DiamondIntegral.Section6.ContractedApertureData.ApertureIncrease.released_nonempty}
\leanok
\uses{defn-special-bijection}
Suppose two admissible aperture functions on the same locus satisfy $\gamma'(u)>\gamma(u)$ for every label $u$, keeping the full label and distinct core tags. On a polydisc $\mathfrak p$, increasing the aperture replaces the old endpoint $t_{\mathfrak p}$ by the relative annulus and the new endpoint, $\rv(\mathcal{M}_{\gamma_{\mathfrak p}}\mi\mathcal{M}_{\gamma'_{\mathfrak p}})\uplus\{t'_{\mathfrak p}\}$, distinguished by finite tags; outside the old core the two transformations agree. The comparison is fibrewise over $u$. If the old contracted clause occurs, the released part is nonempty; otherwise the two images agree.
\end{remark}

\begin{proof}
\leanok
If the old contracted clause occurs, invariance makes the whole old core present, so the strict increase of the aperture releases a nonempty part.
\end{proof}

\setcounter{theorem}{12}
\begin{remark}[Zero-centred preparation]
\label{spec-addendum}
\lean{DiamondIntegral.Section6.ZeroCenteredData, DiamondIntegral.Section6.ZeroCenteredData.image_vfDoublyBounded, DiamondIntegral.Section6.ZeroCenteredData.image_properInvariant, DiamondIntegral.Section6.ZeroCenteredData.properHomeomorphism}
\leanok
\uses{defn-special-bijection, prop-pseu, bounded}
Suppose that $A$ is proper $\delta$-invariant but not doubly bounded, so that there are no special covariant transformations on $A$. Increase $\delta$ to an attainable definable radius, choose a definable point $t_\delta\in\delta^\sharp$ and a definable $c\in\VF$ with $\rv(c)=t_\delta$. In each open polydisc of radius $\delta$ in $A$, replace every coordinate disc equal to $\mathcal{M}_\delta$ by $t_\delta$; this is the contracted preparation. Its bijective lift replaces $x_i\in\mathcal{M}_\delta$ by $x_i+c\in t_\delta^\sharp$, with finite disjoint-union tags. The lifted set is doubly bounded and remains $\delta$-invariant. These zero-centred translations, with derivative $1$, are included among the permitted preliminary operations.

\emph{Formalization.} Lean takes the nonzero definable centre $c$ with every $\VF$-fiber $\vv(c)$-invariant, which is the remark's enlarged attainable radius.
\end{remark}

\setcounter{theorem}{13}
\begin{lemma}[Special covariant bijections are proper homeomorphisms]
\label{spec-inve}
\lean{DiamondIntegral.Section6.lemma_6_14_bounded_variable_apertures, DiamondIntegral.Section6.lemma_6_14_bounded_variable_apertures_discharged, DiamondIntegral.Section6.lemma_6_14, DiamondIntegral.Section6.lemma_6_14_discharged}
\leanok
\uses{defn-special-bijection, conv-unif}
Special covariant bijections are proper covariant homeomorphisms. If $n$ is the $\VF$-arity, they are totally differentiable with Jacobian $1$ outside a definable set of $\VF$-dimension less than $n$.

\emph{Formalization.} \texttt{lemma\_6\_14} is the constant-aperture subclass and \texttt{lemma\_6\_14\_bounded\_variable\_apertures} the full class of Definition 6.11; in $\VF$-arity $0$ the unit derivative holds everywhere.
\end{lemma}

\begin{proof}
\uses{conti-covar, conv-unif, pr-co-home-com, spec-addendum}
\leanok
Each lifted centripetal branch is a shear in independent coordinates on a clopen branch domain, with continuous inverse, so the graph fibers have open domains and ranges; Lemmas 6.7, 6.8 and 6.9 give proper covariance and the homeomorphism property of the composite. Off the cylinders over the exceptional sets of the focus and endpoint maps, each branch has a triangular derivative with unit diagonal, and the chain rule outside the union of the pulled-back exceptional sets gives Jacobian $1$.
\end{proof}

\setcounter{theorem}{14}
\begin{remark}[Special covariant bijections as isomorphisms]
\label{spec-is-iso}
\lean{DiamondIntegral.Section6.DiamondSpecialIsomorphisms.remark_6_15, DiamondIntegral.Section6.DiamondSpecialIsomorphisms.remark_6_15_discharged}
\leanok
\uses{defn-VF-mu-G, defn-special-bijection, mor-equi}
Let $(A,\omega)$, $(B,\sigma)$ be objects of $\mVF[k]$ and $F:A\fun B$ a special covariant bijection. Then $F$ represents a $\mVF[k]$-isomorphism if $\omega(x)=\sigma(F(x))$ for all $x\in A$ outside a definable subset of $\VF$-dimension less than $k$.

\emph{Formalization.} Lean states it for objects of $\mVF^{\diamond}[k]$ and relative aperture chains, concluding a morphism of $\mVF^{\diamond}[k]$; the version for general objects of $\mVF[k]$ was not pursued (author's triage of November 18, section 6/GAP-REPAIRS.md; FINAL-CHECK.md, section 3).
\end{remark}

\begin{proof}
\uses{spec-inve, mor-equi}
Include the exceptional set of Lemma 6.14, off which $F$ has Jacobian $1$, in the given small set; the measured morphism convention of Remark 3.7 permits such a set.
\end{proof}

\setcounter{theorem}{15}
\begin{lemma}[Yin, Special transformations, Theorem 5.5]
\label{special-bi-term-constant}
\lean{DiamondIntegral.Section6.lemma_6_16}
\leanok
\uses{defn-corr-cont}
Suppose that $A$ is a definable set in $\VF$ and $f:A\fun\mathcal{P}(\RV^m)$ is a definable function. Then there is a special bijection $T$ on $A$ such that $A^\flat$ is an $\RV$-pullback and the function $f\circ T^{-1}$ is $\rv$-contractible.

\emph{Formalization.} Unary: $A\sub\VF$, which is the paper's reading of ``a set in $\VF$'' (FINAL-CHECK.md, section 3). The multi-dimensional extension (Remark 6.17) was dropped by author decision (section 6/GAP-REPAIRS.md).
\end{lemma}

\begin{proof}
\leanok
Apply quantifier elimination to the membership relation $R\sub A\times\RV^m$ and prepare simultaneously the finitely many $\VF$-polynomials in a quantifier-free formula for $R$ and for $A$ by the cited preparation theorem. Their $\RV$-values are then constant on each resulting $\RV$-polydisc, so every atomic formula, and hence the whole formula, has constant truth value there for each output $\RV$-tuple.
\end{proof}

\setcounter{theorem}{16}
\begin{remark}[Special covariant bijections that make a set an RV-pullback]
\label{special-pseudo}
\uses{special-bi-term-constant}
The proof of Lemma 6.16 shows that for every definable set $A$ there is a special bijection $T$ on $\mathrm{RVH}(A)$ such that $A^\flat$ is an $\RV$-pullback. For proper invariant sets one wants $T$ special covariant; Lemma 6.16 does not guarantee this, since its focus maps need not be covariant within each $\RV$-polydisc, except when $A_{\VF} \sub \VF$, where $T$ can be turned into a special covariant bijection whose components all have the same aperture.

\emph{Formalization.} Not formalized: the multi-dimensional form of Lemma 6.16 and this remark were dropped by the author (FINAL-CHECK.md, "Not formalized, and who decided"). Lemma 6.16 is formalized for unary sets, the case the remark singles out.
\end{remark}

\setcounter{theorem}{18}
\begin{proposition}[Simultaneous special covariant preparation]
\label{simplex-with-hole-rvproduct}
\lean{DiamondIntegral.Section6.proposition_6_19, DiamondIntegral.Section6.proposition_6_19_discharged, DiamondIntegral.Section6.proposition_6_19_vfGenerated, DiamondIntegral.Section6.lemma_6_19_zero_arity}
\leanok
\uses{prop-pseu, defn-special-bijection}
Let $H$ be a proper invariant $\RV$-pullback, not necessarily doubly bounded, and $(A_i)_i$ a definable finite partition of $H$ such that each $A_i$ is proper invariant. Then there is a special covariant bijection $T$ on $H$ such that every $T\rest A_i$ is a special covariant bijection and every $A_i^\flat\sub H^\flat$ is a doubly bounded $\RV$-pullback.

\emph{Formalization.} Lean allows any finite family of proper invariant subsets (not necessarily a partition) and produces a chain of the constant-aperture subclass, as the proof does.
\end{proposition}

\begin{proof}
\uses{spec-addendum, special-bi-term-constant, int-inv, inv-crit, vol-par-bounded, db-to-db, RV-fiber-dim-same, part-rv-cons, ball-to-ac, dcl-to-ac, prop-continuity-fiberwise, special-pseudo}
\leanok
After the zero-centred preparation of Remark 6.13, induct on the $\VF$-dimension using the constant-aperture subclass; the unary case is Lemma 6.16 with Remark 6.17. In the inductive step the fibrewise transformations over the first coordinate are fused: a special bijection from Lemma 6.16 makes their $\RV$-data constant, FMT (Terminology 6.18) freezes finitely many discs, and Lemmas 2.53 and 2.41 give uniform bounds and a common aperture. An inner induction on the components handles successive aperture changes on the released blocks, using Lemma 6.5 for the projected pieces.
\end{proof}

\setcounter{theorem}{19}
\begin{remark}[Contractible powerset-valued data]
\label{cons-extra}
\lean{DiamondIntegral.Section6.remark_6_20, DiamondIntegral.Section6.remark_6_20_discharged, DiamondIntegral.Section6.remark_6_20_unary}
\leanok
\uses{simplex-with-hole-rvproduct, prop-pseu, defn-corr-cont}
If $f:H\fun\mathcal{P}(\RV^m)$ is a proper covariant function, then one can add to the conclusion of Proposition 6.19 that $f\circ T^{-1}$ is $\rv$-contractible.

\emph{Formalization.} Proper covariance of $f$ is given as proper invariance of its membership relation, as in the text after Definition 6.2 (section 6/GAP-REPAIRS.md).
\end{remark}

\begin{proof}
\uses{simplex-with-hole-rvproduct, special-bi-term-constant}
\leanok
Include this clause in the inductive hypothesis of the proof of Proposition 6.19, as in Lemma 6.16.
\end{proof}

\setcounter{theorem}{20}
\begin{lemma}[Measured special covariant preparation]
\label{L-measure-surjective}
\lean{DiamondIntegral.Section6.lemma_6_21, DiamondIntegral.Section6.lemma_6_21_discharged, DiamondIntegral.Section6.lemma_6_21_vfGenerated, DiamondIntegral.Section6.lemma_6_21_zero_arity}
\leanok
\uses{prop-pseu, simplex-with-hole-rvproduct, defn-VF-mu-G, db-RV, RVcat-tag, defn-corr-cont}
A pair $(A,\omega)$ with $\omega:A\fun\Gamma$ definable is proper invariant if $A$ is proper invariant and every fiber of $\omega$ is proper invariant. Suppose $A$ is an $\RV$-pullback and $(A_i)_i$ is a definable finite partition of $A$ such that each $(A_i,\omega\rest A_i)$ is proper invariant. Then there is a special covariant bijection $T$ as in Proposition 6.19 such that $\omega\circ T^{-1}$ is $\rv$-contractible. In particular, if $(A_i,\omega\rest A_i)\in\mVF[k]$ then $T$ induces a morphism $(A_i,\omega\rest A_i)\fun\mL(\mathbf{U}_i,\pi_i)$, where $(\mathbf{U}_i,\pi_i)\in\mRV^{\db}[k]$.
\end{lemma}

\begin{proof}
\uses{simplex-with-hole-rvproduct, cons-extra, RV-fiber-dim-same, 0coor, spec-is-iso}
\leanok
View $\omega$ as a map $A\fun\mathcal{P}(\RV)$. Proposition 6.19 with Remark 6.20 gives $T$ such that each $(\omega\circ T^{-1})\rest A_i^\flat$ is constant on $\RV$-polydiscs and induces $\pi_i$ on $\rv(A_i^\flat)$. If $(A_i,\omega\rest A_i)\in\mVF[k]$, Lemma 3.3 gives $\mathbf{U}_i=\rv(A_i^\flat)_{\leq k}\in\RV^{\db}[k]$ (Notation 3.9), and Remark 6.15 identifies $(A_i,\omega\rest A_i)$ with $(\mathbb{L}\mathbf{U}_i,\sigma_i)$.
\end{proof}

\setcounter{theorem}{21}
\begin{lemma}[Locally constant $\Gamma$-valued functions are bounded]
\label{invar-db}
\lean{DiamondIntegral.Section6.lemma_6_22, DiamondIntegral.Section6.lemma_6_22_discharged}
\leanok
\uses{bounded}
Let $A$ be a definable set with $A_{\VF}$ bounded and $A_{\RV}$ doubly bounded. Let $f:A\fun\Gamma$ be a definable function that is constant on $A\cap\mathfrak{a}$ for all open polydiscs $\mathfrak{a}$ of radius $\gamma\in\Gamma$. Then $f(A)$ is doubly bounded.
\end{lemma}

\begin{proof}
\uses{L-measure-surjective, db-to-db}
\leanok
Replacing $A$ by its $\gamma$-tubular neighbourhood makes $(A,f)$ proper invariant. By Lemma 6.21, $f$ factors through a doubly bounded definable set in $\RV$, so $f(A)$ is doubly bounded by Lemma 2.41.
\end{proof}

\setcounter{theorem}{22}
\begin{corollary}[Values of a proper covariant form]
\label{s6-form-values-db}
\lean{DiamondIntegral.Section6.corollary_6_23, DiamondIntegral.Section6.corollary_6_23_discharged}
\leanok
\uses{L-measure-surjective}
If $(A,\omega)$ is proper invariant then $\omega(A)$ is doubly bounded.
\end{corollary}

\begin{proof}
\uses{vol-par-bounded, invar-db}
\leanok
By Lemma 2.53, for all sufficiently large $\beta\in\Gamma$ every fiber of $\omega$ is $\beta$-invariant, and Lemma 6.22 applies.
\end{proof}

\setcounter{theorem}{23}
\begin{definition}[Relatively unary bijections and morphisms]
\label{rela-unary}
\lean{DiamondIntegral.Section6.VFUnaryRelativeAt, DiamondIntegral.Section6.RVUnaryRelativeAt}
\leanok
\uses{defn-special-bijection, defn-c-RV-cat}
Let $A\sub\VF^n\times\RV^m$ and $B\sub\VF^n\times\RV^{m'}$ be definable. A definable bijection $G:A\fun B$ is \emph{relatively unary in the $i$th coordinate}, where $i\in[n]$, if $(\operatorname{pr}_{\tilde i}\circ G)(x)=\operatorname{pr}_{\tilde i}(x)$ for all $x\in A$. If, in addition, $G\rest A_a$ is a special covariant bijection for every $a\in\operatorname{pr}_{\tilde i}(A)$, then $G$ is \emph{relatively special covariant} in the $i$th coordinate. For objects $\mathbf{U}=(U,f)$, $\mathbf{V}=(V,g)$ of $\RV[k]$, a morphism $F:\mathbf{U}\fun\mathbf{V}$ is \emph{relatively unary in the $i$th coordinate}, where $i\in[k]$, if $(\operatorname{pr}_{\tilde i}\circ g\circ F)(u)=(\operatorname{pr}_{\tilde i}\circ f)(u)$ for all $u\in U$.

\emph{Formalization.} ``Relatively special covariant'' (the second sentence) is not formalized: it is not used after Example 6.25 and was dropped by author decision of 2026-11-19 (section 6/GAP-REPAIRS.md). The single-coordinate form of the definition is the author's decision of 2026-11-20, applied in the paper and in Lean.
\end{definition}

\setcounter{theorem}{25}
\begin{definition}[Standard contraction]
\label{defn-standard-contraction}
\lean{DiamondIntegral.Section6.StandardContraction, DiamondIntegral.Section6.OrderedStandardContractionProgram, DiamondIntegral.Section6.OrderedStandardContractionLift}
\leanok
\uses{prop-pseu, rela-unary, spec-addendum, conv-unif, pr-co-home-com, conv-can}
Let $A\sub\VF^n\times\RV^m$ be proper invariant, made doubly bounded by Remark 6.13 if necessary. For $i\in[n]$ let $T_i:A\fun A^\flat$ be a continuous proper covariant bijection, relatively unary in the $i$th coordinate, totally differentiable with Jacobian $1$ off a definable set of $\VF$-dimension less than $n$, passing on the $\RV$-coordinates of $A$, and such that every $T_i(A_{a,t})$ is an $\RV$-pullback. Contracting the $i$th $\VF$-coordinate of $A^\flat$ by $\rv$ gives a doubly bounded proper invariant set $A_{-i}\sub\VF^{n-1}\times\RV^{m_i}$ and a map $\hat T_i:A\fun A_{-i}$. Iterating along a permutation $\sigma$ of $[n]$ gives $\hat T_\sigma:A\fun\RV^l$, whose lift $T_\sigma=T_{\sigma(n)}\circ\cdots\circ T_{\sigma(1)}$ is a proper covariant homeomorphism. The function $\hat T_\sigma$, or the set $\hat T_\sigma(A)$, is called a \emph{standard contraction} of $A$.

\emph{Formalization.} Lean's \texttt{StandardContraction} starts with the permitted zero-centred normalization of Remark 6.13, followed by the ordered program and its lift (section 6/verification/status.json, scope).
\end{definition}

\setcounter{theorem}{26}
\begin{lemma}[Existence of standard contractions]
\label{exit-stand}
\lean{DiamondIntegral.Section6.lemma_6_27, DiamondIntegral.Section6.lemma_6_27_discharged, DiamondIntegral.Section6.lemma_6_27_normalized}
\leanok
\uses{defn-standard-contraction, prop-pseu}
Every proper invariant set $A\sub\VF^n\times\RV^m$ admits a standard contraction with respect to any permutation $\sigma$ of $[n]$.
\end{lemma}

\begin{proof}
\uses{simplex-with-hole-rvproduct, inv-crit, conti-covar, spec-inve}
\leanok
Make $\sigma(n)$ the base coordinate in the proof of Proposition 6.19 and strengthen its fibrewise induction to prescribe the order $\sigma(1),\ldots,\sigma(n-1)$, keeping the labels of contracted coordinates. Graft the $\RV$-coordinates of the deferred base preparation $R$ (a proper covariant homeomorphism by Lemmas 6.5 and 6.7), transport the fibrewise relatively unary maps along $R$, which keeps Jacobian $1$ as in Lemma 6.14, and apply $R$ to the last coordinate.
\end{proof}

\setcounter{theorem}{27}
\begin{remark}[Standard contractions of measured pairs]
\label{stand-vol}
\lean{DiamondIntegral.Section6.DiamondStandardContraction.remark_6_28, DiamondIntegral.Section6.DiamondStandardContraction.remark_6_28_exists, DiamondIntegral.Section6.DiamondStandardContraction.headStartSourceObject, DiamondIntegral.Section6.DiamondStandardContraction.remark_6_28_discharged}
\leanok
\uses{defn-standard-contraction, db-RV, RVcat-tag, defn-VF-mu-G, rela-unary, L-measure-surjective}
Suppose there is $k\in\{0\}\cup[m]$ (the \emph{head start}) with $(A_a)_{\leq k}\in\RV^{\db}[k]$ for every $a\in A_{\VF}$. Then $\hat T_\sigma(A)_{\leq n+k}\in\RV^{\db}[n+k]$. If $(A,\omega)$ is a proper invariant pair and $\omega\circ T_\sigma^{-1}$ is invariant on $\RV$-polydiscs, inducing $\omega_{\hat T_\sigma}$, then $(\hat T_\sigma(A),\omega_{\hat T_\sigma})_{\leq n+k}\in\mRV^{\db}[n+k]$ is a \emph{standard contraction} of $(A,\omega)$, and such $\hat T_\sigma$ exist. With $A^\sharp=\bigcup_{a\in A_{\VF}}a\times\mathbb{L}(A_a)_{\leq k}$, $T_\sigma$ induces a $\mVF[n+k]$-morphism $(A^\sharp,\omega)\fun\mL\hat T_\sigma(A,\omega)_{\leq n+k}$ that is a composition of relatively unary proper covariant homeomorphisms.

\emph{Formalization.} Lean gives the morphism as an isomorphism of $\mVF^{\diamond}[(n+1)+k]$ onto the lift of the contraction, for head-start objects with finite $\RV$-fibers; the head start with infinite $\RV$-fibers at the $\mVF^{\diamond}$ level was not pursued (author's triage of November 18, section 6/GAP-REPAIRS.md; FINAL-CHECK.md, section 3).
\end{remark}

\begin{proof}
\uses{RV-fiber-dim-same, exit-stand, L-measure-surjective, simplex-with-hole-rvproduct}
The first claim follows from Lemma 3.3. Existence follows by using Lemma 6.27 instead of Proposition 6.19 in the proof of Lemma 6.21.
\end{proof}

\setcounter{theorem}{28}
\begin{lemma}[Factoring RV morphisms into unary steps]
\label{bijection-partitioned-unary}
\lean{DiamondIntegral.Section6.AllDegreeUnaryFactorizations.lemma_6_29_finiteVolume_all_degrees, DiamondIntegral.Section6.AllDegreeUnaryFactorizations.lemma_6_29_measured_all_degrees, DiamondIntegral.Section6.AllDegreeUnaryFactorizations.lemma_6_29_doublyBounded_all_degrees, DiamondIntegral.Section6.AllDegreeUnaryFactorizations.lemma_6_29_measured_doublyBounded_all_degrees}
\leanok
\uses{rela-unary, defn-c-RV-cat, defn-RV-mu, db-RV}
Every morphism in $\RV[k]$ can be written as a composition of relatively unary morphisms, and similarly in $\mRV[k]$, $\RV^{\db}[k]$ and $\mRV^{\db}[k]$.

\emph{Formalization.} Every degree $k\geq0$. For $\RV[k]$ the faithful form is \texttt{lemma\_6\_29\_finiteVolume\_all\_degrees}; \texttt{lemma\_6\_29\_all\_degrees} omits the finite-volume conclusion (section 6/GAP-REPAIRS.md).
\end{lemma}

\begin{proof}
\uses{finite-fib-def, mor-ftf}
\leanok
Working piecewise, reduce to $F=\operatorname{id}$ and induct on $\dim_{\RV}(U)$, adjusting $f(u)$ towards $g(F(u))$ one coordinate at a time while keeping the coordinate maps finite-to-one; Lemma 2.15 controls the exceptional fibers of the finitary maps induced by $F^{\dagger}$ (Notation 3.11). For the measured versions the intermediate forms $\omega_h=\omega+\Sigma\vrv(f)-\Sigma\vrv(h)$ telescope, and in the doubly bounded case the intermediate coordinates are original coordinates or constant pads.
\end{proof}

\setcounter{theorem}{29}
\begin{lemma}[Lifting finitary RV functions]
\label{RVlift}
\lean{DiamondIntegral.Section6.lemma_6_30, DiamondIntegral.Section6.lemma_6_30_tuple, DiamondIntegral.Section6.lemma_6_30_tuple_vfGenerated}
\leanok
\uses{defn-corr-cont}
Let $U$, $V$ be sets in $\RV$ and $f:U\fun\mathcal{P}(V)$ a definable finitary function. Then there is a definable finitary function $f^\sharp:U^\sharp\fun\mathcal{P}(V^\sharp)$ that $\rv$-contracts to $f$. Moreover, for every $(u,v)\in f$, $f^\sharp_{u,v}=(u,v)^\sharp\cap f^\sharp$ is a partially differentiable function $u^\sharp\fun v^\sharp$.

\emph{Formalization.} \texttt{lemma\_6\_30} is the literal unary case; \texttt{lemma\_6\_30\_tuple} covers tuple arity with continuous, differentiable branches, the strengthened form used later (section 6/GAP-REPAIRS.md).
\end{lemma}

\begin{proof}
\uses{algebraic-balls-definable-centers, fun-almost-cont, diff-almost-every, RV-bou-dim}
\leanok
Induct on $\dim_{\RV}(U)$. Over a $\VF$-tuple $a$ the $\RV$-tuples in $f(\rv(a))$ are algebraic over $a$, so Lemma 2.20 over the $\VF$-generated base, compactness and averaging give one output in each prescribed $\RV$-polydisc. The non-continuity and non-differentiability loci have $\VF$-dimension less than $n$ (Lemmas 2.36 and 2.38), so their $\RV$-image has smaller $\RV$-dimension (Lemma 2.35); replace the lift there by the inductive one and glue along clopen sets.
\end{proof}

\setcounter{theorem}{30}
\begin{lemma}[Hrushovski--Kazhdan, Lemma 6.3]
\label{iso-lifted-vol}
\lean{DiamondIntegral.Section6.lemma_6_31_vfGenerated, DiamondIntegral.Section6.lemma_6_31_unmeasured_zero_arity, DiamondIntegral.Section6.lemma_6_31_measured_vfGenerated, DiamondIntegral.Section6.lemma_6_31_zero_arity}
\leanok
\uses{defn-c-RV-cat, RVcat-tag, mor-ftf, jcbG-op, defn-corr-cont}
Let $F:\mathbf{U}\fun\mathbf{V}$ be an $\RV[k]$-morphism, where $\mathbf{U}=(U,f)$ and $\mathbf{V}=(V,g)$. A \emph{lift} of $F$ is a definable bijection $F^\sharp:\mathbb{L}\mathbf{U}\fun\mathbb{L}\mathbf{V}$ that $\rv$-contracts to the function $U_f\fun V_g$ induced by $F$. There is a lift $F^\sharp$ of $F$ such that $\vv(\jcb_{\VF}F^\sharp(a,u))=\jcb_\Gamma F^{\dagger}(u)$ for all $(a,u)\in\mathbb{L}\mathbf{U}$.

\emph{Formalization.} Stated in the paper with a citation and no proof; proved in Lean, measured and unmeasured, in every degree, with no cited input (section 6/verification/status.json, scope).
\end{lemma}

\setcounter{theorem}{31}
\begin{remark}[Stretching RV-discs]
\label{rvdisc-stre}
\lean{DiamondIntegral.Section6.remark_6_32}
\leanok
\uses{defn-disc}
Let $t,s\in\RV$ be such that the $\RV$-discs $t^\sharp$, $s^\sharp$ contain definable points. For any definable $c\in\VF^\times$ with $\vv(c)=\vrv(s/t)$ there is a definable bijection $f:t^\sharp\fun s^\sharp$ such that $\frac{d}{dx}f=c$.
\end{remark}

\setcounter{theorem}{32}
\begin{lemma}[Lifting measured RV morphisms]
\label{L-measure-class-lift}
\lean{DiamondIntegral.Section6.lemma_6_33_vfGenerated, DiamondIntegral.Section6.lemma_6_33_doublyBounded_vfGenerated, DiamondIntegral.Section6.lemma_6_33_all_degrees_vfGenerated, DiamondIntegral.Section6.DiamondLiftHomomorphism.lemma_6_33_diamond}
\leanok
\uses{defn-RV-mu, RVcat-tag, rela-unary, defn-VF-mu-G, db-RV, conv-unif}
Every $\mRV[k]$-morphism $F:(\mathbf{U},\omega)\fun(\mathbf{V},\sigma)$ admits a lift $F^\sharp:\mL(\mathbf{U},\omega)\fun\mL(\mathbf{V},\sigma)$ that is piecewise a composition of partially differentiable relatively unary $\mVF[k]$-morphisms. Moreover, if $F$ is a $\mRV^{\db}[k]$-morphism then every component of the composition is a proper covariant homeomorphism, and hence so is $F^\sharp$.

\emph{Formalization.} Positive degree: \texttt{lemma\_6\_33\_vfGenerated} and its doubly bounded form; the first clause in every degree including $0$ is \texttt{lemma\_6\_33\_all\_degrees\_vfGenerated}, and the doubly bounded clause in every degree, as a $\mVF^{\diamond}$ morphism, is \texttt{lemma\_6\_33\_diamond} (section 6/verification/status.json).
\end{lemma}

\begin{proof}
\uses{rvdisc-stre, bijection-partitioned-unary, RVlift, iso-lifted-vol, fun-almost-cont, diff-almost-every, RV-bou-dim, conti-covar, conv-unif, pr-co-home-com}
\leanok
Induct on $\dim_{\RV}(U)$, requiring continuous branches; degree zero uses the affine lifts of Remark 6.32. Triangular corrections built from Lemma 6.30 reduce to $n=k$, and by Lemma 6.29 it suffices to lift one relatively unary morphism, done fibrewise by Lemma 6.31; over the bad set, of smaller $\RV$-dimension by Lemmas 2.36, 2.38 and 2.35, the lift is replaced by the inductive one. In the doubly bounded case Lemmas 6.7, 6.8 and 6.9 make every component and the composite proper covariant homeomorphisms.
\end{proof}

\setcounter{theorem}{33}
\begin{definition}[{The category $\mVF^{\diamond}[k]$}]
\label{defn-diamond}
\lean{DiamondIntegral.Section6.DiamondCategory.VFObject, DiamondIntegral.Section6.DiamondCategory.IsStep, DiamondIntegral.Section6.DiamondCategory.ObjectFactorization, DiamondIntegral.Section6.DiamondCategory.Morphism}
\leanok
\uses{defn-VF-mu-G, prop-pseu, rela-unary, conv-unif}
The subcategory $\mVF^{\diamond}[k]$ of $\mVF[k]$ consists of the proper invariant objects, that is, pairs $(A,\omega)\in\mVF[k]$ with $A$ proper invariant and $\omega$ proper covariant, and of the following morphisms. A morphism $F:\mathbf{A}\fun\mathbf{B}$ is a composition $F=F_l\circ\cdots\circ F_1$ of $\mVF[k]$-morphisms $F_m:\mathbf{A}_m\fun\mathbf{A}_{m+1}$, where $\mathbf{A}_1=\mathbf{A}$, $\mathbf{A}_{l+1}=\mathbf{B}$, every $\mathbf{A}_m$ is a proper invariant object, and every $F_m$ is a relatively unary proper covariant homeomorphism. Thus each intermediate object carries its own proper covariant volume form, and each factor $F_m$ is $\Gamma$-measure-preserving outside a definable set of $\VF$-dimension less than $k$.

\emph{Formalization.} Factorization through objects, rather than bare definable sets, is the author's ``fix (1)'' of 2026-11-20, applied to the paper and to Lean (section 6/GAP-REPAIRS.md).
\end{definition}

\setcounter{theorem}{34}
\begin{remark}[{$\mVF^{\diamond}[k]$ is a groupoid}]
\label{dia-cat-well}
\lean{DiamondIntegral.Section6.DiamondGroupoid.isomorphic_equivalence, DiamondIntegral.Section6.DiamondGroupoid.Morphism.inverse_exists, DiamondIntegral.Section6.DiamondGroupoid.Morphism.comp_exists, DiamondIntegral.Section6.DiamondGroupoid.IsMorphismMap}
\leanok
\uses{defn-diamond, mor-equi, simplex-with-hole-rvproduct, spec-inve, stand-vol, L-measure-class-lift}
The composition law holds in $\mVF^{\diamond}[k]$, the identity being the composition with no factors. Every morphism is an honest bijection with an inverse morphism: the inverse of each factor $F_m$ is a relatively unary, $\Gamma$-measure-preserving proper covariant homeomorphism $\mathbf{A}_{m+1}\fun\mathbf{A}_m$, so $F^{-1}=F_1^{-1}\circ\cdots\circ F_l^{-1}$. Thus $\mVF^{\diamond}[k]$ is a groupoid and no quotient as in Remark 3.7 is needed. Proposition 6.19, Lemma 6.14 and Remark 6.28 give nontrivial morphisms: for a special covariant bijection on a proper invariant object, the intermediate images with the transported forms are proper invariant objects, and similarly for the lifts of Lemma 6.33. The intermediate objects cannot be replaced by definable sets without volume forms: an explicit relatively unary proper covariant homeomorphism of $\mathcal{O}^2$ admits no proper covariant form on its target making it $\Gamma$-measure-preserving.

\emph{Formalization.} Formalized: composition, inverses, isomorphism as an equivalence relation in every degree, and a groupoid instance whose morphisms are the underlying maps (section 6/GAP-REPAIRS.md). The counterexample on $\mathcal{O}^2$ is not formalized, a choice of a formalization session with no author decision on record (FINAL-CHECK.md, section 3).
\end{remark}

\begin{proof}
\uses{spec-inve, conv-unif, bijection-partitioned-unary}
For a special covariant bijection $T=T_r\circ\cdots\circ T_1$ the transported forms $\omega\circ(T_i\circ\cdots\circ T_1)^{-1}$ have the proper invariant fibers $T_i\circ\cdots\circ T_1(A_\gamma)$ by Lemmas 6.14 and 6.8, and each stage is $\Gamma$-measure-preserving by Lemma 6.14. For the counterexample, $F(x,z)=(x,x\kappa(z/x))$ with $\kappa$ locally affine of slopes $1$, $t$, $1/t$ would force the form to take three values in every open polydisc around the fixed point $(0,0)$.
\end{proof}

\setcounter{theorem}{35}
\begin{remark}[Why morphisms are compositions of unary ones]
\label{s6-relatively-unary-morphisms}
\uses{L-measure-class-lift, simul-special-dim-1}
The lift in Lemma 6.33 must have its particular form because the morphisms of $\mVF^{\diamond}[k]$ are explicit compositions of relatively unary proper covariant morphisms, Lemma 6.38 having no higher-dimensional form. Every definable bijection between definable sets with the same number of $\VF$-coordinates is a composition of relatively unary bijections (Yin, \emph{Integration in algebraically closed valued fields}, Lemma 5.6). It is suspected, without a counterexample, that not every proper covariant homeomorphism is such a composition of relatively unary proper covariant homeomorphisms.

\emph{Formalization.} Not formalized: a cited claim that nothing later uses, and a conjecture (author's triage of November 18, FINAL-CHECK.md).
\end{remark}

\setcounter{theorem}{36}
\begin{corollary}[Surjectivity of the lifting map]
\label{L-sur-c}
\lean{DiamondIntegral.Section6.DiamondLiftHomomorphism.liftHom, DiamondIntegral.Section6.DiamondLiftHomomorphism.liftHom_surjective, DiamondIntegral.Section6.DiamondLiftHomomorphism.liftHom_surjective_discharged, DiamondIntegral.Section7.ChartOfMeasured.liftCls_surjective}
\leanok
\uses{def-Kgroup, db-RV, defn-diamond, RVcat-tag}
The lifting maps $\mL_k$ induce surjective homomorphisms of Grothendieck semigroups $\mL:\gsk\mRV^{\db}[k]\twoheadrightarrow\gsk\mVF^{\diamond}[k]$.

\emph{Formalization.} Every $k\geq0$, over a $\VF$-generated base (the section's standing assumption); \texttt{liftCls\_surjective} is the same statement on the Section 7 chart classes (section 7/COVERAGE.md).
\end{corollary}

\begin{proof}
\uses{L-measure-class-lift, L-measure-surjective, stand-vol}
\leanok
By Lemma 6.33 every $\mRV^{\db}[k]$-morphism lifts to a $\mVF^{\diamond}[k]$-morphism, so $\mL_k$ induces a map on isomorphism classes, which is a semigroup homomorphism. It is surjective by Lemma 6.21 (applied with the auxiliary piece $\operatorname{RVH}(A)\mi A$) or Remark 6.28.
\end{proof}

\setcounter{theorem}{37}
\begin{lemma}[Simultaneous preparation of a unary morphism]
\label{simul-special-dim-1}
\lean{DiamondIntegral.Section6.lemma_6_38_bounded_variable_apertures, DiamondIntegral.Section6.lemma_6_38_bounded_variable_apertures_contractible, DiamondIntegral.Section6.DiamondSimultaneous.lemma_6_38_diamond, DiamondIntegral.Section6.DiamondSimultaneous.lemma_6_38_diamond_contractible}
\leanok
\uses{defn-diamond, defn-special-bijection, db-RV, 0coor, defn-corr-cont}
Let $f:(A,\omega)\fun(B,\sigma)$ be a $\mVF^{\diamond}[1]$-morphism. Then there are special covariant bijections $T_A:A\fun A^\flat$, $T_B:B\fun B^\flat$, using the bounded variable-aperture data of Definition 6.11, such that $A^\flat$, $B^\flat$ are doubly bounded $\RV$-pullbacks, $f^\flat=T_B\circ f\circ T_A^{-1}$ is bijective, and $f^\flat$ $\rv$-contracts to a bijection $f^\flat_\downarrow:\rv(A^\flat)\fun\rv(B^\flat)$. Moreover $\omega'=\omega\circ T_A^{-1}$ and $\sigma'=\sigma\circ T_B^{-1}$ are $\rv$-contractible, and $f^\flat_\downarrow$ is a $\mRV^{\db}[1]$-morphism $(\rv(A^\flat),\omega')_1\fun(\rv(B^\flat),\sigma')_1$.
\end{lemma}

\begin{proof}
\uses{L-measure-surjective, lim-weak, db-to-db, spec-inve}
\leanok
Prepare both sides by Lemma 6.21 and take the ordinary unary matching preparations of Yin, Integration in ACVF, Lemma 5.2. Around paired centres $p$ and $q=f(p)$ choose definable whole discs with $f(D_p)=D_q$ (Lemma 2.51), and modify the preparations so that each centripetal step capturing a centre uses its offset as a bounded variable-aperture endpoint, bounded by Lemma 2.41. The literal conjugate then maps whole label fibers onto label fibers, and Lemma 6.14 with the chain rule gives the volume identity at every $\RV$-label.
\end{proof}

\section{Integrating proper invariant sets}
\label{section-7}
Blowups, 2-cells and the bounded integral: the ring isomorphism from $\ggk\mVF^{\diamond}[*]$ onto $\ggk\mRV^{\db}[*]/J_{\mathrm{var}}$ (Theorem 7.26), with kernel the relative annular ideal of Section 4, and the commutative squares of Corollary 7.27.

\setcounter{theorem}{0}
\begin{definition}[Shell and endpoint]
\label{punc-RVd}
\lean{DiamondIntegral.Section4.AnnularBlowup.shellPoints, DiamondIntegral.Section4.AnnularBlowup.endpointPoints}
\leanok
\uses{defn-disc}
Write $\RV^{\circ\circ}_\gamma = \rv(\mathcal{M}_\gamma \mi 0)$ as in Notation 2.5. For a finite $\gamma \in \Gamma$ with $\gamma \geq 0$ and $t \in \gamma^\sharp$, put\[ \RV^{\circ\circ}(t)_\gamma = (\RV^{\circ\circ} \mi \RV^{\circ\circ}_\gamma) \uplus t = \{ s \in \RV \mid 0 < \vrv(s) \leq \gamma \} \uplus t. \]The first part is the shell and the second the endpoint; they carry distinct tags, so for $\gamma > 0$ the endpoint is not identified with the element $t$ of the shell. If $\gamma = 0$ the shell is empty and $\RV^{\circ\circ}(t)_0 = t$. In all cases $\RV^{\circ\circ}(t)_\gamma$ is doubly bounded and does not contain $\infty$.

\emph{Formalization.} Formalized as data of Definition 7.2: the shell and the tagged endpoint are defined fibrewise over the domain of an aperture, with no separate declaration for a single $t$ (section 7/COVERAGE.md, ``done (as data)'').
\end{definition}

\setcounter{theorem}{1}
\begin{definition}[Annular blowups]
\label{defn-blowup-coa}
\lean{DiamondIntegral.Section4.RelativeAnnularGenerators.AnnularDatum, DiamondIntegral.Section4.RelativeAnnularGenerators.AnnularDatum.annularObject, DiamondIntegral.Section7.BlowupObjects.Blowup, DiamondIntegral.Section7.BlowupCongruence.BlowupSteps}
\leanok
\uses{punc-RVd, defn-RV-mu, defn-special-bijection, db-RV, relative-annular-ideal}
Let $\boldsymbol{U} = (U, f, \omega) \in \mRV[k]$, $k > 0$, with $\operatorname{pr}_{\tilde j} \rest f(U)$ finite-to-one for some $j \in [k]$. An aperture on $\boldsymbol{U}$ is a definable $\tau : U \fun \RV$ with a definable $\Delta \in \Gamma$ such that $0 \leq \gamma(u) := \vrv(\tau(u)) \leq \Delta$ on $U$. The elementary annular blowup in the $j$th coordinate of aperture $\tau$ is $\boldsymbol{U}^\flat = (U^\flat, f^\flat, \omega^\flat) \in \mRV[k]$, where $U^\flat = \bigcup_{u \in U} u \times \RV^{\circ\circ}(\tau(u))_{\gamma(u)}$, $f^\flat_j(u, s) = s f_j(u)$, $f^\flat_i(u, s) = f_i(u)$ for $i \neq j$, and $\omega^\flat(u, s) = \omega(u)$. For $\boldsymbol{V} = (V, g, \sigma) \in \mRV[k]$, pairwise disjoint definable $C_i \sub V$ with subobjects $\boldsymbol{C}_i = (C_i, g \rest C_i, \sigma \rest C_i)$, morphisms $F_i : \boldsymbol{U}_i \fun \boldsymbol{C}_i$ and elementary blowups $\boldsymbol{U}_i^\flat$, the object $(\boldsymbol{V} \mi \boldsymbol{C}) \uplus \biguplus_i \boldsymbol{U}_i^\flat$, where $\boldsymbol{C} = \bigcup_i \boldsymbol{C}_i$, is an annular blowup of $\boldsymbol{V}$ with locus $C = \bigcup_i C_i$. A blowup of length $n$ is a composition of $n$ annular blowups. For doubly bounded $\boldsymbol{U}$, $\boldsymbol{U}^\flat$ is the relative annular object $\boldsymbol{U}^{\tau, j}$ of Notation 4.17.

\emph{Formalization.} Formalized for doubly bounded objects, the only ones the paper uses (Remark 7.4), with $\infty$ allowed as a label tag (section 7/GAP-REPAIRS.md, item 3). Blowups of length one and $n$ are defined on isomorphism classes; the literal form with disjoint subobjects, \texttt{BlowupObjects.Blowup}, is equivalent by \texttt{step\_iff\_blowup} (GAP-REPAIRS item 2).
\end{definition}

\setcounter{theorem}{2}
\begin{remark}[Infinite apertures are excluded]
\label{s7-infinite-apertures}
\lean{DiamondIntegral.Section7.Remark73.remark_7_3, DiamondIntegral.Section7.Remark73.remark_7_3_lt}
\leanok
\uses{defn-blowup-coa, 0coor}
Infinite apertures are excluded: the aperture $\infty$ would replace $\RV^{\circ\circ}(t)_\gamma$ by $\rv(\mathcal{M})$ and give an elementary blowup in the sense of Yin's integration paper, whose endpoint lies in $\RV[{\leq}k]$ rather than $\RV[k]$ (Notation 3.9) and whose shell has unbounded valuation; this comparison is only heuristic. If $\boldsymbol{U} \in \mRV[k]$ has an elementary annular blowup then $\dim_{\RV}(U) < k$. For a given coordinate there is at most one elementary blowup of $\boldsymbol{U}$ but possibly many annular ones; the blown-up coordinate is part of the data but is usually suppressed from the notation.

\emph{Formalization.} Only the dimension claim is formalized, as $\dim_{\RV}(U) \leq k < 1 + k$ for an object of degree $1 + k$; the heuristic comparison with ordinary blowups is not. Author decision of 2026-11-19 (section 7/PLAN.md, section 5).
\end{remark}

\setcounter{theorem}{3}
\begin{remark}[Blowups of doubly bounded objects]
\label{db-blowup-bound}
\lean{DiamondIntegral.Section7.Remark74.remark_7_4_bounded, DiamondIntegral.Section7.Remark74.remark_7_4_named}
\leanok
\uses{defn-blowup-coa, db-RV}
If $\boldsymbol{U} \in \mRV^{\db}[k]$ then every elementary annular blowup $\boldsymbol{U}^\flat$ lies in $\mRV^{\db}[k]$. Conversely, for such $\boldsymbol{U}$ the bound $\Delta$ of an aperture is automatic: every definable $\tau : U \fun \RV$ has doubly bounded image, and $\vrv \circ \tau$ has a definable upper bound. From now on only annular blowups in $\mRV^{\db}[k]$ are considered, and they are simply called blowups.

\emph{Formalization.} In Lean an aperture carries some bound in $\Gamma$, not required to be definable; \texttt{remark\_7\_4\_named} produces a bound named over the base (section 7/GAP-REPAIRS.md, item 4).
\end{remark}

\begin{proof}
\uses{db-to-db}
\leanok
The new coordinate satisfies $0 < \vrv(s) \leq \Delta$. For the converse, Lemma 2.41 gives a doubly bounded image, and a bounded definable subset of $\Gamma$ has a definable upper bound.
\end{proof}

\setcounter{theorem}{4}
\begin{lemma}[Elementary blowups of isomorphic objects]
\label{elementary-blowups-preserves-iso-vol}
\lean{DiamondIntegral.Section7.BlowupResults.lemma_7_5}
\leanok
\uses{defn-blowup-coa, db-RV}
Let $\boldsymbol{U}, \boldsymbol{V} \in \mRV^{\db}[k]$ with $[\boldsymbol{U}] = [\boldsymbol{V}]$, and let $\boldsymbol{U}^\flat$, $\boldsymbol{V}^\flat$ be elementary blowups of $\boldsymbol{U}$, $\boldsymbol{V}$, possibly in different coordinates. Then there are blowups $\boldsymbol{U}^{\flat\flat}$, $\boldsymbol{V}^{\flat\flat}$ of $\boldsymbol{U}^\flat$, $\boldsymbol{V}^\flat$ of length $1$ such that $[\boldsymbol{U}^{\flat\flat}] = [\boldsymbol{V}^{\flat\flat}]$.
\end{lemma}

\begin{proof}
\uses{defn-RV-mu, defn-blowup-coa}
\leanok
Fix an isomorphism $F : \boldsymbol{U} \fun \boldsymbol{V}$ and split $U$ according to whether the depth of $\tau_U$ is at most that of $\tau_V \circ F$. Over each part, blow up the endpoints of the shallower side with the ratio aperture; since morphisms preserve total weight (Definition 3.12), explicit definable bijections onto the other side are isomorphisms, and they glue.
\end{proof}

\setcounter{theorem}{5}
\begin{corollary}[Length-one blowups of isomorphic objects]
\label{eblowup-same-loci-iso}
\lean{DiamondIntegral.Section7.BlowupResults.corollary_7_6}
\leanok
\uses{defn-blowup-coa, db-RV}
Let $F : \boldsymbol{U} \fun \boldsymbol{V}$ be a morphism of $\mRV^{\db}[k]$ and $\boldsymbol{U}^\flat$, $\boldsymbol{V}^\flat$ blowups of $\boldsymbol{U}$, $\boldsymbol{V}$ of length $1$. Then there are blowups $\boldsymbol{U}^{\flat\flat}$, $\boldsymbol{V}^{\flat\flat}$ of $\boldsymbol{U}^\flat$, $\boldsymbol{V}^\flat$ of length $1$ such that $[\boldsymbol{U}^{\flat\flat}] = [\boldsymbol{V}^{\flat\flat}]$.
\end{corollary}

\begin{proof}
\uses{elementary-blowups-preserves-iso-vol}
\leanok
Refine the two loci into the pieces $C_i \cap F^{-1}(D_l)$ and $F(C_i) \cap D_l$ and apply Lemma 7.5 on each. Blow up the part of each locus not met by the other through the transported elementary blowups; $F$ matches the parts outside all loci.
\end{proof}

\setcounter{theorem}{6}
\begin{lemma}[Confluence of blowups]
\label{blowup-equi-class-coa}
\lean{DiamondIntegral.Section7.BlowupResults.lemma_7_7}
\leanok
\uses{defn-blowup-coa, db-RV}
Let $\boldsymbol{U}, \boldsymbol{V} \in \mRV^{\db}[k]$ with $[\boldsymbol{U}] = [\boldsymbol{V}]$, and let $\boldsymbol{U}_1$, $\boldsymbol{V}_1$ be blowups of $\boldsymbol{U}$, $\boldsymbol{V}$ of lengths $m$, $n$. Then there are blowups $\boldsymbol{U}_2$, $\boldsymbol{V}_2$ of $\boldsymbol{U}_1$, $\boldsymbol{V}_1$ of lengths $n$, $m$ such that $[\boldsymbol{U}_2] = [\boldsymbol{V}_2]$.
\end{lemma}

\begin{proof}
\uses{eblowup-same-loci-iso}
\leanok
The formal argument of Lemma 6.5 of Yin's integration paper, which depends only on the analogue of Corollary 7.6.
\end{proof}

\setcounter{theorem}{7}
\begin{corollary}[Isomorphic blowups depend only on classes]
\label{blowup-equi-class}
\lean{DiamondIntegral.Section7.BlowupResults.corollary_7_8}
\leanok
\uses{defn-blowup-coa, db-RV}
Let $\boldsymbol{U}, \boldsymbol{U}', \boldsymbol{V}, \boldsymbol{V}' \in \mRV^{\db}[k]$ with $[\boldsymbol{U}] = [\boldsymbol{U}']$ and $[\boldsymbol{V}] = [\boldsymbol{V}']$. If $\boldsymbol{U}$, $\boldsymbol{V}$ have isomorphic blowups then so do $\boldsymbol{U}'$, $\boldsymbol{V}'$.
\end{corollary}

\setcounter{theorem}{8}
\begin{definition}[The blowup congruence]
\label{s7-mu-isp-db}
\lean{DiamondIntegral.Section7.BlowupCongruence.MuIsp, DiamondIntegral.Section7.FiberwiseInput.IspRel}
\leanok
\uses{defn-blowup-coa, db-RV, blowup-equi-class}
Let $\mispdb[k]$ be the set of pairs $(\boldsymbol{U}, \boldsymbol{V})$ of objects of $\mRV^{\db}[k]$ such that there exist isomorphic blowups $\boldsymbol{U}^\flat$, $\boldsymbol{V}^\flat$, and let $\mispdb[*] = \bigoplus_k \mispdb[k]$. By Corollary 7.8 these sets may be regarded as binary relations on isomorphism classes.

\emph{Formalization.} Defined on isomorphism classes in degree $1 + k$; \texttt{FiberwiseInput.IspRel} covers every degree, reading $\mispdb[0]$ as equality since there are no blowups in degree $0$ (section 7/GAP-REPAIRS.md, item 10).
\end{definition}

\setcounter{theorem}{9}
\begin{lemma}[The blowup relation is a congruence]
\label{isp-congruence-vol}
\lean{DiamondIntegral.Section7.ProductBlowups.lemma_7_10, DiamondIntegral.Section7.ProductBlowups.lemma_7_10_graded, DiamondIntegral.Section7.BlowupResults.lemma_7_10_addCon}
\leanok
\uses{s7-mu-isp-db}
$\mispdb[k]$ is a semigroup congruence relation on $\gsk \mRV^{\db}[k]$, and $\mispdb[*]$ is a semiring congruence relation on $\gsk \mRV^{\db}[*]$.

\emph{Formalization.} The semiring part is stated on homogeneous classes (an equivalence compatible with sums and products), with $\mispdb[0]$ read as equality; associativity of the class product is neither needed nor proved (section 7/COVERAGE.md; GAP-REPAIRS.md, item 10).
\end{lemma}

\begin{proof}
\uses{blowup-equi-class-coa}
\leanok
An easy consequence of Lemma 7.7, as in the proof of Lemma 6.8 of Yin's integration paper.
\end{proof}

\setcounter{theorem}{10}
\begin{lemma}[Special covariant bijections yield blowups]
\label{special-to-blowup-coa}
\lean{DiamondIntegral.Section7.SpecialChainInduction.lemma_7_11, DiamondIntegral.Section7.Lemma711Partial.lemma_7_11_partial}
\leanok
\uses{defn-special-bijection, defn-blowup-coa, RVcat-tag, db-RV}
Let $\boldsymbol{U} = (U, f, \omega) \in \mRV^{\db}[k]$ and let $T$ be a special covariant bijection on $\mathbb{L}(U, f)$. Let $\boldsymbol{U}_T = (U_T, \omega_T)_{\leq k} \in \mRV^{\db}[k]$, where $(U_T)_{\leq k} = (\mathbb{L}(U, f)^\flat_{\RV})_{\leq k}$ is the $\RV$-part of the image of $T$ and $\omega_T$ is the volume form induced by $\omega$. Then $\boldsymbol{U}_T$ is isomorphic to a blowup of $\boldsymbol{U}$ of the same length as $T$.

\emph{Formalization.} Lean gives the class of a blowup of some length, not of the same length as $T$; no later result uses the length (FINAL-CHECK.md, Notes on Section 7). The partial-lift generalization was approved by the author on 2026-11-19 (section 7/PLAN.md).
\end{lemma}

\begin{proof}
\uses{blowup-equi-class-coa, dcl-to-ac, defn-special-bijection, defn-blowup-coa}
Induction on the length of $T$ and Lemma 7.7 reduce to length $1$ with locus $\mathbb{L}(U, f)$. Lemma 2.13 and compactness show that $\operatorname{pr}_{\tilde j} \rest f(U)$ is finite-to-one. Then $(u, s) \mapsto (u, s f_j(u))$ is an isomorphism from the elementary blowup of aperture $\tau$ onto $\boldsymbol{U}_T$, by the explicit form of the centripetal transformation (Definition 6.11).
\end{proof}

\setcounter{theorem}{11}
\begin{lemma}[Blowups have isomorphic lifts]
\label{blowup-same-RV-coa}
\lean{DiamondIntegral.Section7.BlowupLiftGeneral.lemma_7_12, DiamondIntegral.Section7.Lemma712Objects.lemma_7_12_objects, DiamondIntegral.Section7.BlowupLiftGeneral.steps_lift}
\leanok
\uses{defn-blowup-coa, RVcat-tag, defn-diamond, db-RV}
Let $\boldsymbol{U} \in \mRV^{\db}[k]$ and let $\boldsymbol{U}^\flat$ be a blowup of $\boldsymbol{U}$. Then $\mL \boldsymbol{U}^\flat$ is isomorphic to $\mL \boldsymbol{U}$ in $\mVF^\diamond[k]$.

\emph{Formalization.} Object level only for elementary blowups; for blowups with a locus, and of length $n$, only the class form $[\mL \boldsymbol{U}^\flat] = [\mL \boldsymbol{U}]$ in $\gsk \mVF^\diamond[k]$, which is all that later results use. Left at class level by author decision of 2026-11-20 (section 7/STATUS.md; receipt field unfinished; GAP-REPAIRS.md, item 19).
\end{lemma}

\begin{proof}
\uses{L-measure-class-lift, RVlift, conti-covar, conv-unif, spec-inve}
By induction on the length, the disjoint-union structure of $\mL \boldsymbol{U}^\flat$ and Lemma 6.33, it suffices to treat an elementary blowup. Centres $b_0$, $b_1$ are lifted by the continuous-branch form of Lemma 6.30, and the two-branch translation $\Phi$ is a proper covariant homeomorphism by Lemmas 6.7 and 6.8, with triangular Jacobian as in Lemma 6.14; so it is a $\mVF^\diamond[k]$-morphism $\mL \boldsymbol{U} \fun \mL \boldsymbol{U}^\flat$.
\end{proof}

\setcounter{theorem}{12}
\begin{lemma}[The kernel in dimension one]
\label{kernel-dim-1-coa}
\lean{DiamondIntegral.Section7.KernelDimOne.lemma_7_13, DiamondIntegral.Section7.KernelBaseChange.lemma_7_13_of_support}
\leanok
\uses{defn-diamond, defn-standard-contraction, stand-vol, s7-mu-isp-db, db-RV}
Let $[\boldsymbol{A}] = [\boldsymbol{B}]$ in $\gsk \mVF^\diamond[1]$ and let $\boldsymbol{U}, \boldsymbol{V} \in \mRV^{\db}[1]$ be standard contractions of $\boldsymbol{A}$, $\boldsymbol{B}$. Then $([\boldsymbol{U}], [\boldsymbol{V}]) \in \mispdb$.

\emph{Formalization.} The Lean hypothesis is an isomorphism $\boldsymbol{A} \cong \boldsymbol{B}$ in $\mVF^\diamond[1]$ rather than equality in $\gsk \mVF^\diamond[1]$ (section 7/GAP-REPAIRS.md, item 8); the printed form follows from Proposition 7.24 (FINAL-CHECK.md, Notes on Section 7).
\end{lemma}

\begin{proof}
\uses{simul-special-dim-1, special-to-blowup-coa}
Immediate from Lemma 6.38, applied to $\mL \boldsymbol{U}$ and $\mL \boldsymbol{V}$, and Lemma 7.11.
\end{proof}

\setcounter{theorem}{13}
\begin{lemma}[Fibrewise blowup congruence]
\label{isp-VF-fiberwise-contract}
\lean{DiamondIntegral.Section7.Lemma714Reduce.lemma714, DiamondIntegral.Section7.FiberwiseInput.Lemma714}
\leanok
\uses{prop-pseu, defn-standard-contraction, stand-vol, s7-mu-isp-db}
Let $A'$, $A''$ be definable sets with $A'_{\VF} = A''_{\VF} = A \sub \VF^n$ and $\omega'$, $\omega''$ definable functions from $A'$, $A''$ into $\Gamma$, such that $\boldsymbol{A}' = (A', \omega')$ and $\boldsymbol{A}'' = (A'', \omega'')$ are proper invariant. Suppose there is $k \in \mathbb{N}$ such that, for every $a \in A$, over $\mathbb{S}\langle a \rangle$,\[ ([(A'_a, \omega')]_{\leq k}, [(A''_a, \omega'')]_{\leq k}) \in \mispdb. \]Then for any standard contractions $\hat T_\sigma$, $\hat R_\sigma$ of $\boldsymbol{A}'$, $\boldsymbol{A}''$,\[ ([\hat T_\sigma(\boldsymbol{A}')]_{\leq n+k}, [\hat R_\sigma(\boldsymbol{A}'')]_{\leq n+k}) \in \mispdb. \]

\emph{Formalization.} Formal reading (section 7/GAP-REPAIRS.md, items 9 and 10): over a substructure with support data $V \sub T \sub \operatorname{acl}(V)$, which covers $\mathbb{S}\langle a \rangle$, with $[\cdot]_{\leq k}$ read through head-start labels and $n \geq 1$. Formerly a temporary input (author decision of 2026-11-19), now proved (receipt field proved\_inputs).
\end{lemma}

\begin{proof}
\uses{simplex-with-hole-rvproduct, special-bi-term-constant, db-blowup-bound, simul-special-dim-1, special-to-blowup-coa, kernel-dim-1-coa, isp-congruence-vol}
\leanok
Induction on $n$ reduces to $n = 1$: apply that case over $\mathbb{S}\langle b \rangle$ to the first stage, then the inductive hypothesis with head start $k + 1$. For $n = 1$, following Lemma 6.14 of Yin's integration paper: choose uniform witnesses by compactness, freeze them on the discs of a special covariant bijection (Proposition 6.19, Lemma 6.16), glue them into $\mathbb{S}$-definable blowups with bounds from Remark 7.4, and compare with the given contractions by Lemmas 6.38 and 7.11 as in Lemma 7.13; conclude by Lemma 7.10.
\end{proof}

\setcounter{theorem}{14}
\begin{corollary}[Contractions across a relatively unary morphism]
\label{contraction-same-perm-isp}
\lean{DiamondIntegral.Section7.ProvedResults.corollary_7_15, DiamondIntegral.Section7.KernelDimOne.lemma_7_13}
\leanok
\uses{defn-diamond, rela-unary, defn-standard-contraction, stand-vol, s7-mu-isp-db}
Let $\boldsymbol{A}', \boldsymbol{A}'' \in \mVF^\diamond[n]$ and let $F : \boldsymbol{A}' \fun \boldsymbol{A}''$ be a morphism that is relatively unary in the $i$th $\VF$-coordinate. Then for any permutation $\sigma$ of $[n]$ with $\sigma(1) = i$ and any standard contractions $\hat T_\sigma$, $\hat R_\sigma$ of $\boldsymbol{A}'$, $\boldsymbol{A}''$,\[ ([\hat T_\sigma(\boldsymbol{A}')]_{\leq n}, [\hat R_\sigma(\boldsymbol{A}'')]_{\leq n}) \in \mispdb. \]

\emph{Formalization.} The Lean statement is for $n \geq 2$, over a substructure with support data, with $\sigma(1)$ written as the image of the last position under Lean's coordinate order (section 7/GAP-REPAIRS.md, items 9, 11, 12); the case $n = 1$ is Lemma 7.13, applied to the isomorphism $F$.
\end{corollary}

\begin{proof}
\uses{kernel-dim-1-coa, isp-VF-fiberwise-contract}
\leanok
For $n = 1$ this is Lemma 7.13. Otherwise Lemma 7.13 applies to the fibres over good points, where the exceptional set of $F$ is finite; the other points lie in the zero set of a nonzero polynomial, so a grid point near any $a$ is good, and by invariance the fibre condition holds at every $a$. Lemma 7.14 with head start $1$ gives the claim.
\end{proof}

\setcounter{theorem}{15}
\begin{definition}[$\vv$-affine and $\rv$-affine]
\label{rvaffine}
\lean{DiamondIntegral.Section7.TwoCells.VAffine, DiamondIntegral.Section7.TwoCells.RVAffine, DiamondIntegral.Section7.TwoCells.HasSlope}
\leanok
Let $\mathfrak{a}$ be an open disc and $f : \mathfrak{a} \fun \VF$ an injection. Then $f$ is $\vv$-affine if there is a (necessarily unique) $\gamma \in \Gamma$, the shift of $f$, such that $\vv(f(a) - f(a')) = \gamma + \vv(a - a')$ for all $a, a' \in \mathfrak{a}$; and $f$ is $\rv$-affine if there is a (necessarily unique) $t \in \RV$, the slope of $f$, such that $\rv(f(a) - f(a')) = t \rv(a - a')$ for all $a, a' \in \mathfrak{a}$.
\end{definition}

\setcounter{theorem}{16}
\begin{definition}[Balanced bijections]
\label{defn-balance}
\lean{DiamondIntegral.Section7.TwoCells.Balanced, DiamondIntegral.Section7.TwoCells.DiscToDisc}
\leanok
\uses{rvaffine, defn-dtdp}
Let $A \sub \VF^2$ be a definable set such that $\mathfrak{a}_1 := \operatorname{pr}_1(A)$ and $\mathfrak{a}_2 := \operatorname{pr}_2(A)$ are open discs, and let $f : \mathfrak{a}_1 \fun \mathfrak{a}_2$ be a definable bijection with the disc-to-disc property. Then $f$ is balanced in $A$ if $f$ is $\rv$-affine and there are $t_1, t_2 \in \RV_\infty$, the paradigms of $f$, such that, for every $a \in \mathfrak{a}_1$, $A_a = t_2^\sharp + f(a)$ and $f^{-1}(A_a) = a - t_1^\sharp$.

\emph{Formalization.} Geometric definition, with definability imposed where it is used (section 7/COVERAGE.md); the disc-to-disc condition is encoded in the forward (concentric) direction only, as noted for Definition 7.19 in FINAL-CHECK.md.
\end{definition}

\setcounter{theorem}{17}
\begin{remark}[Internal symmetry of balanced sets]
\label{2cell-linear}
\lean{DiamondIntegral.Section7.TwoCells.paradigm_zero_iff, DiamondIntegral.Section7.TwoCells.graph_of_zero, DiamondIntegral.Section7.TwoCells.slope_eq}
\leanok
\uses{defn-balance}
Let $f$ be balanced in $A$ with paradigms $t_1$, $t_2$. If one paradigm is $\infty$ then so is the other, and $A$ is the graph of $f$. Otherwise, for $a_1 \in \mathfrak{a}_1$ let $\mathfrak{o}_1$ be the maximal open subdisc containing $a_1$ of the closed disc of radius $\vrv(t_1)$ around $a_1$, and $\mathfrak{o}_2 = A_{a_1}$; then $\mathfrak{o}_2 = t_2^\sharp + f(\mathfrak{o}_1)$ and the fibre of $A$ over each $a_2 \in \mathfrak{o}_2$ is $f^{-1}(\mathfrak{o}_2) + t_1^\sharp = \mathfrak{o}_1$, so the slope of $f$ is $-t_2/t_1$. In the residue torsors $A$ is a line, and by Lemma 6.38 the obvious bijection $A_1 \times t_2^\sharp \to t_1^\sharp \times A_2$ is the lift of an $\RV^{\db}[2]$-morphism modulo special covariant bijections.

\emph{Formalization.} Formalized: the $\infty$-symmetry of the paradigms, $A$ as the graph of $f$, and the slope relation $s t_1 = -t_2$; the torsor line and the closing sentence are not. Author decision of 2026-11-19 (section 7/PLAN.md, section 5).
\end{remark}

\begin{proof}
\uses{defn-balance, simul-special-dim-1}
The symmetry is read off from the two fibre conditions of Definition 7.17, and $\rv$-affinity of $f$ then forces the slope; the last sentence is an application of Lemma 6.38, detailed in Lemma 7.21.
\end{proof}

\setcounter{theorem}{18}
\begin{definition}[1-cells and 2-cells]
\label{def-units}
\lean{DiamondIntegral.Section7.TwoCells.IsOneCell, DiamondIntegral.Section7.TwoCells.IsTwoCell, DiamondIntegral.Section7.TwoCells.PositioningCase}
\leanok
\uses{defn-balance, defn-dtdp}
A set $A$ is a $1$-cell if it is an open disc contained in a single $\RV$-disc or a point of $\VF$. A set $A$ is a $2$-cell if\begin{itemize}\item $A \sub \VF^2$ is contained in a single $\RV$-polydisc and $A_1 = \operatorname{pr}_1(A)$ is a $1$-cell;\item there are $\epsilon : A_1 \fun \VF$ and $t \in \RV$ with $A_a = t^\sharp + \epsilon(a)$ for every $a \in A_1$;\item $\epsilon$ is constant, or $\epsilon$ is injective with the disc-to-disc property and $\rad(\epsilon(A_1)) \geq \vrv(t)$, or $\epsilon$ is balanced in $A$.\end{itemize}Then $\epsilon$ is the positioning function and $t$ the paradigm of $A$. For sets with $\RV$-coordinates both notions are defined fibrewise over the $\RV$-coordinates.

\emph{Formalization.} Geometric definitions, with definability imposed where used. The disc-to-disc condition is forward only in Lean, so Lean's $2$-cells of the second kind are a priori more general (FINAL-CHECK.md, Notes on Section 7); the parametrized version enters through the uniform cell data of Lemma 7.20.
\end{definition}

\setcounter{theorem}{19}
\begin{lemma}[Decomposition into 2-cells]
\label{decom-into-2-units}
\lean{DiamondIntegral.Section7.Lemma720.lemma720, DiamondIntegral.Section7.TwoCellInput.Lemma720, DiamondIntegral.Section7.TwoCellInput.UniformTwoCells}
\leanok
\uses{prop-pseu, defn-standard-contraction, def-units}
Let $A \sub \VF^2$ be a proper invariant set. Then there is a standard contraction $\hat T_\sigma$ of $A$ such that $A_s$ is a $2$-cell for every $s \in \hat T_\sigma(A)$.

\emph{Formalization.} Formal reading (section 7/GAP-REPAIRS.md, item 13): retained $\RV$ labels allowed, fibres read in the normalized coordinates of the contraction, cell data uniform in $s$ and definable, $2$-cells as in Lean's Definition 7.19. Formerly a temporary input (author decision of 2026-11-19), now proved (receipt field proved\_inputs).
\end{lemma}

\begin{proof}
\uses{simplex-with-hole-rvproduct, special-bi-term-constant, rvaffine, def-units}
\leanok
A variation of Lemma 4.8 of Yin's integration paper, with FMT cutting out finitely many exceptional discs as in Proposition 6.19; this gives $\vv$-affine positioning functions. A refinement of the final unary stage, preparing the slope relation by Lemma 6.16 and FMT, makes them $\rv$-affine and keeps the kinds of cells of Definition 7.19.
\end{proof}

\setcounter{theorem}{20}
\begin{lemma}[Contracting a 2-cell in both orders]
\label{2-unit-contracted}
\lean{DiamondIntegral.Section7.ProductCells.cls_product_orders, DiamondIntegral.Section7.BalancedPiece.cls_balanced_orders, DiamondIntegral.Section7.ProductKinds.product_of_kind}
\leanok
\uses{def-units, defn-diamond, defn-standard-contraction}
Let $\boldsymbol{A} = (A, \omega) \in \mVF^\diamond[2]$, where $A \sub \VF^2$ is a $2$-cell and $\omega$ is constant. Then there are standard contractions $\hat T_{12}$, $\hat R_{21}$ of $\boldsymbol{A}$ such that $[\hat T_{12}(\boldsymbol{A})]_{\leq 2} = [\hat R_{21}(\boldsymbol{A})]_{\leq 2}$.

\emph{Formalization.} Proved in the uniform form that Lemma 7.22 uses: for families of product cells over $\RV$ labels (cells of the first two kinds are products), and for balanced cells on the balanced piece of the object $\boldsymbol{A}^*$ of Lemma 7.22, not for one arbitrary balanced $2$-cell (section 7/COVERAGE.md; FINAL-CHECK.md, Notes on Section 7).
\end{lemma}

\begin{proof}
\uses{def-units, simul-special-dim-1, 2cell-linear}
If the positioning function is constant or of the second kind, $A$ is a product and the same unary contractions in either order differ only by the order of coordinates. In the balanced case, translate by $\epsilon^{-1}$, respectively $\epsilon$, and apply Lemma 6.38 to the $\mVF^\diamond[1]$-morphism $\epsilon$; the map $(u, t_2) \mapsto (t_1, \psi(u))$ has $\Gamma$-Jacobian $0$, using the shift from Remark 7.18.
\end{proof}

\setcounter{theorem}{21}
\begin{lemma}[Reduction to 2-cells]
\label{subset-partitioned-2-unit-contracted}
\lean{DiamondIntegral.Section7.ProvedResults.lemma_7_22, DiamondIntegral.Section7.Lemma722Objects.prepObject, DiamondIntegral.Section7.Lemma722Cells.cellObject}
\leanok
\uses{defn-diamond, rela-unary, defn-standard-contraction}
Let $\boldsymbol{A} = (A, \omega) \in \mVF^\diamond[2]$. There are a morphism $\boldsymbol{A} \fun \boldsymbol{A}^*$, relatively unary in both coordinates, and standard contractions $\hat T_{12}$, $\hat R_{21}$ of $\boldsymbol{A}^*$ such that $[\hat T_{12}(\boldsymbol{A}^*)]_{\leq 2} = [\hat R_{21}(\boldsymbol{A}^*)]_{\leq 2}$.

\emph{Formalization.} Lean builds $\boldsymbol{A}^*$ on the normalized image, so the printed morphism becomes a chain graft, normalization, graft, and proves that any two standard contractions of $\boldsymbol{A}^*$ in the orders $12$, $21$ are $\mispdb$-related rather than that two have equal classes; this is the form Lemma 7.23 uses (section 7/GAP-REPAIRS.md, item 14; FINAL-CHECK.md, Notes on Section 7).
\end{lemma}

\begin{proof}
\uses{L-measure-surjective, cons-extra, decom-into-2-units, exit-stand, 2-unit-contracted, simul-special-dim-1}
Make $\omega$ $\rv$-contractible by Lemma 6.21 and Remark 6.20, apply Lemma 7.20 to the graph of the resulting $\RV$-data with retained labels, and graft the labels as in Lemma 6.27 to get $\boldsymbol{A}^*$, whose label fibres are $2$-cells with constant form. The constructions of Lemma 7.21 are uniform in the labels and glue by compactness.
\end{proof}

\setcounter{theorem}{22}
\begin{lemma}[Swapping the first two contracted coordinates]
\label{contraction-perm-pair-isp}
\lean{DiamondIntegral.Section7.Lemma723General.lemma_7_23, DiamondIntegral.Section7.ProvedResults.lemma_7_23_n2}
\leanok
\uses{defn-diamond, defn-standard-contraction, stand-vol, s7-mu-isp-db}
Let $\boldsymbol{A} = (A, \omega) \in \mVF^\diamond[n]$, let $i \neq j$ in $[n]$, and let $\sigma_1$, $\sigma_2$ be permutations of $[n]$ with $\sigma_1(1) = \sigma_2(2) = i$, $\sigma_1(2) = \sigma_2(1) = j$ and $\sigma_1 \rest \{3, \ldots, n\} = \sigma_2 \rest \{3, \ldots, n\}$. Then for any standard contractions $\hat T_{\sigma_1}$, $\hat T_{\sigma_2}$ of $\boldsymbol{A}$,\[ ([\hat T_{\sigma_1}(\boldsymbol{A})]_{\leq n}, [\hat T_{\sigma_2}(\boldsymbol{A})]_{\leq n}) \in \mispdb. \]

\emph{Formalization.} In Lean's indexing the orders are $\sigma$ and $\sigma \circ \tau$ with $\tau$ the swap of the last two positions, which are contracted first; the base is a substructure with support data (section 7/GAP-REPAIRS.md, item 9).
\end{lemma}

\begin{proof}
\uses{isp-VF-fiberwise-contract, contraction-same-perm-isp, subset-partitioned-2-unit-contracted}
\leanok
By Lemma 7.14 it suffices to treat the two-dimensional fibres $\boldsymbol{A}_a$, $a \in A_E$, $E = [n] \mi \{i, j\}$. By Corollary 7.15 and Lemma 7.14 it is then enough to find a morphism $\boldsymbol{A}_a \fun \boldsymbol{B}$, relatively unary in both coordinates, with two standard contractions of $\boldsymbol{B}$ in the two orders of equal class, which is Lemma 7.22.
\end{proof}

\setcounter{theorem}{23}
\begin{proposition}[The kernel of the lift]
\label{kernel-L}
\lean{DiamondIntegral.Section7.Proposition724.proposition_7_24, DiamondIntegral.Section7.Proposition724.proposition_7_24_muIsp, DiamondIntegral.Section7.Proposition724Zero.proposition_7_24_zero}
\leanok
\uses{db-RV, RVcat-tag, defn-diamond, s7-mu-isp-db}
For $\boldsymbol{U}, \boldsymbol{V} \in \mRV^{\db}[k]$,\[ [\mL \boldsymbol{U}] = [\mL \boldsymbol{V}] \text{ in } \gsk \mVF^\diamond[k] \quad \text{if and only if} \quad ([\boldsymbol{U}], [\boldsymbol{V}]) \in \mispdb. \]
\end{proposition}

\begin{proof}
\uses{blowup-same-RV-coa, L-sur-c, contraction-same-perm-isp, contraction-perm-pair-isp, exit-stand, isp-VF-fiberwise-contract, isp-congruence-vol}
\leanok
For $k = 0$ both conditions say $\boldsymbol{U} \cong \boldsymbol{V}$. The implication from right to left follows from Lemma 7.12 and Corollary 6.37. For the converse, first show by induction on $k$, with Corollary 7.15, Lemma 7.23, Lemma 6.27 and Lemma 7.14, that the class of a standard contraction modulo $\mispdb$ depends neither on the contraction nor on its order; then apply Corollary 7.15 to each relatively unary arrow from $\mL \boldsymbol{U}$ to $\mL \boldsymbol{V}$ and use transitivity (Lemma 7.10).
\end{proof}

\setcounter{theorem}{24}
\begin{remark}[The kernel is the relative annular ideal]
\label{thekernel}
\lean{DiamondIntegral.Section7.Theorem726.remark_7_25_kernel, DiamondIntegral.Section7.RingClasses.ispRel_mem, DiamondIntegral.Section4.RelativeAnnularScalar.radiusIdeal_le_manuscriptIdeal}
\leanok
\uses{s7-mu-isp-db, defn-blowup-coa, RVcat-tag, relative-annular-ideal, pgammabis}
The congruence on $\gsk \mRV^{\db}[k]$ induced by $\mL$ is $\mispdb[k]$, which is generated by the pairs $([\boldsymbol{U}], [\boldsymbol{U}^\flat])$ with $\boldsymbol{U}^\flat$ an elementary blowup. Consequently the kernel of the graded ring homomorphism $\ggk \mRV^{\db}[*] \fun \ggk \mVF^\diamond[*]$ induced by $\mL$ is the relative annular ideal $J_{\mathrm{var}}$ of Notation 4.17; no cancellation in $\gsk \mRV^{\db}[k]$ is used. For a constant aperture $t \in \gamma^\sharp$ in the $j$th coordinate, $[\boldsymbol{U}^\flat] - [\boldsymbol{U}] = [\boldsymbol{W}] \boldsymbol{P}_\gamma$ with $\boldsymbol{W} = (U, f_{\tilde j}, \omega + \vrv(f_j)) \in \mRV^{\db}[k-1]$, hence $(\boldsymbol{P}_\Gamma) \sub J_{\mathrm{var}}$ (equality is not asserted).

\emph{Formalization.} $\ggk \mVF^\diamond[*]$ is Lean's ring presentation (see Theorem 7.26). The inclusion $(\boldsymbol{P}_\Gamma) \sub J_{\mathrm{var}}$ is proved in Section 4 with Notation 4.17; the closing sentence on variable apertures refers to Lemma 4.18.
\end{remark}

\begin{proof}
\uses{kernel-L, defn-blowup-coa, blowup-same-RV-coa, L-sur-c, relative-annular-compatibility}
\leanok
The generators $[\boldsymbol{U}^\flat] - [\boldsymbol{U}]$ of $J_{\mathrm{var}}$ lie in the kernel by Lemma 7.12, and any blowup differs from its source by an element of $J_{\mathrm{var}}$. Conversely, if $[\boldsymbol{U}] - [\boldsymbol{V}]$ is in the kernel, Corollary 6.37 and Proposition 7.24 give isomorphic blowups of $\boldsymbol{U} \uplus \boldsymbol{W}$ and $\boldsymbol{V} \uplus \boldsymbol{W}$ for some $\boldsymbol{W}$, so the difference lies in $J_{\mathrm{var}}$.
\end{proof}

\setcounter{theorem}{25}
\begin{theorem}[Bounded integral as a quotient]
\label{main-prop}
\lean{DiamondIntegral.Section7.Theorem726.theorem_7_26, DiamondIntegral.Section7.Theorem726.theorem_7_26_graded, DiamondIntegral.Section7.Proposition724.integralEquiv, DiamondIntegral.Section8.Closed.theorem_7_26}
\leanok
\uses{defn-diamond, db-RV, s7-mu-isp-db, relative-annular-ideal, RVcat-tag}
For each $k \geq 0$ there is a canonical isomorphism of Grothendieck semigroups\[ \int^\diamond_+ : \gsk \mVF^\diamond[k] \fun \gsk \mRV^{\db}[k] / \mispdb \]such that $\int^\diamond_+ [\boldsymbol{A}] = [\boldsymbol{U}] / \mispdb$ if and only if $[\boldsymbol{A}] = [\mL \boldsymbol{U}]$. Taking the direct sum of the groupifications yields a canonical isomorphism of graded rings\[ \int^\diamond : \ggk \mVF^\diamond[*] \fun \ggk \mRV^{\db}[*] / J_{\mathrm{var}}. \]

\emph{Formalization.} The source ring is a presentation by objects and relations; that it is the direct sum of the groupifications, and a grading on it, are not Lean statements, and gradedness is stated as: the integral of a degree-$k$ object lies in the degree-$k$ piece (FINAL-CHECK.md, Notes on Section 7). \texttt{Closed.theorem\_7\_26} is the instance in the Puiseux ambient of Section 5, with every hypothesis supplied.
\end{theorem}

\begin{proof}
\uses{L-sur-c, kernel-L, isp-congruence-vol, thekernel}
\leanok
The semigroup isomorphism is the inverse of the map induced by $\mL$, which is surjective by Corollary 6.37 and injective by Proposition 7.24. Groupifying the quotient by a congruence divides by the differences of congruent pairs, which by Lemma 7.10 and Remark 7.25 generate the homogeneous ideal $J_{\mathrm{var}}$ in the direct sum.
\end{proof}

\setcounter{theorem}{26}
\begin{corollary}[Compatibility with the other two integrals]
\label{db-to-novol}
\lean{DiamondIntegral.Section7.Corollary727Rings.corollary_7_27_rings, DiamondIntegral.Section7.Corollary727Rings.corollary_7_27_left_ring, DiamondIntegral.Section7.Corollary727Rings.corollary_7_27_right_ring}
\leanok
\uses{main-prop, main-prop-HK, relative-annular-compatibility, dia-cat-well}
Consider the diagram with top row $\ggk \VF_* \leftarrow \ggk \mVF^\diamond[*] \rightarrow \ggk \mVF[*]$, bottom row $\ggk \RV[*] / (\boldsymbol{P} - 1) \leftarrow \ggk \mRV^{\db}[*] / J_{\mathrm{var}} \rightarrow \ggk \mRV[*] / (\boldsymbol{P})$ and vertical arrows $\int$, $\int^\diamond$, $\int^\mu$, where the top arrows are induced by the forgetful functor and by the subcategory relation and the bottom arrows come from Lemma 4.18. Both squares commute.

\emph{Formalization.} Over the real part of a Section 3 ground context, with $\int$ taken for the proved lift record; the ring laws of the top-right arrow are read through $\int^\mu$, which is injective by Theorem 3.22 (section 7/GAP-REPAIRS.md, item 18).
\end{corollary}

\begin{proof}
\uses{dia-cat-well, relative-annular-compatibility, main-prop, main-prop-HK}
\leanok
The forgetful arrow exists because morphisms of $\mVF^\diamond[*]$ are bijections (Remark 6.35), and the bottom row exists by Lemma 4.18. Both squares commute because the image of $[\boldsymbol{A}]$ under each of the three integrals is represented by a standard contraction of $\boldsymbol{A}$.
\end{proof}

\setcounter{theorem}{27}
\begin{remark}[Geometric versus bounded retraction]
\label{eg-eb-factor}
\lean{DiamondIntegral.Section7.Remark728.remark_7_28, DiamondIntegral.Section7.Remark728.remark_7_28_zero}
\leanok
\uses{mov-vol, eb-no-eg, prop-eu-retr-k, main-prop-HK, db-RV, RVcat-tag}
The motivic volume operator of Notation 5.3 may be defined with $\mathbb{E}_g$ instead of $\mathbb{E}_b$, inverting $[\A]$ (Remark 4.4); for proper invariant objects this only changes the image by a factor. Namely, every proper invariant $A \in \VF_*$ with $\dim_{\VF}(A) = n$ is isomorphic to some $\mathbb{L} \boldsymbol{U}$ with $(\boldsymbol{U}, 0) \in \mRV^{\db}[n]$, and\[ \mathbb{E}_g\Big(\int [A]\Big) = \mathbb{E}_b\Big(\int [A]\Big) [\A]^{-n}. \]

\emph{Formalization.} Stated for the forgetful image of an object of $\mVF^\diamond$ over the real part of a Section 3 ground context, which carries strong homogeneity (bridge approved by the author on 2026-11-19, section 7/STATUS.md), for any accepted lift input defining $\int$, in degrees $n + 1$ and $0$ separately.
\end{remark}

\begin{proof}
\uses{main-prop, L-sur-c, prop-eu-retr-k}
\leanok
Theorem 7.26, or rather Corollary 6.37, gives $A \cong \mathbb{L} \boldsymbol{U}$ with $\boldsymbol{U}$ doubly bounded; the identity then follows from the isomorphism $\mu\Psi$ of Section 3 and the defining conditions of $\mathbb{E}_g$, $\mathbb{E}_b$ (Proposition 4.7).
\end{proof}

\section{Motivic Milnor fiber}
\label{section-8}
The zeta function of a polynomial over the complex Puiseux series, its rationality and its limit: the motivic Milnor fiber equals the bounded integral of the analytic Milnor fiber after $\Theta$ and $\mathbb{E}^{\diamond}$ (Theorem 8.11), and its motivic volume (Remarks 8.14 and 8.15).

\setcounter{theorem}{0}
\begin{definition}[Puiseux polynomials and evaluation of T]
\label{loc-rings}
\lean{DiamondIntegral.Section8.PuiseuxPolynomial, DiamondIntegral.Section8.PuiseuxPolynomial.laurentSubring, DiamondIntegral.Section8.PuiseuxPolynomial.eta, DiamondIntegral.Section8.Hm.Rloc}
\leanok
\uses{bipo-gdv}
For any ring $R$, $R[T^{\Q}]$ denotes the ring of Puiseux polynomials over $R$, that is, the group ring of $\Q$ over $R$. Let $\gmv[[\A]^{-1}][T,T^{-1}]$ be the obvious subring of $\gmv[[\A]^{-1}][T^{\Q}]$; the canonical image of $\gmv[T^{\Q}]$ in $\gmv[[\A]^{-1}][T^{\Q}]$ is denoted the same way. The assignment $T \efun [\A]$ determines a ring homomorphism \[ \eta : \gmv[[\A]^{-1}][T,T^{-1}] \fun \gmv[[\A]^{-1}]. \]
\end{definition}

\setcounter{theorem}{1}
\begin{lemma}[Lattice-fibre formula for the piecewise retraction]
\label{muhat-section}
\lean{DiamondIntegral.Section8.Hm.hmValue, DiamondIntegral.Section8.Hm.hmValue_iso, DiamondIntegral.Section8.Hm.hmValue_decomp, DiamondIntegral.Section8.Hm.hm}
\leanok
\uses{loc-rings, bipo-gdv, comp-form, defn-RV-mu, db-RV, coh-decom}
Let $m \in \Z^+$ and $\mathbf U = (U, f, \omega) \in \mRV^{\db}[k]$ with $U \sub \RV^n$, and write $\omega_f(u) = \omega(u) + \sum_i \vrv(f_i(u))$. Take a monomial decomposition $(U_i)_i$ of $U$ (Lemma 5.1), $h_{M_i}(U_i) = V_i \times I_i^\sharp$ with $\vrv(V_i) = \{\beta_i\}$ and $I_i \sub \Gamma^{\ell_i}$, such that $\omega_f \rest U_i$ $\vrv$-contracts to $\sigma_i : D_i = \vrv(U_i) \fun \Gamma$, and put $h_m(\mathbf U) = \sum_i \Theta([V_i])[\G_m]^{\ell_i} \sum_{\gamma \in D_{i,m}} T^{-m\sigma_i(\gamma)} \in \gmv[[\A]^{-1}][T^{\Q}]$, where $D_{i,m} = D_i \cap (m^{-1}\Z)^n$. Then \[ h_m(\mathbf U) = \sum_w \Theta([U^{(m)}_w])\, T^{-mw}, \qquad U^{(m)}_w = \{u \in U : \vrv(u) \in (m^{-1}\Z)^n,\ \omega_f(u) = w\}. \] In particular $h_m(\mathbf U)$ does not depend on the monomial decomposition and is invariant on isomorphism classes.

\emph{Formalization.} Lean takes the right-hand side as the definition of $h_m$ (\texttt{hmValue}) and proves invariance under isomorphism (\texttt{hmValue\_iso}), a formula from any finite decomposition of the weight fibres into classes (\texttt{hmValue\_decomp}) and that $h_m$ is a ring homomorphism (\texttt{hm}); the monomial-decomposition expression (8.2) itself is not stated. This is a choice of a formalization session with no author decision on record (FINAL-CHECK.md, section 3).
\end{lemma}

\begin{proof}
\uses{coh-decom, bipo-gdv, defn-RV-mu, jcbG-op, cut-to-hensel-substru}
For $\gamma \in D_{i,m}$ with $M_i \gamma = (\beta_i, \delta)$, the map $h_{M_i}$ is a definable bijection from $U_i \cap \vrv^{-1}(\gamma)$ onto $V_i \times \delta^\sharp$, and $\Theta([\delta^\sharp]) = [\G_m]^{\ell_i}$; summing over the pieces that make up $U^{(m)}_w$ gives the formula. For invariance, a morphism $F : \mathbf U \fun \mathbf V = (V, g, \pi)$ satisfies $\omega_f = \pi_g \circ F$ (Definition 3.12, Notation 2.8), and since $\C(\!(t^{1/m})\!)$ is functionally closed (Lemma 2.11 and the opening of Section 8), $F$ restricts to definable bijections $U^{(m)}_w \fun V^{(m)}_w$.
\end{proof}

\setcounter{theorem}{2}
\begin{lemma}[Vanishing on the relative annular ideal]
\label{PG-van}
\lean{DiamondIntegral.Section8.Hm.lemma_8_3, DiamondIntegral.Section8.Hm.hm_generator, DiamondIntegral.Section8.Hm.naturalRing, DiamondIntegral.Section8.Hm.etaHm}
\leanok
\uses{muhat-section, loc-rings, relative-annular-ideal, db-RV}
Let $h_m : \ggk \mRV^{\db}[*] \fun \gmv[[\A]^{-1}][T^{\Q}]$ be the ring homomorphism given by Lemma 8.2, and let $\ggk^\natural_m \mRV^{\db}[*] = h_m^{-1}(\gmv[[\A]^{-1}][T,T^{-1}])$. The homomorphism \[ \eta \circ h_m : \ggk^\natural_m \mRV^{\db}[*] \fun \gmv[[\A]^{-1}] \] vanishes on $J_{\mathrm{var}} \cap \ggk^\natural_m \mRV^{\db}[*]$, where $J_{\mathrm{var}}$ is the relative annular ideal of Notation 4.17.
\end{lemma}

\begin{proof}
\uses{muhat-section, relative-annular-ideal, bipo-gdv, cut-to-hensel-substru}
\leanok
For a generator $[\mathbf U^{\tau,j}] - [\mathbf U]$ of $J_{\mathrm{var}}$, partition the points of $U$ with values in the lattice into pieces $R$ on which $\omega_f$ and $\vrv \circ \tau$ are constant, with values $w$ and $M/m$ (functional closedness of $\C(\!(t^{1/m})\!)$). The endpoint graph over $R$ and the shell slices $R \times (i/m)^\sharp$, $1 \leq i \leq M$, contribute $\Theta([R])\, T^{-mw}([\A]T^{-1} - 1)\sum_{i=0}^{M-1} T^{-i}$, so $h_m(x) = ([\A]T^{-1} - 1)Q$ for every $x \in J_{\mathrm{var}}$. If moreover $x \in \ggk^\natural_m \mRV^{\db}[*]$, the projection $\operatorname{pr}_0$ onto integer exponents gives $h_m(x) = ([\A]T^{-1} - 1)\operatorname{pr}_0(Q)$, which $\eta$ sends to $0$.
\end{proof}

\setcounter{theorem}{3}
\begin{remark}[Natural subrings and well-defined coefficients]
\label{dia-to-db-dag}
\lean{DiamondIntegral.Section8.Hm.zetaQ, DiamondIntegral.Section8.Hm.annular_kills, DiamondIntegral.Section8.Truncation.integral_lift, DiamondIntegral.Section8.Truncation.lift_form_one}
\leanok
\uses{muhat-section, loc-rings, PG-van, db-RV, defn-diamond, main-prop, relative-annular-ideal}
If $\mathbf U = (U, f, l) \in \mRV^{\db}[*]$ has constant volume form $l \in m^{-1}\Z$, the exponents in $h_m(\mathbf U)$ are integers, so $[\mathbf U] \in \ggk^\natural_m \mRV^{\db}[*]$; put $\ggk^\natural \mRV^{\db}[*] = \bigcap_{m \in \Z^+} \ggk^\natural_m \mRV^{\db}[*]$. If $\mathbf A = (A, l) \in \mVF^\diamond[*]$ with $l$ constant, then $\int^\diamond [\mathbf A] = [(U, f, l)]/J_{\mathrm{var}}$ for some $(U, f)$; if $l \in \Z$, it lies in the image $\ggk^\natural \mRV^{\db}[*]/J_{\mathrm{var}}$ of $\ggk^\natural \mRV^{\db}[*]$, and, by Lemma 8.3, for every $m \in \Z^+$ the expression $(\eta \circ h_m \circ \int^\diamond)([\mathbf A])$ is a well-defined element of $\gmv[[\A]^{-1}]$. Finally, $\ggk^{\Z} \mRV^{\db}[*]$ and $\ggk^{\Z} \mVF^\diamond[*]$ are the subrings generated by the classes $[(U, f, l)]$, respectively $[(A, l)]$, with $l \in \Z$; the first is contained in $\ggk^\natural \mRV^{\db}[*]$, and $\int^\diamond$ maps the second into its image modulo $J_{\mathrm{var}}$.

\emph{Formalization.} The representation of $\int^\diamond[\mathbf A]$ by a lift with the same constant form is formalized for the form $0$ used by $\mathcal X_f$ (\texttt{integral\_lift}, \texttt{lift\_form\_one}), and well-definedness is covered by the stronger \texttt{zetaQ} and \texttt{annular\_kills}, which define the zeta series on all of $\ggk \mRV^{\db}[*]/J_{\mathrm{var}}$ through the integer-exponent part of $h_m$. The rings $\ggk^{\Z}$ are notation with no Lean counterpart (FINAL-CHECK.md, section 3, New in the final review).
\end{remark}

\begin{proof}
\uses{main-prop, PG-van}
Integrality of the exponents is read off from the formula defining $h_m$. A representative $[(U, f, l)]$ of $\int^\diamond[\mathbf A]$ is given by Theorem 7.26, and two representatives in $\ggk^\natural_m \mRV^{\db}[*]$ differ by an element of $J_{\mathrm{var}} \cap \ggk^\natural_m \mRV^{\db}[*]$, on which $\eta \circ h_m$ vanishes by Lemma 8.3.
\end{proof}

\setcounter{theorem}{4}
\begin{lemma}[Change of cutoff for truncations]
\label{s8-truncation-cutoff}
\lean{DiamondIntegral.Section8.Truncation.lemma_8_5_sets, DiamondIntegral.Section8.Truncation.lemma_8_5, DiamondIntegral.Section8.Truncation.fullClass_cutoff, DiamondIntegral.Section8.Truncation.truncation_definable}
\leanok
\uses{defn-binv, prop-pseu, bipo-gdv}
Let $m \in \Z^+$, $\beta \in (m^{-1}\Z)^n$ with all $\beta_i \geq 0$, and let $A \sub \mathcal O^n \times \RV^l$ be a proper $\beta$-invariant definable set such that $\operatorname{pr}_{\VF} \rest A$ is finite-to-one. Let $A[m;\beta] \sub \prod_i \C[t^{1/m}]/t^{\beta_i + 1/m} \times \RV(\C(\!(t^{1/m})\!))^l$ be its truncation: for every $r \in \RV(\C(\!(t^{1/m})\!))^l$ the fibre $A_r(\C(\!(t^{1/m})\!))$ is the preimage of $A[m;\beta]_r$; it is an object of $\RES_m$, the full subcategory of $\RES$ of sets whose coordinates are $\infty$ or have value in $m^{-1}\Z$. If $\beta' \in (m^{-1}\Z)^n$ and $\beta_i \leq \beta'_i$ for all $i$, then \[ [A[m;\beta']] = [A[m;\beta]]\,[\A]^{m \Sigma(\beta' - \beta)} \in {!\ggk}\RES. \]

\emph{Formalization.} Lean proves, in the concrete Puiseux model, the set identity $A[m;\beta'] = A[m;\beta] \times \prod L_\gamma$ (\texttt{lemma\_8\_5\_sets}), the definability of truncations (\texttt{truncation\_definable}) and the class identity after $\Theta$ in the localization at $[\A]$, as equality of normalized classes (\texttt{lemma\_8\_5}, \texttt{fullClass\_cutoff}); the identity in ${!\ggk}\RES$ itself is not stated. This is a choice of a formalization session with no author decision on record (FINAL-CHECK.md, section 3).
\end{lemma}

\begin{proof}
\uses{bipo-gdv}
By $\beta$-invariance, the coefficients of $t^\gamma$, $\gamma > \beta_i$, in the $i$-th coordinate of a $\C(\!(t^{1/m})\!)$-rational point of $A$ may be changed freely, so $A[m;\beta'] = A[m;\beta] \times \prod_i \prod_\gamma L_\gamma$, with $\gamma$ running over the $m(\beta'_i - \beta_i)$ elements of $m^{-1}\Z$ in $(\beta_i, \beta'_i]$. The equality follows since $[L_\gamma] = [\A]$ in ${!\ggk}\RES$ (Notation 5.2).
\end{proof}

\setcounter{theorem}{5}
\begin{lemma}[Zeta coefficients as truncation classes]
\label{inv-dir-com}
\lean{DiamondIntegral.Section8.Truncation.zeta_coeff_of_form_one, DiamondIntegral.Section8.Truncation.chain_teq, DiamondIntegral.Section8.Truncation.lift_class}
\leanok
\uses{defn-diamond, main-prop, loc-rings, muhat-section, s8-truncation-cutoff, bipo-gdv}
Suppose that $(A, 0) \in \mVF^\diamond[n]$ and $A \sub \mathcal O^n \times \RV^l$, and let $m \in \Z^+$. Let $\tilde A[m] = [A[m;\beta]]\,[\A]^{-m\Sigma\beta} \in {!\ggk}\RES[[\A]^{-1}]$, for $\beta \in (m^{-1}\Z)^n$ such that $A$ is $\beta$-invariant; by Lemma 8.5 it does not depend on $\beta$. Then, in $\gmv[[\A]^{-1}]$, \[ (\eta \circ h_m)\Big(\int^\diamond [(A, 0)]\Big) = \Theta(\tilde A[m]), \] where $\Theta$ also denotes the localization at $[\A]$ of the isomorphism of Notation 5.2.

\emph{Formalization.} Lean proves it for $n \geq 1$ (FINAL-CHECK.md, section 3), in the concrete Puiseux model: the $m$-th coefficient of the zeta series of $\int^\diamond[(A, 0)]$ equals the normalized truncation class of the trace of $A$ (\texttt{zeta\_coeff\_of\_form\_one}), under the side conditions of the Lean statement (cutoffs at which the trace is invariant, label patterns covering the labels, nonzero labels, integral rational points). Special covariant chains preserve normalized truncation classes (\texttt{chain\_teq}), and the lift case is \texttt{lift\_class}.
\end{lemma}

\begin{proof}
\uses{defn-special-bijection, conti-covar, conv-unif, spec-inve, spec-addendum, L-measure-surjective, muhat-section, s8-truncation-cutoff, cut-to-hensel-substru}
First, $\tilde A[m]$ is unchanged by each step of a special covariant bijection: for a lifted centripetal transformation followed by regularization (Definition 6.11), with cutoffs chosen large enough, first for the moved coordinate and then for the others (the source and target are proper invariant by Lemmas 6.14 and 6.8, and the focus maps are covariant as in Lemma 6.7), the transformation and its inverse induce a bijection of truncations, because $\C(\!(t^{1/m})\!)$ is functionally closed; the same holds for regularization and the translations of Remark 6.13. By Lemma 6.21 one reduces to $A = \mathbb L \mathbf U$ with $\mathbf U = (U, f) \in \RV^{\db}[n]$. Over each piece $R$ of the lattice points of $U$ on which $\vrv \circ f = \alpha$, the truncation is $R \times \prod_i \prod_\gamma L_\gamma$, so $\Theta(\tilde A[m]) = \sum_R \Theta([R])[\A]^{-m\Sigma\alpha}$, which is $(\eta \circ h_m)([(U, f, 0)])$ by Lemma 8.2.
\end{proof}

\begin{assumption}[\'Etale chart at the origin]
\label{hyp-etale-chart}
\lean{DiamondIntegral.Section8.Milnor.EtaleChart, DiamondIntegral.Section8.Milnor.etaleChart_affine}
\leanok
Let $X \sub \A^N$ be cut out near $0$ by $F_1, \ldots, F_k \in \C[y_1, \ldots, y_N]$, with tube $X(\mathcal M) = \{y \in \mathcal M^N : F(y) = 0\}$, and let $q = (q_1, \ldots, q_d)$ with $q_i \in \C[y]$ and $q(0) = 0$. Hypothesis: $q$ maps $X(\mathcal M)$ onto $\mathcal M^d$ and its inverse is $1$-Lipschitz, that is, for $y, y' \in X(\mathcal M)$ and $\gamma \in \Gamma$, if $\vv(q_i(y) - q_i(y')) > \gamma$ for all $i$ then $\vv(y_j - y'_j) > \gamma$ for all $j$. This is what Hensel's lemma gives when $q$ is \'etale at $0$ (Section 8.2: $q$ induces a bijection between $X(\mathcal M)$ and $\mathcal M^d$). It holds for $X = \A^d$ and $q$ the identity.

\emph{Formalization.} An explicit hypothesis of the smooth-$X$ results (\texttt{milnorObjectX}, \texttt{theorem\_8\_11\_smooth}, \texttt{milnor\_fiber\_vol\_smooth}, \texttt{remark\_8\_15\_milnor\_smooth}), not proved in general, by the author's instruction of 2026-11-21 (section 8/CITED-INPUTS.md; receipt field \texttt{accepted\_inputs}). It is proved for $X = \A^d$ (\texttt{etaleChart\_affine}).
\end{assumption}

\begin{assumption}[\'Etale coordinates at the origin]
\label{hyp-etale-coords}
\lean{DiamondIntegral.Section8.Truncation.EtaleCoords, DiamondIntegral.Section8.Truncation.EtaleCoords.chart, DiamondIntegral.Section8.Truncation.etaleCoords_affine}
\leanok
Let $X \sub \A^N$ be cut out near $0$ by $F = (F_1, \ldots, F_k)$, $F_j \in \C[y_1, \ldots, y_N]$, with tube $X(\mathcal M) = \{y \in \mathcal M^N : F(y) = 0\}$, and let $q = (q_1, \ldots, q_d)$ with $q_i \in \C[y]$. Hypothesis: $q(0) = 0$, $q$ maps $X(\mathcal M)$ onto $\mathcal M^d$, and $(F, q) : \mathcal M^N \fun \mathcal M^k \times \mathcal M^d$ has a $1$-Lipschitz inverse: for $y, y' \in \mathcal M^N$ and $\gamma \in \Gamma$, if $\vv(F_j(y) - F_j(y')) > \gamma$ and $\vv(q_i(y) - q_i(y')) > \gamma$ for all $i, j$, then $\vv(y_l - y'_l) > \gamma$ for all $l$. This follows from the multivariable Hensel lemma when $X = V(F)$ is smooth of dimension $d$ at $0$, $F$ generates its ideal near $0$ and $q \rest X$ is \'etale at $0$. It implies the \'etale chart hypothesis and holds for $X = \A^d$.

\emph{Formalization.} An explicit hypothesis of \texttt{z\_coeff\_change\_smooth} only, not proved in general, by the author's instruction of 2026-11-21 (section 8/CITED-INPUTS.md; receipt field \texttt{accepted\_inputs}). It implies the \'etale chart hypothesis (\texttt{EtaleCoords.chart}) and is proved for $X = \A^d$ (\texttt{etaleCoords\_affine}).
\end{assumption}

\setcounter{formula}{5}
\begin{formula}[Coefficients of the motivic zeta function]
\label{z-coeff-change}
\lean{DiamondIntegral.Section8.Truncation.z_coeff_change, DiamondIntegral.Section8.Truncation.milnor_zeta_coeff, DiamondIntegral.Section8.Truncation.z_coeff_change_jets, DiamondIntegral.Section8.Truncation.z_coeff_change_smooth}
\leanok
\uses{defn-diamond, main-prop, loc-rings, muhat-section, bipo-gdv, inv-dir-com, hyp-etale-chart, hyp-etale-coords}
Let $X$ be a smooth connected complex algebraic variety of dimension $d$, $f : X \fun \A^1$ a nonconstant regular function and $0 \in X(\C)$ with $f(0) = 0$. Let $\mathcal X_f = \{x \in X(\mathcal M) : \rv(f(x)) = \rv(t)\}$, identified through an \'etale chart at $0$ with a subset of $\mathcal M^d$ and regarded as an object of $\mVF^\diamond[d]$ with volume form $0$. For $m \geq 1$ let $\mathcal X_m = \{\varphi \in X(\C[t]/(t^{m+1})) : \varphi(0) = 0,\ f(\varphi) = t^m \bmod t^{m+1}\}$, with its natural $\mu_m$-action. Then, in $\gmv[[\A]^{-1}]$, \[ \Big(\eta \circ h_m \circ \int^\diamond\Big)([\mathcal X_f]) = \Theta(\tilde{\mathcal X}_f[m]) = [\mathcal X_m][\A]^{-md}. \] So the coefficients of the motivic zeta function $Z_f(T) = \sum_{m \in \Z^+} (\eta \circ h_m \circ \int^\diamond)([\mathcal X_f])\, T^m$ are the usual ones.

\emph{Formalization.} Equation (8.6) of the running text of Section 8.2, not a numbered item. Lean proves both equalities for $X = \A^d$ (\texttt{milnor\_zeta\_coeff}, \texttt{z\_coeff\_change\_jets}, \texttt{z\_coeff\_change}; checked in FINAL-CHECK.md, section 3) and, under the \'etale coordinates hypothesis taken by the author's instruction of 2026-11-21 (section 8/CITED-INPUTS.md), for smooth $X$ given near $0$ by polynomial equations in $\A^N$ (\texttt{z\_coeff\_change\_smooth}).
\end{formula}

\begin{proof}
\uses{dia-to-db-dag, inv-dir-com}
\leanok
The first equality is Lemma 8.6 together with Remark 8.4, since $\mathcal X_f$ is $\beta$-invariant for every $\beta \geq 1$. For the second, put $s = t^{1/m}$ and $\beta = (1, \ldots, 1)$: the chart identifies the truncated arcs at $0$ with $(s\C[s]/(s^{m+1}))^d$, and for a $\C(\!(s)\!)$-rational point the condition $\rv(f(x)) = \rv(t)$ reads $f(x) = s^m \bmod s^{m+1}$; so $\mathcal X_f[m;\beta]$ is $\mathcal X_m$ in the variable $s$, with $\mu_m$ acting by $s \efun \zeta s$, and the normalizing factor is $[\A]^{-m\Sigma\beta} = [\A]^{-md}$.
\end{proof}

\setcounter{theorem}{6}
\begin{remark}[Hrushovski--Kazhdan, Proposition 10.10(2), fails]
\label{rem-HK-fail}
\lean{DiamondIntegral.Section8.Remark87.remark_8_7, DiamondIntegral.Section8.Remark87.germ, DiamondIntegral.Section8.Hm.map_lattice}
\leanok
\uses{jcbG-op, red-D-iso, z-coeff-change}
The tensor decomposition of $\ggk\,\mathrm{vol}\RV[*]$ claimed in Hrushovski--Kazhdan, Proposition 10.10(2), and the surjectivity of the morphisms (8.1.4), (8.1.6) and (8.1.7) of Hrushovski--Loeser do not hold. Let $X = \{u \in \RV : 0 < \vrv(u) < 1\}$ and $f(u) = u^2$. Then $(X, f)$ is an object of $\mathrm{vol}\RV[1]$ whose class is not the class of a finite disjoint union of products $W \times \Delta^\sharp$ with $W$ in $\mathrm{vol}\RES$ and $\Delta \sub \Gamma^l$ definable. In $\mRV[1]$, by contrast, the identity of $X$ is an isomorphism from $(X, f, 0)$ onto $(X, \mathrm{id}, \vrv)$, which lies in the image of $\mu\Psi$.

\emph{Formalization.} Lean proves it for any target covered by a part on which the weight takes finitely many values and by products $W \times \Delta^\sharp$ with $W$ finite and $\Delta$ definable, contained in it, which is more general than the paper's disjoint unions; the proof replaces rational points by their values. The closing comparison with $\mRV[1]$ is not stated (FINAL-CHECK.md).
\end{remark}

\begin{proof}
\uses{cut-to-hensel-substru}
\leanok
Choose $m$ so that the coordinates of the points of the finite sets $W$ have values in $(1/m)\Z$, $2/m$ is not a weight of a product of the first kind, and $2/m$ lies below the point after which each $\Delta$ is constant (a definable subset of $\Gamma$ is a finite union of points and intervals with rational endpoints). A point of $X$ of value $1/m$ has weight $2/m$, so its image lies in a product $W \times \Delta^\sharp$ of the second kind with $2/m \in \Delta$, hence $1/m \in \Delta$, and the point $(p, \rv(t^{1/m}))$ of $Y$ has weight $1/m$ and all coordinates with values in $(1/m)\Z$. Definable maps preserve this lattice of values (functional closedness of $\C((t^{1/m}))$), so its preimage in $X$ has value in $(1/m)\Z$ and weight in $(2/m)\Z$, which does not contain $1/m$.
\end{proof}

\setcounter{theorem}{7}
\begin{definition}[Rational series and their limit at infinity]
\label{dag-loc}
\lean{DiamondIntegral.Section8.RationalSeries.daggerDenominators, DiamondIntegral.Section8.RationalSeries.IsDagger, DiamondIntegral.Section8.RationalSeries.HasLimit}
\leanok
\uses{bipo-gdv}
$\gmv[[\A]^{-1}][T]_{\dagger}$ is the localization of $\gmv[[\A]^{-1}][T]$ by the multiplicative family generated by the elements $1 - [\A]^a T^b$, $a \in \Z$, $b \in \Z^+$; it may be regarded as a subring of $\gmv[[\A]^{-1}][\![T]\!]$. The text after the notation makes the same construction $R[T]_{\dagger} \sub R[\![T]\!]$ for a commutative ring $R$ and a unit $A \in R$. For $h = P/Q \in R[T]_{\dagger}$ put $\deg h = \deg P - \deg Q$; if $\deg h \leq 0$, then $\lim_{T \to \infty} h$ is the constant term of the expansion of $h$ in powers of $T^{-1}$, that is, $0$ if $\deg h < 0$ and the quotient of the leading coefficients of $P$ and $Q$ if $\deg h = 0$. The elements of degree at most $0$ form a subgroup on which this limit is additive.

\emph{Formalization.} Lean encodes membership in $R[T]_{\dagger}$ as a predicate on power series (\texttt{IsDagger}) and bundles membership, degree at most $0$ and the value of the limit in \texttt{HasLimit}.
\end{definition}

\setcounter{theorem}{8}
\begin{lemma}[Limit of lattice-point zeta series]
\label{lat-zeta}
\lean{DiamondIntegral.Section8.PuiseuxValues.lemma_8_9, DiamondIntegral.Section8.Prism.hasLimit_latticeSeries}
\leanok
\uses{dag-loc}
Let $R$ be a commutative ring, $A \in R$ a unit, and $R[T]_{\dagger}$ and $\lim_{T \to \infty}$ as in Notation 8.8; identify $\Gamma$ with $\Q$. Let $I \sub \Gamma^n$ be a bounded definable set and $\sigma : I \fun \Gamma$ a function that, on each piece of a finite definable partition of $I$, is of the form $\gamma \efun \sum_i a_i \gamma_i + b$ with $a_i$, $b$ in $\Z$. For $m \in \Z^+$ let $I_m = I \cap (m^{-1}\Z)^n$, and put \[ Z(\sigma)(T) = \sum_{m \in \Z^+} \sum_{\gamma \in I_m} A^{-m\sigma(\gamma)}\, T^m \in R[\![T]\!]. \] Then $Z(\sigma)(T) \in R[T]_{\dagger}$, $\deg Z(\sigma)(T) \leq 0$, and $\lim_{T \to \infty} Z(\sigma)(T) = -\chi(I)$.
\end{lemma}

\begin{proof}
\uses{dag-loc}
\leanok
A band between two affine functions over an open rational simplex is a finite disjoint union of open simplices; by induction on $n$, a bounded Boolean combination of affine sign conditions is a finite disjoint union of open simplices on which the given functions have constant sign. For an open simplex $D$ of dimension $k$ on which $\sigma$ is affine, the lattice points of the cone over $D$ are uniquely of the form $p + \sum_j n_j u_j$ with $p$ in a finite set, so $Z(\sigma \rest D)$ is a quotient by $\prod_{j=0}^{k}(1 - A^{-L(u_j)}T^{m_j})$ of degree $0$ with limit $(-1)^{k+1}$. Summing over the simplices, with $\chi(D) = (-1)^{\dim D}$ and the additivity of the limit, gives $-\chi(I)$. This is a direct proof of Hrushovski--Loeser, Lemma 8.5.2.
\end{proof}

\setcounter{theorem}{9}
\begin{remark}[Evaluation by monomial decompositions]
\label{s8-theta-ediamond}
\lean{DiamondIntegral.Section8.Hm.Ediamond, DiamondIntegral.Section8.Hm.EdiamondQ, DiamondIntegral.Section8.Hm.Ediamond_sharp, DiamondIntegral.Section8.Hm.Ediamond_res}
\leanok
\uses{prop-eu-retr-k-db, db-RV, coh-decom, bipo-gdv}
Let $\mathbb E^\diamond : \ggk \mRV^{\db}[*] \fun {!\ggk}\RES$ be the composite of the homomorphism $\mgE^{\db}$ of Proposition 4.15 with the forgetful map ${!\ggk}\RES[*] \fun {!\ggk}\RES$, as in diagram (4.8). By Lemma 5.1 and Notation 5.2, the homomorphism \[ \Theta \circ \mathbb E^\diamond : \ggk \mRV^{\db}[*] \fun \gmv \] can be constructed directly by the expression (5.1) with respect to monomial decompositions: if $h_{M_i}(U_i) = V_i \times I_i^\sharp$ is a monomial decomposition of $\mathbf U$ with $\ell_i$ free $\RV$-coordinates, then $(\Theta \circ \mathbb E^\diamond)([\mathbf U]) = \sum_i \chi_b(I_i)\, \Theta([V_i])[\G_m]^{\ell_i}$.

\emph{Formalization.} Lean defines $\mathbb E^\diamond$ as the evaluation on $\ggk \mRV^{\db}[*]/J_{\mathrm{var}}$ followed by the forgetful map (\texttt{Ediamond}, \texttt{EdiamondQ}) and computes it on product generators $[g^\sharp \times V]$ (\texttt{Ediamond\_sharp}, \texttt{Ediamond\_res}). The global monomial-decomposition expression (5.1) is not stated (section 8/LEAN-INTERFACE.md; FINAL-CHECK.md, section 3, row for Section 5).
\end{remark}

\setcounter{theorem}{10}
\begin{theorem}[Motivic Milnor fiber via the bounded integral]
\label{direct-mil}
\lean{DiamondIntegral.Section8.Closed.theorem_8_11, DiamondIntegral.Section8.Milnor.theorem_8_11_milnor, DiamondIntegral.Section8.Milnor.theorem_8_11_smooth, DiamondIntegral.Section8.Hm.theorem_8_11_rv_ring}
\leanok
\uses{loc-rings, muhat-section, dia-to-db-dag, dag-loc, s8-theta-ediamond, defn-binv, defn-diamond, main-prop, bipo-gdv, hyp-etale-chart}
Let $X$ be a smooth connected complex algebraic variety of dimension $d$, $f : X \fun \A^1$ a nonconstant regular function and $0 \in X(\C)$ with $f(0) = 0$. Let $\mathcal X_f = \{x \in X(\mathcal M) : \rv(f(x)) = \rv(t)\}$, identified through an \'etale chart at $0$ with a subset of $\mathcal M^d$; it is $\beta$-invariant for every $\beta \geq 1$, hence an object of $\mVF^\diamond[d]$ with constant volume form $0$. Let $Z_f(T) = \sum_{m \in \Z^+} (\eta \circ h_m \circ \int^\diamond)([\mathcal X_f])\, T^m \in \gmv[[\A]^{-1}][\![T]\!]$, well defined by Remark 8.4. Then $Z_f(T)$ lies in $\gmv[[\A]^{-1}][T]_{\dagger}$ and has degree at most $0$, and the motivic Milnor fiber $\mathcal S_f = -\lim_{T \to \infty} Z_f(T)$ satisfies \[ \mathcal S_f = \Big(\Theta \circ \mathbb E^\diamond \circ \int^\diamond\Big)([\mathcal X_f]). \]

\emph{Formalization.} Lean proves it for $X = \A^d$ with any polynomial $f$ (\texttt{theorem\_8\_11\_milnor}, checked in FINAL-CHECK.md, section 3) and, under the \'etale chart hypothesis taken by the author's instruction of 2026-11-21 (section 8/CITED-INPUTS.md), for smooth $X$ given near $0$ by polynomial equations in $\A^N$ (\texttt{theorem\_8\_11\_smooth}); the RV-side form \texttt{theorem\_8\_11\_rv\_ring} holds for every element of $\ggk \mRV^{\db}[*]$, its zeta series using the integer-exponent part of $h_m$. In these theorems the encoding hypotheses (a valuation compatible with the language and nontrivial on the base) are variables; \texttt{Closed.theorem\_8\_11} supplies them, with the Puiseux ambient of Section 5, so it holds for every complex polynomial with no hypothesis.
\end{theorem}

\begin{proof}
\uses{dia-to-db-dag, coh-decom, comp-form, GLZ-char, bipo-gdv, s8-theta-ediamond, muhat-section, lat-zeta}
\leanok
Represent $\int^\diamond[\mathcal X_f]$ by $[\mathbf U] = [(U, g, l)] \in \ggk^\natural \mRV^{\db}[*]$ with $l \in \Z$ (Remark 8.4), so that $Z_f(T) = \sum_m (\eta \circ h_m)([\mathbf U])\, T^m$. Both sides of the desired equality are additive, so by a monomial decomposition (Lemma 5.1) one may assume $h_M(U) = V \times I^\sharp$ with $I$ bounded and $l_g$ contracting to a function $\sigma$, piecewise $\Z$-affine by Remark 3.15; with $v = \Theta([V])[\G_m]^\ell$, (5.1) gives $(\Theta \circ \mathbb E^\diamond)([\mathbf U]) = \chi(I)v$ and Lemma 8.2 gives $(\eta \circ h_m)([\mathbf U]) = v \sum_{\gamma \in D_m} [\A]^{-m\sigma(\gamma)}$. Rescaling the lattice by the least $N \in \Z^+$ with $N\beta$ integral, where $\vrv(V) = \{\beta\}$, gives $Z_f(T) = v\, Z(\sigma')(T^N)$, and Lemma 8.9 yields membership, degree at most $0$ and limit $-\chi(I)v$.
\end{proof}

\setcounter{theorem}{12}
\begin{remark}[Theorem 8.11 fails for bounded objects]
\label{sec-hm-bounded}
\uses{direct-mil}
The homomorphisms $\eta \circ h_m$ are defined on doubly bounded objects, and Theorem 8.11 does not extend to bounded objects in the obvious way. For the bounded object $\bm U = (\RV^{\circ\circ}, \mathrm{id}, 0)$ the natural value of each $h_m$ is $1$, so its zeta function is $T/(1-T)$ and $-\lim_{T \to \infty}$ of it is $1$, whereas $(\Theta \circ \mathbb{E}_b)([\bm U]) = 0$. For $(\RV^{\circ\circ}, \mathrm{id}, -\vrv)$ the defining expression is not summable.

\emph{Formalization.} Not formalized. The general claims of this remark were deleted from the paper by the author in the final review; leaving the remaining example unformalized was a choice of the formalization session (FINAL-CHECK.md).
\end{remark}

\setcounter{theorem}{13}
\begin{remark}[Milnor fiber without inverting the affine line]
\label{rem-comp-milnor-fib}
\lean{DiamondIntegral.Section8.VolComparison.remark_8_14, DiamondIntegral.Section8.VolComparison.milnor_fiber_vol, DiamondIntegral.Section8.Milnor.milnor_fiber_vol_smooth, DiamondIntegral.Section8.Closed.remark_8_14}
\leanok
\uses{direct-mil, mov-vol, defn-diamond, main-prop, s8-theta-ediamond, hyp-etale-chart}
Let $X$, $f$, $0$, $\mathcal X_f$, $Z_f(T)$ and $\mathcal S_f$ be as in Theorem 8.11. Since the coefficients of $Z_f(T)$ involve $[\A]^{-1}$ and $\gmv \fun \gmv[[\A]^{-1}]$ is not injective ($[\A]$ is a zero divisor in $\ggk \Var_{\C}$ by Borisov), $\mathcal S_f$ cannot be taken in $\gmv$ through $Z_f(T)$; this is why $\int^\diamond$ and the proper invariant $\mathcal X_f$ are used instead of $\mathcal X_t = \{x \in X(\mathcal M) : f(x) = t\}$, which is not proper invariant. By Theorem 8.11, however, $\mathcal S_f = (\Theta \circ \mathbb E^\diamond \circ \int^\diamond)([\mathcal X_f])$ with no need to invert $[\A]$, and \[ \mathcal S_f = \Vol([\mathcal X_f]) \] (Notation 5.3), because the left half of diagram (1.2) commutes (Corollary 7.27).

\emph{Formalization.} Lean proves $\Theta \circ \mathbb E^\diamond \circ \int^\diamond = \Vol$ on every object of $\mVF^\diamond[k]$ (\texttt{remark\_8\_14}), hence $\mathcal S_f = \Vol([\mathcal X_f])$ for $X = \A^d$ (\texttt{remark\_8\_14\_milnor}, \texttt{milnor\_fiber\_vol}) and, under the \'etale chart hypothesis of 2026-11-21 (section 8/CITED-INPUTS.md), for smooth $X$ (\texttt{milnor\_fiber\_vol\_smooth}). The commentary of the first paragraph (Borisov's zero divisor, $\mathcal X_t$ not proper invariant) has no Lean statement. \texttt{Closed.remark\_8\_14} states it for $X = \A^d$ with every hypothesis supplied.
\end{remark}

\begin{proof}
\uses{direct-mil, db-to-novol}
\leanok
Corollary 7.27 identifies $\mathbb E^\diamond \circ \int^\diamond$ on $[\mathcal X_f]$ with $\mathbb E_b \circ \int$ of the underlying definable set, so $(\Theta \circ \mathbb E^\diamond \circ \int^\diamond)([\mathcal X_f]) = \Vol([\mathcal X_f])$; combine with Theorem 8.11.
\end{proof}

\setcounter{theorem}{14}
\begin{remark}[Milnor fiber through the fibre over t]
\label{rem-comp-milnor-fib-bis}
\lean{DiamondIntegral.Section8.Remark815.remark_8_15_integral, DiamondIntegral.Section8.Remark815.remark_8_15, DiamondIntegral.Section8.Remark815.remark_8_15_milnor_smooth, DiamondIntegral.Section8.Closed.remark_8_15}
\leanok
\uses{direct-mil, mov-vol, rem-comp-milnor-fib, hyp-etale-chart}
Let $X$, $f$, $0$, $\mathcal X_f$ and $\mathcal S_f$ be as in Theorem 8.11, and for $t' \in \rv(t)^\sharp$ let $\mathcal X_{t'} = \{x \in X(\mathcal M) : f(x) = t'\}$, so that $\mathcal X_f = \bigcup_{t' \in \rv(t)^\sharp} \mathcal X_{t'}$. All $t' \in \rv(t)^\sharp$ have the same type over $\C \cup \RV$. Consequently there are an object $\mathbf U$ of $\RV[*]$ and a bijection $\mathcal X_f \fun \rv(t)^\sharp \times \mathbb L \mathbf U$ definable over $\mathbb S$, and \[ \int [\mathcal X_f] = [1] \int [\mathcal X_t], \qquad \mathcal S_f = \Vol([\mathcal X_f]) = \Vol([\mathcal X_t]). \]

\emph{Formalization.} Lean proves the integral and volume identities for $X(\mathcal M)$ cut out by any polynomial equations (\texttt{remark\_8\_15\_integral}, \texttt{remark\_8\_15}) and links them to $\mathcal S_f$ for $X = \A^d$ (\texttt{remark\_8\_15\_milnor}) and, under the \'etale chart hypothesis of 2026-11-21 (section 8/CITED-INPUTS.md), for smooth $X$ (\texttt{remark\_8\_15\_milnor\_smooth}). The closing paragraph on earlier proofs in the literature is commentary. \texttt{Closed.remark\_8\_15} states the identity with $\mathcal S_f$ for $X = \A^d$ with every hypothesis supplied.
\end{remark}

\begin{proof}
\uses{main-prop-HK, mov-vol, rem-comp-milnor-fib}
\leanok
By quantifier elimination, $\rv(P(t')) = \rv(a_j)\rv(t)^j$ for the lowest term $a_j x^j$ of a nonzero $P \in \C[x]$, which gives the type statement. Over the substructure generated by $\C$ and $t$, $\mathbb L$ is surjective on isomorphism classes (see the discussion before Theorem 3.22), so there is a definable bijection $F_t : \mathcal X_t \fun \mathbb L \mathbf U$; replacing $t$ by $t'$ in its defining formulas gives $F_{t'}$, and $x \efun (f(x), F_{f(x)}(x))$ is the required bijection. Hence $\int[\mathcal X_f] = [1]\int[\mathcal X_t]$, and $\mathbb E_b([1]) = 1$ with Remark 8.14 gives the equalities of volumes.
\end{proof}
